% III : formuler les idées essentielles d’un texte à résumer — Français Tle Tech — Cours signé OnBuch+
% Programme officiel MINESEC. Compiler avec : tectonic N15-iii-formuler-idees-essentielles-texte-resumer.tex
\newcommand{\DOCMATIERE}{Français}
\newcommand{\DOCNIVEAU}{Tle Tech}
\newcommand{\DOCMODULE}{Français – classe de Terminale technique}
\newcommand{\DOCLECON}{N15}
\newcommand{\DOCTITRE}{III : formuler les idées essentielles d’un texte à résumer}
% OnBuch+ — préambule « cours signés » ENRICHI (compilé avec tectonic / XeLaTeX)
% Système visuel « Pop » d'OnBuch+ : fond crème (jamais blanc), bordures encre
% épaisses, ombres « dures » décalées, titres Archivo Black en capitales.
% Version MATHÉMATIQUES : graphes (pgfplots/tikz), boîtes pédagogiques
% (définition, propriété, méthode, attention, le savais-tu, exercice résolu,
% exercice, TP, bilan), page de garde et signature OnBuch+ ancrée en bas.
%
% Macros attendues dans main.tex AVANT \input{preamble} :
%   \DOCMATIERE  \DOCNIVEAU  \DOCMODULE  \DOCLECON (numéro)  \DOCTITRE
\documentclass[11pt]{article}

\usepackage{fontspec}
\usepackage[french]{babel}
\usepackage[a4paper,margin=2cm,top=3.4cm,bottom=2.8cm,headheight=1.4cm,footskip=1.3cm]{geometry}
\usepackage{amsmath,amssymb,amsfonts}
\usepackage{mathtools} % \xmapsto, \coloneqq... souvent utilisés par les modèles
\usepackage{cancel} % simplifications barrées (bilans de forces)
% \abs{z} (module, valeur absolue) et \norm{u} (norme), utilisés par les modèles
\providecommand{\abs}{}\let\abs\relax\DeclarePairedDelimiter{\abs}{\lvert}{\rvert}
\providecommand{\norm}{}\let\norm\relax\DeclarePairedDelimiter{\norm}{\lVert}{\rVert}
\usepackage[table]{xcolor} % [table] : fournit aussi \rowcolors (alternance de lignes)
\usepackage{pagecolor}
\usepackage{tikz}
\usetikzlibrary{arrows.meta,positioning,calc,shapes.geometric,shapes.misc,shapes.arrows,decorations.pathmorphing,decorations.pathreplacing,decorations.markings,patterns,shadows,mindmap,trees,fit,backgrounds,matrix,intersections,angles,quotes}
\usepackage{pgfplots}
\usepackage[version=4]{mhchem} % équations nucléaires \ce{^{235}_{92}U}
\usepackage[european,siunitx]{circuitikz} % schémas électriques
\pgfplotsset{compat=1.18}
\usepgfplotslibrary{fillbetween,groupplots} % aires entre courbes (calcul intégral)
\usepackage{enumitem}
\usepackage{fancyhdr}
% [most] charge toutes les bibliothèques tcolorbox (listings, minted, raster,
% documentation...) et gonfle inutilement la mémoire TeX sur un document long
% (~300 boîtes) : on ne charge que ce qui est réellement utilisé.
\usepackage[skins,breakable]{tcolorbox}
\usepackage{graphicx}
\usepackage{colortbl}
\usepackage{booktabs}
\usepackage{tabularx}
\usepackage{multirow}
\usepackage{diagbox} % case d'en-tête diagonale des tableaux à double entrée
% Colonnes extensibles centrée / gauche / droite (C, L, R), souvent utilisées par les modèles
\newcolumntype{C}{>{\centering\arraybackslash}X}
\newcolumntype{L}{>{\raggedright\arraybackslash}X}
\newcolumntype{R}{>{\raggedleft\arraybackslash}X}
\usepackage{array}
\usepackage{multicol}
\usepackage{siunitx}
% parse-numbers=false : accepte aussi \SI{2,5 \times 10^{-3}}{...}
\sisetup{locale=FR,output-decimal-marker={,},group-minimum-digits=4,parse-numbers=false}
\DeclareSIUnit\ohm{\ensuremath{\symup{\Omega}}} % U+2126 absent de la police
\DeclareSIUnit\division{div} % réglages d'oscilloscope (ms/div, V/div)
\DeclareSIUnit\tr{tr} % tours (vitesse de rotation, ex. \SI{1500}{\tr\per\minute})
\DeclareSIUnit\molar{M} % concentration molaire (ex. \SI{3}{\molar}, \SI{1}{\micro\molar})
\DeclareSIUnit\an{an} % année (le modèle confond parfois avec \year, un compteur TeX) — ex. \SI{5}{\kilo\meter\per\an}
\DeclareSIUnit\FCFA{FCFA} % franc CFA (ex. \SI{500}{\FCFA\per\kilo\gram})
\DeclareSIUnit\degreeBrix{\textdegree Brix} % teneur en sucre d'un sirop/jus (réfractométrie)
\DeclareSIUnit\month{mois} % le modèle utilise parfois \month, un compteur TeX, comme unité de durée
\DeclareSIUnit\microliter{\micro\liter} % le modèle écrit parfois \microliter en un seul token
\DeclareSIUnit\million{millions} % dénombrement (ex. \SI{15}{\million\per\mL} spermatozoïdes)
\usepackage{qrcode}
% \degree : commande fréquemment utilisée par les modèles pour un angle ou une
% température en dehors de \SI{}{} (non fournie par les paquets chargés).
\providecommand{\degree}{^{\circ}}
% \qed (fin de démonstration), utilisé par les modèles sans amsthm
\providecommand{\qed}{\hfill$\square$}
% \barre : double barre des valeurs interdites dans un tableau de variations (tkz-tab)
\providecommand{\barre}{\ensuremath{\|}}
% environnement proof (amsthm non chargé) utilisé par les modèles
\ifdefined\proof\else\newenvironment{proof}[1][Démonstration]{\par\noindent\textit{#1.}\ }{\qed\par}\fi
\usepackage{lastpage}
\usepackage{hyperref}
\hypersetup{hidelinks}

% ============================================================
% Palette « Pop » OnBuch+ — mêmes hex que la classe P (plus_ui.dart)
% ============================================================
\definecolor{poporange}{HTML}{F59321}
\definecolor{popdark}{HTML}{E07A0C}
\definecolor{popcream}{HTML}{FFF6E8}
\definecolor{popink}{HTML}{1C1714}
\definecolor{poppurple}{HTML}{7A5AE0}
\definecolor{popgreen}{HTML}{2E9E6A}
\definecolor{poppink}{HTML}{E0457B}
\definecolor{popblue}{HTML}{2F7DE1}
\definecolor{popgold}{HTML}{C98A12}
\definecolor{popmuted}{HTML}{8A7F76}
\definecolor{popline}{HTML}{EADFD2}
\definecolor{popgray}{HTML}{8A7F76}
\colorlet{popgrey}{popgray}
% Alias tolérants (le modèle nomme parfois les couleurs différemment)
\colorlet{popwhite}{white}
\colorlet{popblack}{popink}
\colorlet{popred}{poppink}
\colorlet{popyellow}{popgold}
% Teintes claires (fonds de tableaux/encadrés) — le modèle nomme parfois
% les couleurs avec un suffixe L (ex. popblueL) sans les avoir définies.
\colorlet{poporangeL}{poporange!20}
\colorlet{popdarkL}{popdark!20}
\colorlet{poppurpleL}{poppurple!20}
\colorlet{popgreenL}{popgreen!20}
\colorlet{poppinkL}{poppink!20}
\colorlet{popblueL}{popblue!20}
\colorlet{popgoldL}{popgold!20}
\colorlet{popredL}{popred!20}
\colorlet{popgrayL}{popgray!20}
% Teintes claires pour fonds de tableaux / zones
\colorlet{popblueL}{popblue!10!white}
\colorlet{poporangeL}{poporange!14!white}
\colorlet{popgreenL}{popgreen!12!white}
\colorlet{poppurpleL}{poppurple!12!white}
\colorlet{poppinkL}{poppink!12!white}
\colorlet{popgoldL}{popgold!14!white}

\pagecolor{popcream}

% ============================================================
% Typographie
% ============================================================
\defaultfontfeatures{Ligatures=TeX}
\setmainfont{PlusJakartaSans-Medium.ttf}[Path=../../../fonts/,BoldFont=PlusJakartaSans-Bold.ttf]
% Maths en sans-serif (Fira Math) avec chiffres et lettres droites de Plus
% Jakarta, pour que les formules se fondent dans le texte.
\usepackage[math-style=ISO,bold-style=ISO]{unicode-math}
\setmathfont{FiraMath-Regular.otf}[Path=../../../fonts/]
\setmathfont{PlusJakartaSans-Medium.ttf}[Path=../../../fonts/,range={up/num,up/{Latin,latin},\mathup}]
\usepackage{icomma} % virgule décimale sans espace en mode math
% Caractères mathématiques Unicode absents des polices de texte (Plus Jakarta,
% Archivo Black) : rendus par leur équivalent mathématique, en texte comme en
% formule — sinon ils s'affichent en carré vide (ex. « Arithmétique dans ℤ »).
\usepackage{newunicodechar}
\newunicodechar{ℝ}{\ensuremath{\mathbb{R}}}
\newunicodechar{ℤ}{\ensuremath{\mathbb{Z}}}
\newunicodechar{ℂ}{\ensuremath{\mathbb{C}}}
\newunicodechar{ℕ}{\ensuremath{\mathbb{N}}}
\newunicodechar{ℚ}{\ensuremath{\mathbb{Q}}}
\newunicodechar{𝔻}{\ensuremath{\mathbb{D}}}
\newunicodechar{α}{\ensuremath{\alpha}}
\newunicodechar{β}{\ensuremath{\beta}}
\newunicodechar{γ}{\ensuremath{\gamma}}
\newunicodechar{δ}{\ensuremath{\delta}}
\newunicodechar{ε}{\ensuremath{\varepsilon}}
\newunicodechar{θ}{\ensuremath{\theta}}
\newunicodechar{λ}{\ensuremath{\lambda}}
\newunicodechar{μ}{\ensuremath{\mu}}
\newunicodechar{σ}{\ensuremath{\sigma}}
\newunicodechar{τ}{\ensuremath{\tau}}
\newunicodechar{φ}{\ensuremath{\varphi}}
\newunicodechar{ψ}{\ensuremath{\psi}}
\newunicodechar{ω}{\ensuremath{\omega}}
\newunicodechar{Σ}{\ensuremath{\Sigma}}
% majuscules produites par \MakeUppercase dans les titres (α-aminés -> Α)
\newunicodechar{Α}{\ensuremath{\alpha}}
\newunicodechar{Β}{\ensuremath{\beta}}
\newunicodechar{Γ}{\ensuremath{\Gamma}}
\newunicodechar{Δ}{\ensuremath{\Delta}}
\newunicodechar{Ε}{\ensuremath{\varepsilon}}
\newunicodechar{Θ}{\ensuremath{\Theta}}
\newunicodechar{Λ}{\ensuremath{\Lambda}}
\newunicodechar{Μ}{\ensuremath{\mu}}
\newunicodechar{Τ}{\ensuremath{\tau}}
\newunicodechar{Φ}{\ensuremath{\Phi}}
\newunicodechar{Ψ}{\ensuremath{\Psi}}
\newunicodechar{Ω}{\ensuremath{\Omega}}
\newunicodechar{↦}{\ensuremath{\mapsto}}
\newunicodechar{∈}{\ensuremath{\in}}
\newunicodechar{∉}{\ensuremath{\notin}}
\newunicodechar{∧}{\ensuremath{\wedge}}
\newunicodechar{⇔}{\ensuremath{\Leftrightarrow}}
\newunicodechar{⟺}{\ensuremath{\Longleftrightarrow}}
\newunicodechar{⇒}{\ensuremath{\Rightarrow}}
\newunicodechar{⟹}{\ensuremath{\Longrightarrow}}
\newunicodechar{∘}{\ensuremath{\circ}}
\newunicodechar{∪}{\ensuremath{\cup}}
\newunicodechar{∩}{\ensuremath{\cap}}
\newunicodechar{≡}{\ensuremath{\equiv}}
\newunicodechar{⊂}{\ensuremath{\subset}}
\newunicodechar{⊥}{\ensuremath{\perp}}
\newunicodechar{∃}{\ensuremath{\exists}}
\newunicodechar{∀}{\ensuremath{\forall}}
\newunicodechar{∛}{\ensuremath{\sqrt[3]{\,}}}
\newunicodechar{∜}{\ensuremath{\sqrt[4]{\,}}}
\newunicodechar{⃗}{\ensuremath{{}^{\to}}}
\newunicodechar{ⁿ}{\ifmmode^{n}\else\textsuperscript{\ensuremath{n}}\fi}
\newunicodechar{ˣ}{\ifmmode^{x}\else\textsuperscript{\ensuremath{x}}\fi}
\newunicodechar{ᵇ}{\ifmmode^{b}\else\textsuperscript{\ensuremath{b}}\fi}
\newunicodechar{ʳ}{\ifmmode^{r}\else\textsuperscript{\ensuremath{r}}\fi}
\newunicodechar{ᵘ}{\ifmmode^{u}\else\textsuperscript{\ensuremath{u}}\fi}
\newunicodechar{ʸ}{\ifmmode^{y}\else\textsuperscript{\ensuremath{y}}\fi}
\newunicodechar{ᵃ}{\ifmmode^{a}\else\textsuperscript{\ensuremath{a}}\fi}
\newunicodechar{ᵉ}{\ifmmode^{e}\else\textsuperscript{\ensuremath{e}}\fi}
\newunicodechar{ᵏ}{\ifmmode^{k}\else\textsuperscript{\ensuremath{k}}\fi}
\newunicodechar{ᵖ}{\ifmmode^{p}\else\textsuperscript{\ensuremath{p}}\fi}
\newunicodechar{ᴺ}{\ifmmode^{N}\else\textsuperscript{\ensuremath{N}}\fi}
\newunicodechar{ᶜ}{\ifmmode^{c}\else\textsuperscript{\ensuremath{c}}\fi}
\newunicodechar{ʰ}{\ifmmode^{h}\else\textsuperscript{\ensuremath{h}}\fi}
\newunicodechar{ᵠ}{\ifmmode^{q}\else\textsuperscript{\ensuremath{q}}\fi}
\newunicodechar{ₙ}{\ifmmode_{n}\else\textsubscript{\ensuremath{n}}\fi}
\newunicodechar{ₐ}{\ifmmode_{a}\else\textsubscript{\ensuremath{a}}\fi}
\newunicodechar{ᵢ}{\ifmmode_{i}\else\textsubscript{\ensuremath{i}}\fi}
\newunicodechar{ᵣ}{\ifmmode_{r}\else\textsubscript{\ensuremath{r}}\fi}
\newunicodechar{ₖ}{\ifmmode_{k}\else\textsubscript{\ensuremath{k}}\fi}
\newunicodechar{ᵦ}{\ifmmode_{\beta}\else\textsubscript{\ensuremath{\beta}}\fi}
\newunicodechar{ₓ}{\ifmmode_{x}\else\textsubscript{\ensuremath{x}}\fi}
\newfontfamily\popdisplay{ArchivoBlack-Regular.ttf}[Path=../../../fonts/]
\newfontfamily\poplogo{SpaceGrotesk-Bold.ttf}[Path=../../../fonts/]
\newfontfamily\popsemi{PlusJakartaSans-SemiBold.ttf}[Path=../../../fonts/]
\newfontfamily\popxbold{PlusJakartaSans-ExtraBold.ttf}[Path=../../../fonts/]
\newfontfamily\popxboldls{PlusJakartaSans-ExtraBold.ttf}[Path=../../../fonts/,LetterSpace=3.0]

\setlist[enumerate]{leftmargin=1.7em,itemsep=5pt,topsep=4pt}
\setlist[itemize]{leftmargin=1.7em,itemsep=4pt,topsep=4pt,label={\color{poporange}\textbullet}}
\setlength{\parindent}{0pt}
\setlength{\parskip}{5pt}
\renewcommand{\arraystretch}{1.25}

% pgfplots : style Pop par défaut
\pgfplotsset{
  every axis/.append style={axis line style={popink,line width=1pt},tick style={popink},
    grid style={popline,line width=0.5pt},label style={font=\small},tick label style={font=\footnotesize}},
  popcurve/.style={poporange,line width=1.8pt,no markers,smooth},
}
% Même style disponible dans un \draw TikZ simple (hors pgfplots)
\tikzset{popcurve/.style={poporange,line width=1.8pt}}
\tikzset{popgrid/.style={popline,very thin}}
\colorlet{popcurve}{poporange} % le nom du style est parfois utilisé comme couleur

% ============================================================
% Pill / squiggle
% ============================================================
\newcommand{\poppill}[3]{%
  \tikz[baseline=-0.55ex]\node[draw=popink,line width=1.2pt,rounded corners=2.6pt,%
    fill=#2,inner xsep=6.5pt,inner ysep=3pt]{%
    \popxboldls\fontsize{7}{7}\selectfont\color{#3}\MakeUppercase{#1}};%
}
\newcommand{\popsquiggle}[1]{%
  \tikz[baseline=-0.4ex]{
    \draw[#1,line width=2.6pt,line cap=round]
      plot [smooth] coordinates {(0,0.09) (0.34,0.01) (0.68,0.21) (1.02,0.01) (1.36,0.10)};
  }%
}
% Mot-clé surligné (marqueur orange sous le texte)
\newcommand{\cle}[1]{{\popxbold\color{popdark}#1}}

% ============================================================
% En-tête / pied
% ============================================================
\pagestyle{fancy}
\fancyhf{}
\renewcommand{\headrulewidth}{0pt}
\renewcommand{\footrulewidth}{0pt}
\lhead{\raisebox{-0.15ex}{\poplogo\fontsize{13}{13}\selectfont\color{popink}OnBuch\color{poporange}+}}
\rhead{\poppill{\DOCMATIERE\ \textbullet\ \DOCNIVEAU\ \textbullet\ Leçon \DOCLECON}{white}{popink}}
\lfoot{\popxbold\fontsize{7.6}{7.6}\selectfont\color{popmuted}\MakeUppercase{Cours signé OnBuch+}\ \textbullet\ {\popsemi\DOCTITRE}}
\rfoot{\tikz[baseline=-0.6ex]\node[rounded corners=8.5pt,fill=popink,minimum height=17pt,minimum width=17pt,inner xsep=5pt,inner sep=0]{\popxbold\fontsize{8}{8}\selectfont\color{popcream}\thepage\,/\,\pageref*{LastPage}};}

% ============================================================
% Titres
% ============================================================
\newcommand{\chaptitre}[2]{%
  \begin{tcolorbox}[enhanced,colback=poporange,colframe=popink,boxrule=2.6pt,arc=11pt,
      drop shadow=popink,left=15pt,right=15pt,top=12pt,bottom=11pt,before skip=3pt,after skip=18pt]
    \poppill{#2}{popink}{popcream}\par\vspace{9pt}
    {\raggedright\hyphenpenalty=10000\exhyphenpenalty=10000\popdisplay\fontsize{22}{24}\selectfont\color{popcream}\MakeUppercase{#1}\par}
  \end{tcolorbox}
}
% Section numérotée : pastille orange avec le numéro + titre Archivo
\newcounter{popsec}
\newcommand{\coursec}[1]{%
  \stepcounter{popsec}\setcounter{popsub}{0}%
  \par\vspace{16pt}\noindent
  \begin{minipage}{\linewidth}
  \tikz[baseline=-0.7ex]\node[draw=popink,line width=1.6pt,fill=poporange,rounded corners=4pt,minimum size=22pt,inner sep=0]{\popdisplay\fontsize{12}{12}\selectfont\color{popcream}\thepopsec};%
  \hspace{8pt}{\popdisplay\fontsize{15}{17}\selectfont\color{popink}\MakeUppercase{#1}}\par
  \vspace{4pt}\noindent{\color{poporange}\rule{36pt}{3.6pt}}
  \end{minipage}\par\nopagebreak\vspace{8pt}
}
\newcounter{popsub}
\newcommand{\courssub}[1]{%
  \stepcounter{popsub}%
  \par\vspace{9pt}\noindent{\popxbold\fontsize{11.5}{13}\selectfont\color{popink}{\color{poporange}\thepopsec.\thepopsub}\hspace{6pt}#1}\par\nopagebreak\vspace{4pt}
}

% ============================================================
% Boîtes pédagogiques — même recette : fond blanc, bordure épaisse colorée,
% ombre dure encre, badge plein. Titre optionnel [#1].
% ============================================================
\tcbset{popbase/.style={breakable,enhanced,colback=white,boxrule=2.3pt,arc=9pt,
  shadow={3pt}{-3pt}{0pt}{popink},left=14pt,right=14pt,top=11pt,bottom=11pt,before skip=11pt,after skip=15pt,
  fontupper=\popsemi\fontsize{10.6}{14.5}\selectfont\color{popink},
  fontlower=\popsemi\fontsize{10.6}{14.5}\selectfont\color{popink}}}
% Coche dessinée (le glyphe ✓ n'existe pas dans les polices du document)
\newcommand{\popcheck}[1]{\tikz[baseline=-0.5ex]\draw[#1,line width=1.6pt,line cap=round,line join=round](-0.13,0.0)--(-0.04,-0.09)--(0.13,0.1);}
\providecommand{\coche}{\popcheck{popgreen}}
\providecommand{\resultat}[1]{\par\smallskip\noindent\fcolorbox{popgreen}{popgreenL}{\parbox{\dimexpr\linewidth-2\fboxsep-2\fboxrule\relax}{\textbf{Résultat.} #1}}\par\smallskip}
\newcommand{\popboxhead}[3]{\poppill{#1}{#2}{white}\ifx\relax#3\relax\else\hspace{6pt}{\popxbold\fontsize{10.6}{12}\selectfont\color{popink}#3}\fi\par\vspace{7pt}}

\newtcolorbox{aretenir}[1][]{popbase,colframe=popgreen,before upper={\popboxhead{À retenir}{popgreen}{#1}}}
\newtcolorbox{exemplebox}[1][]{popbase,colframe=poporange,before upper={\popboxhead{Exemple}{poporange}{#1}}}
\newtcolorbox{definition}[1][]{popbase,colframe=popblue,before upper={\popboxhead{Définition}{popblue}{#1}}}
\newtcolorbox{propriete}[1][]{popbase,colframe=poppurple,before upper={\popboxhead{Propriété}{poppurple}{#1}}}
\newtcolorbox{methode}[1][]{popbase,colframe=popgold,before upper={\popboxhead{Méthode}{popgold}{#1}}}
\newtcolorbox{attention}[1][]{popbase,colframe=poppink,before upper={\popboxhead{Attention}{poppink}{#1}}}
\newtcolorbox{experience}[1][]{popbase,colframe=popblue,colback=popblueL,before upper={\popboxhead{Expérience}{popblue}{#1}}}
% « Le savais-tu ? » : carte violette pleine (contexte, culture, Cameroun)
\newtcolorbox{savaistu}[1][]{popbase,colback=poppurple,colframe=popink,
  fontupper=\popsemi\fontsize{10.4}{14.2}\selectfont\color{white},
  before upper={\poppill{Le savais-tu\,?}{white}{poppurple}\ifx\relax#1\relax\else\hspace{6pt}{\popxbold\color{white}#1}\fi\par\vspace{7pt}}}
% Exercice résolu : énoncé en haut, \tcblower puis la solution sur fond crème
\newcounter{popexo}\newcounter{popexr}
\newtcolorbox{exoresolu}[1][]{popbase,colframe=poporange,colbacklower=popcream,
  segmentation style={popink,line width=1.2pt,dashed},
  before upper={\stepcounter{popexr}\popboxhead{Exercice résolu \thepopexr}{poporange}{#1}},
  before lower={\poppill{Solution}{popink}{popcream}\par\vspace{6pt}}}
% Exercice (non résolu en place) — la correction est dans la partie Corrigés
\newtcolorbox{exercice}[1][]{popbase,colframe=popink,
  before upper={\stepcounter{popexo}\popboxhead{Exercice \thepopexo}{popink}{#1}}}
% Corrigé
\newtcolorbox{corrige}[1][]{popbase,colframe=popgreen,colback=popgreenL,
  before upper={\popboxhead{Corrigé}{popgreen}{#1}}}
% Objectifs / prérequis / bilan
\newtcolorbox{objectifs}{popbase,colframe=popink,colback=poporangeL,
  before upper={\popboxhead{Objectifs de la leçon}{popink}{}}}
\newtcolorbox{prerequis}{popbase,colframe=popink,colback=popblueL,
  before upper={\popboxhead{Prérequis}{popblue}{}}}
\newtcolorbox{bilan}{popbase,colframe=popink,colback=white,boxrule=2.8pt,
  before upper={\popboxhead{Fiche bilan}{poporange}{L'essentiel à connaître pour l'examen}}}
% Figure encadrée avec légende
\newtcolorbox{popfigure}[1][]{enhanced,breakable=false,colback=white,colframe=popline,boxrule=1.4pt,arc=8pt,
  left=10pt,right=10pt,top=10pt,bottom=8pt,before skip=10pt,after skip=12pt,halign=center,
  lower separated=false,halign lower=center,
  fontlower=\popsemi\fontsize{9}{11.5}\selectfont\color{popmuted},
  #1}
\newcommand{\legende}[1]{\par\vspace{6pt}{\popsemi\fontsize{9}{11.5}\selectfont\color{popmuted}#1}}

% ============================================================
% Signature OnBuch+
% ============================================================
\newcommand{\poplogoinline}[1]{{\poplogo\fontsize{#1}{#1}\selectfont\color{popink}OnBuch\color{poporange}+}}
\newcommand{\signatureonbuch}{%
  \begin{tcolorbox}[enhanced,colback=popink,colframe=popink,boxrule=2.2pt,arc=10pt,
      drop shadow=poporange,left=14pt,right=14pt,top=10pt,bottom=10pt,before skip=0pt,after skip=0pt]
    \tikz[baseline=-0.7ex]\node[circle,fill=popgreen,draw=popcream,line width=1.3pt,minimum size=22pt,inner sep=0]{\popcheck{white}};%
    \hspace{9pt}\begin{minipage}[c]{0.62\linewidth}
      {\popxbold\fontsize{12}{13}\selectfont\color{popcream}Signé\ \textbullet\ {\poplogo OnBuch\color{poporange}+}}\par\vspace{2pt}
      {\raggedright\popsemi\fontsize{8}{10.5}\selectfont\color{popline}\MakeUppercase{Équipe pédagogique OnBuch+}\\\MakeUppercase{Conforme au programme officiel MINESEC}\par}
    \end{minipage}\hfill
    {\poplogo\fontsize{22}{22}\selectfont\color{popcream}OnBuch\color{poporange}+}
  \end{tcolorbox}%
}

% ============================================================
% Page de garde
% #1 titre · #2 matière · #3 niveau · #4 module · #5 n° leçon · #6 durée ·
% #7 description / objectifs (texte ou itemize)
% ============================================================
\newcommand{\pagegarde}[7]{%
  \thispagestyle{empty}
  \newgeometry{a4paper,margin=1.7cm,top=1.6cm,bottom=1.4cm}%
  \begin{tikzpicture}[remember picture,overlay]
    \fill[poporange!18] ([xshift=-3.2cm,yshift=-3.6cm]current page.north east) circle (4.6cm);
    \fill[poppurple!14] ([xshift=2.4cm,yshift=7.6cm]current page.south west) circle (3.8cm);
  \end{tikzpicture}%
  {\poplogo\fontsize{30}{30}\selectfont\color{popink}OnBuch\color{poporange}+}\hfill
  \raisebox{0.6ex}{\poppill{Cours signé \textbullet\ Premium}{popink}{popcream}}\par
  \vspace{1pt}\popsquiggle{poppurple}\par
  \vspace{8pt}
  \begin{tcolorbox}[enhanced,colback=poporange,colframe=popink,boxrule=2.8pt,arc=13pt,
      drop shadow=popink,left=18pt,right=18pt,top=12pt,bottom=13pt,after skip=10pt]
    \poppill{#2 \textbullet\ #3}{popink}{popcream}\hspace{5pt}\poppill{#4}{white}{popink}\par\vspace{8pt}
    {\popdisplay\fontsize{14}{14}\selectfont\color{popink}LEÇON #5}\par\vspace{4pt}
    {\raggedright\hyphenpenalty=10000\exhyphenpenalty=10000\popdisplay\fontsize{25}{27}\selectfont\color{popcream}\MakeUppercase{#1}\par}
    \vspace{8pt}
    {\popxbold\fontsize{9.5}{9.5}\selectfont\color{popink}\MakeUppercase{Durée conseillée : #6}}
  \end{tcolorbox}
  \begin{tcolorbox}[enhanced,colback=white,colframe=popink,boxrule=2.2pt,arc=10pt,
      drop shadow=popink,left=16pt,right=16pt,top=10pt,bottom=10pt,after skip=8pt]
    \poppill{Dans cette leçon}{poppurple}{white}\par\vspace{5pt}
    \popsemi\fontsize{9.3}{12.6}\selectfont\color{popink}#7
  \end{tcolorbox}
  \noindent\begin{minipage}[t]{0.487\linewidth}
    \begin{tcolorbox}[enhanced,colback=popgreen,colframe=popink,boxrule=2.2pt,arc=10pt,
        drop shadow=popink,left=11pt,right=11pt,top=10pt,bottom=10pt,height=3.1cm,valign=center]
      \tcbox[colback=white,colframe=popink,boxrule=1.3pt,arc=5pt,boxsep=0pt,nobeforeafter,
        left=3pt,right=3pt,top=3pt,bottom=3pt,valign=center]{\qrcode[height=1.9cm]{https://play.google.com/store/apps/details?id=com.luvvix.onbuch&hl=fr}}%
      \hspace{8pt}\begin{minipage}[c]{0.5\linewidth}
        {\popdisplay\fontsize{11}{12.5}\selectfont\color{white}\MakeUppercase{Télécharge OnBuch}}\par\vspace{4pt}
        {\raggedright\popsemi\fontsize{8.4}{11}\selectfont\color{white}Gratuit \textbullet\ Google Play\\Tuteur IA, annales, concours, résultats\par}
      \end{minipage}
    \end{tcolorbox}
  \end{minipage}\hfill
  \begin{minipage}[t]{0.487\linewidth}
    \begin{tcolorbox}[enhanced,colback=poppurple,colframe=popink,boxrule=2.2pt,arc=10pt,
        drop shadow=popink,left=11pt,right=11pt,top=10pt,bottom=10pt,height=3.1cm,valign=center]
      \tikz[baseline=-0.7ex]\node[circle,fill=white,draw=popink,line width=1.3pt,minimum size=1.6cm,inner sep=0]{\popdisplay\fontsize{24}{24}\selectfont\color{poppurple}+};%
      \hspace{8pt}\begin{minipage}[c]{0.6\linewidth}
        {\popdisplay\fontsize{11}{12.5}\selectfont\color{white}\MakeUppercase{Passe à OnBuch+}}\par\vspace{4pt}
        {\raggedright\popsemi\fontsize{8.4}{11}\selectfont\color{white}Tous les cours signés, Tuteur IA illimité, fiches PDF — onglet OnBuch+ de l'app\par}
      \end{minipage}
    \end{tcolorbox}
  \end{minipage}
  \vspace{8pt}
  \popcontenu
  \vfill
  \signatureonbuch
  \restoregeometry
  \newpage
}

% Rangée « Ce cours contient » (page de garde)
\newcommand{\poptile}[3]{% #1 couleur #2 gros texte #3 légende
  \begin{tcolorbox}[enhanced,colback=#1,colframe=popink,boxrule=2pt,arc=9pt,shadow={2.5pt}{-2.5pt}{0pt}{popink},
      left=6pt,right=6pt,top=9pt,bottom=9pt,halign=center,valign=center,height=1.75cm,width=\linewidth,nobeforeafter]
    {\popdisplay\fontsize{12.5}{13}\selectfont\color{white}\MakeUppercase{#2}}\par\vspace{3pt}
    {\centering\popsemi\fontsize{7.8}{9.5}\selectfont\color{white}#3\par}
  \end{tcolorbox}}
\newcommand{\popcontenu}{%
  \par\noindent{\popxboldls\fontsize{7.5}{7.5}\selectfont\color{popmuted}CE COURS SIGNÉ CONTIENT}\par\vspace{7pt}
  \noindent\begin{minipage}[t]{0.235\linewidth}\poptile{popblue}{Cours}{complet, illustré, pas à pas}\end{minipage}\hfill
  \begin{minipage}[t]{0.235\linewidth}\poptile{poporange}{Méthodes}{et exercices résolus}\end{minipage}\hfill
  \begin{minipage}[t]{0.235\linewidth}\poptile{poppink}{Exercices}{type examen + corrigés}\end{minipage}\hfill
  \begin{minipage}[t]{0.235\linewidth}\poptile{popgreen}{Fiche bilan}{l'essentiel à retenir}\end{minipage}\par}

% Sommaire
\newtcolorbox{sommaire}{enhanced,colback=white,colframe=popink,boxrule=2.2pt,arc=10pt,
  drop shadow=popink,left=18pt,right=18pt,top=13pt,bottom=13pt,before skip=6pt,after skip=20pt,
  before upper={\poppill{Sommaire}{poppurple}{white}\par\vspace{9pt}\popsemi\fontsize{10.6}{14.5}\selectfont\color{popink}}}

% Signature de fin de cours — poussée tout en bas de la dernière page
\newcommand{\findecours}{%
  \par\vfill\nopagebreak
  \begin{center}\popsquiggle{poporange}\end{center}\vspace{4pt}
  \signatureonbuch
}
\AtBeginDocument{\def\llbracket{\mathopen{[\mkern-3.5mu[}}\def\rrbracket{\mathclose{]\mkern-3.5mu]}}} % intervalles d'entiers (glyphe ⟦ absent de la police)
% Glyphes absents de Fira Math : redéfinis à partir de glyphes disponibles
\AtBeginDocument{%
  \def\setminus{\mathbin{\backslash}}\let\smallsetminus\setminus
  \def\vdots{\mathord{\vbox{\baselineskip3.2pt\lineskiplimit0pt\kern4pt\hbox{.}\hbox{.}\hbox{.}}}}%
  \def\ddots{\mathinner{\mkern1mu\raise6.5pt\hbox{.}\mkern2mu\raise3.6pt\hbox{.}\mkern2mu\raise0.7pt\hbox{.}\mkern1mu}}%
  \def\triangle{\mathord{\scalebox{0.9}{$\symup{\Delta}$}}}%
  \def\nabla{\mathord{\raisebox{\depth}{\scalebox{1}[-1]{$\symup{\Delta}$}}}}%
  \def\checkmark{\popcheck{popdark}}%
  \def\star{\mathbin{\ast}}\def\bigstar{\mathord{\ast}}%
  \def\diamond{\mathbin{\scalebox{0.6}{$\Diamond$}}}%
  \def\bigcup{\mathop{\mathchoice{\raisebox{-0.3em}{\scalebox{1.7}{$\cup$}}}{\scalebox{1.25}{$\cup$}}{\cup}{\cup}}\displaylimits}%
  \def\bigcap{\mathop{\mathchoice{\raisebox{-0.3em}{\scalebox{1.7}{$\cap$}}}{\scalebox{1.25}{$\cap$}}{\cap}{\cap}}\displaylimits}%
  \def\lesseqqgtr{\mathrel{\lessgtr}}\def\bigoplus{\mathop{\oplus}\displaylimits}%
}
\providecommand{\ssi}{si et seulement si} % abréviation fréquente
\AtBeginDocument{\providecommand{\wideparen}{\overparen}} % arc de cercle

%% FIGFIX-BEGIN v3 — correctifs globaux des figures (docs/onbuch-generation/tools/figfix)
\usepackage{adjustbox}\usepackage{etoolbox}
% toute figure TikZ/pgfplots plus large que la ligne est réduite à la largeur disponible
\BeforeBeginEnvironment{tikzpicture}{\begin{adjustbox}{max width=\linewidth}}
\AfterEndEnvironment{tikzpicture}{\end{adjustbox}}
% légendes pgfplots lisibles par-dessus les courbes
\pgfplotsset{every axis/.append style={legend style={fill=white,fill opacity=0.92,text opacity=1,draw=black!15,inner sep=3pt}}}
% étiquettes d'arêtes (syntaxe edge["..."]) avec fond blanc
\tikzset{every edge quotes/.append style={fill=white,inner sep=1.2pt}}
% bulles de cartes mentales : texte à la ligne dans la bulle, bulle un peu plus grande
\tikzset{every concept/.append style={text width=2.0cm,align=center,minimum size=2.3cm,inner sep=2pt}}
%% FIGFIX-END
\begin{document}

\pagegarde{III : formuler les idées essentielles d’un texte à résumer}{Français}{Tle Tech}{Français – classe de Terminale technique}{N15}{10 séances (fiche de progression harmonisée)}{Cette leçon outille l'élève de Terminale technique pour identifier la structure argumentative d'un texte, en dégager les idées essentielles (contenus manifestes et latents) et les reformuler fidèlement dans un résumé personnel, grâce aux techniques de réduction et aux outils de liaison. Elle s'appuie sur les lectures méthodiques n°4 et n°5 du Vieux Nègre et la médaille de Ferdinand Oyono.
\vspace{4pt}
\begin{multicols}{2}\small
\begin{itemize}
\item À la fin de cette leçon, je sais identifier la structure argumentative d'un texte (thèse, arguments, exemples).
\item Je sais distinguer les contenus manifestes des contenus latents dans un énoncé.
\item Je sais repérer les idées essentielles d'un texte et les hiérarchiser.
\item Je sais appliquer les techniques de réduction (suppression, généralisation, condensation).
\item Je sais reformuler des idées sans emprunter les formulations de l'auteur.
\item Je sais utiliser correctement les conjonctions de coordination et de subordination pour lier les idées.
\end{itemize}\end{multicols}}

\begin{sommaire}
\begin{multicols}{2}
\begin{enumerate}
\item Le texte argumentatif : structure et enjeux
\item Contenus manifestes et contenus latents
\item Lecture méthodique du Vieux Nègre et la médaille (n°4 et n°5)
\item Les techniques de réduction du texte
\item Les outils de liaison : conjonctions de coordination et de subordination
\item La reformulation des idées essentielles
\item Compte rendu et remédiation
\item Activité d'intégration
\item Méthodes et exercices résolus
\item Exercices
\item Corrigés des exercices
\item Fiche bilan
\end{enumerate}
\end{multicols}
\end{sommaire}

\begin{objectifs}
\begin{itemize}
\item À la fin de cette leçon, je sais identifier la structure argumentative d'un texte (thèse, arguments, exemples).
\item Je sais distinguer les contenus manifestes des contenus latents dans un énoncé.
\item Je sais repérer les idées essentielles d'un texte et les hiérarchiser.
\item Je sais appliquer les techniques de réduction (suppression, généralisation, condensation).
\item Je sais reformuler des idées sans emprunter les formulations de l'auteur.
\item Je sais utiliser correctement les conjonctions de coordination et de subordination pour lier les idées.
\item Je sais rédiger un résumé respectant la consigne de longueur et la fidélité au texte.
\item Je sais justifier ma démarche de résumé et auto-évaluer ma production.
\end{itemize}
\end{objectifs}

\begin{prerequis}
\begin{itemize}
\item La notion de texte argumentatif (thèse, arguments, exemples) vue en classe de Première.
\item Les connecteurs logiques de base (d'abord, ensuite, cependant, donc...).
\item La lecture méthodique : repérage du thème, du ton et de l'intention de l'auteur.
\item La phrase simple et la phrase complexe (propositions coordonnées et subordonnées).
\item La lecture des premières séquences du Vieux Nègre et la médaille (lectures méthodiques n°1 à 3).
\end{itemize}
\end{prerequis}

\newpage

\coursec{Le texte argumentatif : structure et enjeux}

\courssub{Situation d'accroche et problématique}

Imaginez un journaliste de la CRTV, studio 1, lumière rouge allumée. Il a trente secondes — \emph{trente secondes} — pour résumer à la nation un discours ministériel de vingt minutes sur le développement de la région du Nord. Pas une seconde de trop. S'il noie l'essentiel sous les détails, le message passe à côté. S'il trahit la pensée du ministre, c'est sa crédibilité qui s'effondre.

Dans la salle de classe, à Garoua, un élève de Terminale technique fait face à un texte de six cents mots. Consigne : en rendre l'essentiel en cent cinquante mots, sans copier, sans déformer. Même enjeu : distinguer le squelette de la chair, l'argument de l'illustration, la thèse du bruit.

\begin{aretenir}[Problématique]
Comment identifier la thèse d'un texte argumentatif ? Comment hiérarchiser arguments et exemples pour ne garder que l'indispensable ? Comment reformuler sans trahir ? C'est toute la compétence de \cle{contraction de texte}, épreuve clé du Probatoire et du Baccalauréat technique.
\end{aretenir}

\courssub{Qu'est-ce qu'un texte argumentatif ?}

\begin{definition}[Texte argumentatif]
Un \textbf{texte argumentatif} est un discours qui vise à \textbf{faire admettre une thèse} (une opinion, une prise de position) à un destinataire, en mobilisant des \textbf{arguments} (raisons logiques) et des \textbf{exemples} (illustrations concrètes), dans le but de \textbf{convaincre} (par la raison) ou de \textbf{persuader} (par l'émotion, les valeurs).
\end{definition}

Contrairement au texte explicatif (qui informe de manière neutre) ou au texte narratif (qui raconte), le texte argumentatif \emph{prend parti}. L'auteur n'est pas transparent : il est \textbf{engagé}. On le repère à trois indices majeurs :
\begin{itemize}
    \item Une \cle{thèse} affirmée ou suggérée (l'idée-force).
    \item Des \cle{arguments} qui la soutiennent (raisons générales, causes, conséquences, analogies, arguments d'autorité).
    \item Des \cle{exemples} qui l'ancrent dans le réel (faits précis, statistiques, anecdotes, citations).
\end{itemize}

\begin{exemplebox}[Distinction explicatif / argumentatif]
\textbf{Explicatif} : « Le Cameroun compte dix régions. La région de l'Extrême-Nord est l'une des plus peuplées. » (Constat neutre, sans prise de position.)
\textbf{Argumentatif} : « Parce que l'Extrême-Nord concentre une forte population sur un territoire fragile, l'État doit y prioriser l'investissement hydraulique. » (Thèse + argument causal + conséquence normative : « l'État doit ».)
\end{exemplebox}

\courssub{Les trois composantes de l'argumentation : thèse, arguments, exemples}

Toute argumentation solide repose sur une architecture en trois étages. Visualisons-la :

\begin{popfigure}
\begin{tikzpicture}[scale=0.9, every node/.style={font=\small}]
    \colorlet{thesiscol}{poppurple!80}
    \colorlet{argcol}{popblue!70}
    \colorlet{excol}{popgreen!70}

    \node[draw=thesiscol, fill=thesiscol!15, thick, rounded corners=6pt, minimum width=7cm, minimum height=1.2cm, align=center] (thesis) at (0,3)
        {\textbf{THÈSE}\\(Idée principale défendue)\\\emph{« L'électrification solaire est la solution pour l'Extrême-Nord. »}};

    \node[draw=argcol, fill=argcol!15, thick, rounded corners=6pt, minimum width=6cm, minimum height=1.6cm, align=center] (arg1) at (-3.2,1)
        {\textbf{ARGUMENT 1}\\(Raison générale : causalité)\\\emph{« Le soleil est gratuit, abondant, non polluant. »}};
    \node[draw=argcol, fill=argcol!15, thick, rounded corners=6pt, minimum width=6cm, minimum height=1.6cm, align=center] (arg2) at (3.2,1)
        {\textbf{ARGUMENT 2}\\(Raison générale : économie)\\\emph{« Coût amorti en quelques années, pas de carburant. »}};

    \node[draw=excol, fill=excol!15, thick, rounded corners=6pt, minimum width=4.2cm, minimum height=1.2cm, align=center] (ex1) at (-4.8,-1)
        {\textbf{EXEMPLE 1}\\\emph{Village de Mokolo : foyers équipés.}};
    \node[draw=excol, fill=excol!15, thick, rounded corners=6pt, minimum width=4.2cm, minimum height=1.2cm, align=center] (ex2) at (0,-1)
        {\textbf{EXEMPLE 2}\\\emph{Centre de santé : vaccins conservés.}};
    \node[draw=excol, fill=excol!15, thick, rounded corners=6pt, minimum width=4.2cm, minimum height=1.2cm, align=center] (ex3) at (4.8,-1)
        {\textbf{EXEMPLE 3}\\\emph{École : éclairage des classes le soir.}};

    \draw[thick, -{Stealth[length=3mm]}, thesiscol] (thesis.south) -- ++(0,-0.3) -| (arg1.north);
    \draw[thick, -{Stealth[length=3mm]}, thesiscol] (thesis.south) -- ++(0,-0.3) -| (arg2.north);
    \draw[thick, -{Stealth[length=3mm]}, argcol] (arg1.south) -- ++(0,-0.3) -| (ex1.north);
    \draw[thick, -{Stealth[length=3mm]}, argcol] (arg1.south) -- ++(0,-0.3) -| (ex2.north);
    \draw[thick, -{Stealth[length=3mm]}, argcol] (arg2.south) -- ++(0,-0.3) -| (ex2.north);
    \draw[thick, -{Stealth[length=3mm]}, argcol] (arg2.south) -- ++(0,-0.3) -| (ex3.north);
\end{tikzpicture}
\legende{Schéma pyramidal de l'argumentation : la thèse repose sur des arguments, eux-mêmes étayés par des exemples concrets.}
\end{popfigure}

\begin{definition}[Thèse]
La \textbf{thèse} est l'idée principale que l'auteur cherche à faire admettre au lecteur. C'est le \textbf{point de vue} défendu, la réponse à la question : « \emph{Que veut me faire penser l'auteur ?} »
\end{definition}

\begin{definition}[Argument]
Un \textbf{argument} est une \textbf{raison générale} avancée pour justifier la thèse. Il répond à la question « \emph{Pourquoi ?} ». Il a valeur de loi, de principe, de mécanisme causal.
\end{definition}

\begin{definition}[Exemple]
Un \textbf{exemple} est un \textbf{cas particulier}, une illustration concrète (fait, chiffre, anecdote, citation) qui donne corps à l'argument. Il répond à la question « \emph{Montre-moi.} »
\end{definition}

\begin{attention}[Piège classique : confondre argument et exemple]
Ne pas écrire : « L'argument, c'est que le village de Mokolo a des panneaux solaires. » \textbf{Faux.} C'est un \emph{exemple}.
L'argument, c'est : « L'énergie solaire est adaptée aux zones rurales isolées. »
L'exemple \emph{illustre} l'argument ; il ne le \emph{remplace} pas. Au Bac, si vous citez l'exemple comme argument, vous perdez des points de \cle{structure argumentative}.
\end{attention}

\begin{exemplebox}[Plaidoyer pour l'électrification rurale à l'Extrême-Nord]
\begin{itemize}
    \item \textbf{Thèse} : « L'État doit généraliser le solaire photovoltaïque dans l'Extrême-Nord. »
    \item \textbf{Argument 1 (écologique)} : « Le solaire est une énergie propre, sans émission de CO$_2$. »
    \item \textbf{Exemple 1} : « À Mokolo, des foyers ont remplacé leurs lampes à pétrole par des kits solaires, réduisant leurs émissions de CO$_2$. »
    \item \textbf{Argument 2 (économique)} : « L'investissement initial est amorti en quelques années, puis l'énergie est gratuite. »
    \item \textbf{Exemple 2} : « Un centre de santé équipé a fortement réduit sa facture d'énergie en deux ans. »
\end{itemize}
\end{exemplebox}

\courssub{Thèse explicite et thèse implicite : comment les repérer}

La thèse n'est pas toujours écrite en toutes lettres. Deux cas de figure :

\begin{propriete}[Thèse explicite]
La thèse est \textbf{formulée textuellement} par l'auteur. Repérage : phrases affirmatives au présent de vérité générale, souvent en début ou en fin de texte, introduites par \emph{je pense que, il est évident que, il faut que, nous devons, la solution est...}
\end{propriete}

\begin{propriete}[Thèse implicite]
La thèse \textbf{n'est pas énoncée} ; elle se déduit de l'ensemble des arguments. L'auteur laisse le lecteur « faire le chemin ». Repérage : convergence de tous les arguments vers une même conclusion non dite.
\end{propriete}

\begin{methode}[Repérer la thèse (explicite ou implicite)]
\begin{enumerate}
    \item \textbf{Lisez le titre, l'introduction, la conclusion.} La thèse y est souvent explicite.
    \item \textbf{Repérez les arguments} (connecteurs logiques : \emph{car, parce que, donc, ainsi, en effet...}).
    \item \textbf{Demandez-vous :} « \emph{Tous ces arguments servent-ils à prouver quoi ?} » $\rightarrow$ C'est la thèse.
    \item \textbf{Formulez-la en une phrase complète} (sujet + verbe + complément), au présent, à la 3\textsuperscript{e} personne : « L'auteur soutient que... »
\end{enumerate}
\end{methode}

\begin{exoresolu}[Repérage de thèse implicite]
\textbf{Texte (extrait fictif, style éditorial) :}
\begin{quote}
« Les coupures d'électricité répétées à Maroua paralysent les PME. Les soudeurs, menuisiers, couturières perdent des commandes. Les élèves révisent à la lampe torche. Les vaccins du centre de santé risquent la rupture de la chaîne du froid. Pourtant, le soleil brille presque toute l'année. Des kits solaires existent, abordables, installables en une journée. »
\end{quote}
\tcblower
\textbf{Démarche :}
\begin{enumerate}
    \item Arguments repérés : « coupures paralysent les PME », « élèves révisent à la lampe », « vaccins en danger » (constats négatifs) ; « soleil presque toute l'année », « kits abordables, installables vite » (constats positifs / solution).
    \item Question : « Tout cela prouve quoi ? »
    \item \textbf{Thèse implicite formulée :} « L'auteur soutient que l'énergie solaire est la solution urgente et réaliste à la crise électrique de Maroua. »
\end{enumerate}
\end{exoresolu}

\courssub{Les marques de l'argumentation : connecteurs, champs lexicaux, modalisateurs}

Un texte argumentatif se reconnaît à sa \textbf{texture linguistique}. Trois leviers principaux :

\textbf{1. Les connecteurs logiques (outils de liaison)}\par

Ils organisent le raisonnement, signalent la relation entre les propositions. Le tableau ci-dessous répertorie les plus fréquents par fonction logique.

\begin{center}
\begin{tabularx}{\linewidth}{|l|X|X|}\hline
\rowcolor{popblueL}
\textbf{Fonction logique} & \textbf{Connecteurs principaux} & \textbf{Exemple d'usage} \\ \hline
\textbf{Addition / Enchaînement} & \emph{et, de plus, en outre, par ailleurs, également, non seulement... mais encore, d'une part... d'autre part} & « Le solaire est propre, \textbf{de plus} il est rentable. » \\ \hline
\textbf{Opposition / Concession} & \emph{mais, cependant, pourtant, toutefois, néanmoins, alors que, bien que, malgré, en dépit de} & « \textbf{Bien que} l'investissement initial soit lourd, l'amortissement est rapide. » \\ \hline
\textbf{Cause / Explication} & \emph{car, parce que, comme, puisque, en effet, du fait de, grâce à, à cause de} & « Les PME ferment \textbf{car} l'électricité manque. » \\ \hline
\textbf{Conséquence / Conclusion} & \emph{donc, ainsi, par conséquent, c'est pourquoi, alors, si bien que, en somme, en définitive} & « Le soleil est gratuit ; \textbf{donc} l'énergie solaire coûte moins cher à l'usage. » \\ \hline
\textbf{Illustration / Exemple} & \emph{par exemple, ainsi, notamment, c'est le cas de, tel que, comme en témoigne} & « Des villages s'équipent, \emph{par exemple} Mokolo et Koza. » \\ \hline
\textbf{Justification / Preuve} & \emph{en effet, en fait, effectivement, force est de constater que} & « La solution existe. \textbf{En effet}, des ONG distribuent déjà des kits. » \\ \hline
\end{tabularx}
\end{center}

\begin{aretenir}[Astuce Bac]
Soulignez \textbf{tous les connecteurs} à la première lecture. Ils dessinent le \textbf{squelette logique} du texte. En contraction, vous devez les conserver (ou les remplacer par des équivalents) pour garder la cohérence.
\end{aretenir}

\textbf{2. Les champs lexicaux : l'argumentation par les mots}\par

Le choix du vocabulaire n'est jamais neutre. L'auteur construit des \textbf{champs lexicaux} (ensembles de mots liés au même thème) pour orienter la perception.

\begin{exemplebox}[Champs lexicaux dans un éditorial pro-solaire]
\begin{itemize}
    \item \textbf{Champ lexical de la \emph{pénurie / souffrance} (réalité actuelle) :} coupures, paralysie, obscurité, lampe torche, risque, perte, urgence, criant, insupportable.
    \item \textbf{Champ lexical de l'\emph{abondance / solution} (solaire) :} soleil, rayonnement, gratuit, inépuisable, propre, autonome, rapide, accessible, avenir.
    \item \textbf{Effet :} le contraste lexical crée une \textbf{tension argumentative} : le problème est noir, la solution est lumineuse. Le lecteur \emph{ressent} la nécessité avant même de raisonner.
\end{itemize}
\end{exemplebox}

\textbf{3. Les modalisateurs : la marque de l'engagement}\par

Les \textbf{modalisateurs} expriment le degré de certitude, de doute, d'obligation ou de valeur de l'auteur. Ils signalent : « \emph{Je m'engage.} »

\begin{center}
\begin{tabularx}{\linewidth}{|l|X|X|}\hline
\rowcolor{poppurpleL}
\textbf{Type} & \textbf{Outils linguistiques} & \textbf{Effet argumentatif} \\ \hline
\textbf{Épistémique (certitude / doute)} & \emph{certes, assurément, sans doute, probablement, il semble que, sans conteste} & Affirme la solidité des arguments (« \emph{Sans conteste}, le solaire marche. ») ou ménage une réserve. \\ \hline
\textbf{Déontique (obligation / permission)} & \emph{il faut que, il est nécessaire que, devoir, falloir, impératif, indispensable, urgent} & Appelle à l'action (« \emph{Il est impératif} d'équiper les centres de santé. »). \\ \hline
\textbf{Axiologique (valeur / jugement)} & Adjectifs évaluatifs : \emph{bénéfique, catastrophique, vital, ruineux, prometteur, scandaleux, légitime} & Juge la situation (« \emph{Ce retard est scandaleux.} »). \\ \hline
\textbf{Énonciatif (personne / impersonnel)} & \emph{Je pense que, nous estimons, selon les experts, il est admis que} & Ancre la responsabilité (je) ou la légitime par l'autorité (experts). \\ \hline
\end{tabularx}
\end{center}

\begin{attention}[Piège : modalisateurs \emph{affaiblissants}]
« \emph{Peut-être} le solaire aiderait. » $\rightarrow$ Thèse affaiblie.
Au Bac, repérez si l'auteur \textbf{s'engage franchement} (modalisateurs forts : \emph{il faut, certainement, indispensable}) ou s'il \emph{tergiverse}. Cela change la force de la thèse à résumer.
\end{attention}

\courssub{Schéma de la structure argumentative type}

Un texte argumentatif bien construit suit un parcours balisé. Le connaître permet d'anticiper la localisation de la thèse et des arguments.

\begin{popfigure}
\begin{tikzpicture}[scale=0.85, every node/.style={font=\small}]
    \tikzset{
        box/.style={draw=popdark, thick, fill=popcream, rounded corners=4pt, minimum height=1.1cm, align=center, text width=3.4cm},
        arrow/.style={thick, -{Stealth[length=3mm]}, popdark},
        labelbox/.style={font=\small\bfseries, text=popblue, anchor=south}
    }

    \node[box, fill=poppurple!20, align=center] (intro) at (0,4){\textbf{INTRODUCTION}\\\textit{Accroche + sujet posé + thèse (souvent ici)}};
    \node[box, align=center] (dev1) at (-4.4,1.6){\textbf{DÉVELOPPEMENT 1}\\\textit{Argument 1 + exemple(s)}};
    \node[box, align=center] (dev2) at (0,1.6){\textbf{DÉVELOPPEMENT 2}\\\textit{Argument 2 + exemple(s)}};
    \node[box, align=center] (dev3) at (4.4,1.6){\textbf{DÉVELOPPEMENT 3}\\\textit{Argument 3 (+ contre-argument réfuté)}};
    \node[box, fill=popgreen!20, align=center] (conc) at (0,-0.8){\textbf{CONCLUSION}\\\textit{Bilan + réaffirmation de la thèse + ouverture / appel à l'action}};

    \draw[arrow] (intro.south) -- ++(0,-0.2) -| (dev1.north);
    \draw[arrow] (intro.south) -- (dev2.north);
    \draw[arrow] (intro.south) -- ++(0,-0.2) -| (dev3.north);
    \draw[arrow] (dev1.south) -- ++(0,-0.2) -| (conc.north);
    \draw[arrow] (dev2.south) -- (conc.north);
    \draw[arrow] (dev3.south) -- ++(0,-0.2) -| (conc.north);

    \node[labelbox] at (intro.north) {Annoncer la thèse};
    \node[labelbox] at (dev1.north) {Prouver};
    \node[labelbox] at (dev2.north) {Prouver};
    \node[labelbox] at (dev3.north) {Approfondir / réfuter};
    \node[labelbox, anchor=north] at (conc.south) {Conclure / agir};
\end{tikzpicture}
\legende{Structure canonique d'un texte argumentatif : introduction (thèse), développement (arguments + exemples), conclusion (bilan + réaffirmation).}
\end{popfigure}

\begin{aretenir}[Variantes fréquentes au Bac]
\begin{itemize}
    \item \textbf{Thèse en conclusion seulement} (structure inductive) : arguments d'abord, thèse à la fin comme évidence.
    \item \textbf{Thèse au milieu} (structure encadrée) : arguments pour, puis thèse, puis arguments contre réfutés.
    \item \textbf{Contre-argumentation} : l'auteur anticipe une objection (« \emph{On objectera que...} ») et la réfute (« \emph{Or... / Cependant...} »). Ne la confondez pas avec la thèse !
\end{itemize}
\end{aretenir}

\courssub{Application : repérage dans un éditorial}

Exerçons-nous sur un extrait réaliste (style éditorial, domaine : développement régional / énergie).

\begin{exemplebox}[Extrait d'éditorial — « L'Extrême-Nord ne peut plus attendre » (texte fictif)]
\begin{quote}
\textbf{Titre :} \textit{L'Extrême-Nord ne peut plus attendre}

\textbf{Introduction :} Alors que le pays s'engage résolument dans l'industrialisation, la région de l'Extrême-Nord demeure le parent pauvre de l'électrification. \textbf{Il est temps d'affirmer avec force que l'énergie solaire photovoltaïque constitue la réponse structurelle, rapide et souveraine à ce déficit historique.}

\textbf{Développement :} \textbf{En effet}, le potentiel solaire y est exceptionnel : la région bénéficie d'un ensoleillement parmi les plus élevés du pays. \textbf{De plus}, les technologies de stockage sont désormais matures et leurs coûts ont fortement chuté en dix ans. \textbf{Par conséquent}, l'argument du « coût prohibitif » ne tient plus. \textbf{C'est le cas} du village de \textbf{Mokolo}, où une mini-centrale solaire alimente des foyers, une école et un centre de santé, sans aucune coupure. \textbf{Par ailleurs}, l'approche décentralisée évite les pertes en ligne du réseau national et crée des emplois locaux de maintenance. \textbf{Toutefois}, certains objecteront que la poussière et les pluies endommagent les panneaux. \textbf{Or}, les normes actuelles garantissent leur étanchéité ; \textbf{en outre}, un nettoyage régulier suffit.

\textbf{Conclusion :} \textbf{En définitive}, l'Extrême-Nord a tout pour devenir le laboratoire national de l'autonomie énergétique. \textbf{L'État doit donc lancer, dès la prochaine loi de finances, un programme massif « Solaire Nord ». L'avenir de la région ne saurait attendre le passage du réseau national.}
\end{quote}
\end{exemplebox}

\begin{exoresolu}[Analyse structurelle de l'éditorial]
\textbf{Consigne :} repérez la thèse, les arguments, les exemples et les connecteurs de cet éditorial.
\tcblower
\textbf{Étape 1 — Thèse (explicite, en introduction) :}
\begin{quote}
« L'énergie solaire photovoltaïque constitue la réponse structurelle, rapide et souveraine au déficit électrique de l'Extrême-Nord. »
\end{quote}

\textbf{Étape 2 — Arguments (développement) :}
\begin{enumerate}
    \item \textbf{Argument du potentiel naturel} : « Potentiel solaire exceptionnel. » \emph{[Connecteur : En effet]}
    \item \textbf{Argument technico-économique} : « Technologies matures, coûts en forte baisse. » \emph{[De plus, Par conséquent]}
    \item \textbf{Argument technique (réseau)} : « Décentralisation = moins de pertes, emplois locaux. » \emph{[Par ailleurs]}
    \item \textbf{Réfutation d'objection} : « Poussière et pluies ? Panneaux étanches, maintenance faible. » \emph{[Toutefois, Or, En outre]}
\end{enumerate}

\textbf{Étape 3 — Exemple (ancrage concret) :}
\begin{itemize}
    \item \textbf{Exemple unique mais dense} : village de Mokolo — mini-centrale solaire — foyers + école + centre de santé — zéro coupure. \emph{[C'est le cas de]}
\end{itemize}

\textbf{Étape 4 — Conclusion (réaffirmation + appel à l'action) :}
\begin{quote}
« L'Extrême-Nord = laboratoire de l'autonomie énergétique. $\rightarrow$ Programme "Solaire Nord" dès la prochaine loi de finances. »
\end{quote}

\textbf{Étape 5 — Connecteurs repérés (squelette logique) :}
\emph{En effet, De plus, Par conséquent, C'est le cas, Par ailleurs, Toutefois, Or, En outre, En définitive, Donc.}
\end{exoresolu}

\begin{methode}[Résumer cet éditorial en 3 phrases (entraînement à la contraction)]
\begin{enumerate}
    \item \textbf{Thèse + arguments clés fusionnés :} L'éditorial soutient que l'énergie solaire, portée par un ensoleillement remarquable et des coûts en chute, est la solution structurelle au déficit électrique de l'Extrême-Nord.
    \item \textbf{Preuve par l'exemple + argument technique :} L'expérience de Mokolo (mini-centrale, zéro coupure) prouve la fiabilité ; la décentralisation évite les pertes du réseau et crée de l'emploi local.
    \item \textbf{Conclusion + appel :} Les objections étant réfutées, l'auteur exige un programme massif « Solaire Nord » dès la prochaine loi de finances.
\end{enumerate}
\emph{(Total : environ 85 mots. Essentiel préservé, détails secondaires — normes d'étanchéité, nettoyage — éliminés.)}
\end{methode}

\courssub{Synthèse de la section : ce qu'il faut retenir pour la contraction}

\begin{aretenir}[Check-list « Structure argumentative » — À avoir en tête pendant l'épreuve]
\begin{enumerate}
    \item \textbf{Titre + introduction + conclusion} $\rightarrow$ zone prioritaire pour la \textbf{thèse}.
    \item \textbf{Connecteurs logiques} (car, donc, cependant, en effet...) $\rightarrow$ ils balisent les \textbf{arguments}.
    \item \textbf{Champs lexicaux contrastés} (problème / solution) $\rightarrow$ ils valident l'orientation de la thèse.
    \item \textbf{Modalisateurs forts} (il faut, indispensable, sans conteste) $\rightarrow$ ils mesurent l'engagement de l'auteur.
    \item \textbf{Exemples} (chiffres, noms propres, dates, lieux) $\rightarrow$ ne garder que \textbf{le plus représentatif} par argument.
    \item \textbf{Contre-argumentation} (Toutefois... Or...) $\rightarrow$ la mentionner \textbf{une seule fois} si elle est majeure : « L'auteur réfute l'objection du coût et de la maintenance. »
\end{enumerate}
\end{aretenir}

\begin{exercice}[Entraînement autonome]
À partir de l'éditorial ci-dessus, rédigez une \textbf{contraction de 100 mots $\pm$ 10\,\%} respectant :
\begin{itemize}
    \item la thèse explicite reformulée ;
    \item les 3 arguments principaux (potentiel, coût/technologie, décentralisation) + la réfutation ;
    \item l'exemple de Mokolo (condensé) ;
    \item l'appel final ;
    \item \textbf{zéro copier-coller} : reformulation obligatoire ;
    \item connecteurs logiques conservés ou remplacés par des équivalents.
\end{itemize}
\textit{(Corrigé type en fin de leçon, section 7.)}
\end{exercice}

\courssub{Pont vers la section 2}

Nous savons maintenant \textbf{décortiquer} la structure visible (manifeste) d'un texte argumentatif. Mais tout ne se dit pas explicitement. Derrière les mots, il y a les \textbf{non-dits}, les \textbf{présupposés}, les \textbf{valeurs cachées} qui orientent le raisonnement sans apparaître. C'est le domaine des \textbf{contenus latents}. Section 2 : \emph{Contenus manifestes et contenus latents}.

\coursec{Contenus manifestes et contenus latents}

Dans la section précédente, nous avons découvert la structure du texte argumentatif : une thèse, des arguments, des exemples, et des procédés pour convaincre ou persuader. Mais argumenter, ce n'est pas toujours dire clairement ce que l'on pense. Un auteur habile — et Ferdinand Oyono en est un maître — peut défendre une idée en prétendant défendre son contraire, dénoncer une injustice sans jamais la nommer. Pour résumer correctement un tel texte, il ne suffit donc pas de relever ce qui est écrit : il faut comprendre ce qui est \emph{voulu}. C'est tout l'enjeu de la distinction entre \cle{contenu manifeste} et \cle{contenu latent}.

\courssub{Le contenu manifeste : ce qui est dit}

\begin{definition}[Contenu manifeste]
Le \cle{contenu manifeste} d'un texte est l'ensemble des informations \textbf{explicitement énoncées} : ce que les mots signifient littéralement, au premier degré. C'est le sens que tout lecteur attentif peut relever directement, sans interprétation : les faits rapportés, les personnages nommés, les actions décrites, les propos tenus.
\end{definition}

Prenons une phrase simple : « Le commandant remet une médaille à Meka pour le récompenser de ses services. » Le contenu manifeste est limpide : il y a un commandant, un vieil homme nommé Meka, une médaille, une cérémonie de récompense. Aucun lecteur ne peut contester ces faits : ils sont \emph{dans} le texte.

Le contenu manifeste se repère par des questions simples : \emph{qui ? quoi ? où ? quand ? comment ?} C'est le niveau de lecture qu'exige la première étape du résumé (la lecture analytique), mais ce n'est pas le seul.

\courssub{Le contenu latent : ce qui est suggéré}

\begin{definition}[Contenu latent]
Le \cle{contenu latent} est le \textbf{sens caché} d'un texte, que le lecteur \textbf{reconstruit} à partir d'indices : sous-entendus, connotations, ironie, présupposés, non-dits. Il n'est jamais écrit noir sur blanc, mais il constitue souvent l'\textbf{intention réelle} de l'auteur. On parle aussi de sens \emph{implicite} ou de \emph{second degré}.
\end{definition}

Reprenons notre phrase : « Le commandant remet une médaille à Meka pour le récompenser de ses services. » Si le récit nous apprend que Meka a perdu ses deux fils à la guerre et sa terre au profit de la mission, et que le commandant déclare que la France « récompense toujours ses fidèles serviteurs », le lecteur comprend autre chose : la médaille est une \textbf{compensation dérisoire}, presque une insulte. Ce sens-là n'est écrit nulle part : il est latent, et pourtant c'est le sens véritable du passage.

\begin{popfigure}
\begin{center}
\begin{tikzpicture}[scale=1]
% Ciel et mer
\fill[popblueL!60] (-5,0) rectangle (5,2.6);
\fill[popblue!45] (-5,-3.2) rectangle (5,0);
\draw[popblue, thick] (-5,0) -- (5,0);
\node[popdark, font=\small\itshape] at (4.1,0.25) {surface (le texte)};
% Iceberg partie visible
\fill[white, draw=popblue, thick] (-0.9,0) -- (-0.4,1.5) -- (0.1,2.1) -- (0.6,1.3) -- (0.9,0) -- cycle;
% Iceberg partie immergée
\fill[white, draw=popblue, thick] (-0.9,0) -- (-1.8,-1.2) -- (-1.3,-2.6) -- (0,-3.0) -- (1.5,-2.4) -- (1.9,-1.1) -- (0.9,0) -- cycle;
% Etiquettes
\node[popdark, font=\small\bfseries, align=center] at (0,0.9) {Manifeste};
\node[popdark, font=\scriptsize, align=center] at (2.9,1.4) {ce qui est dit\\(faits, mots, énoncés)};
\draw[-{Stealth}, popdark] (2.0,1.3) -- (0.6,1.1);
\node[popdark, font=\small\bfseries, align=center] at (0,-1.6) {Latent};
\node[popdark, font=\scriptsize, align=center] at (-3.4,-1.9) {ce qui est suggéré\\(ironie, connotations,\\présupposés, non-dits)};
\draw[-{Stealth}, popdark] (-1.9,-1.8) -- (-0.9,-1.7);
% Proportion
\node[popdark, font=\scriptsize\itshape, align=center] at (0,-3.6) {Comme l'iceberg, la partie visible du texte est la plus petite :\\l'essentiel du sens se cache sous la surface des mots.};
\end{tikzpicture}
\legende{Schéma de l'iceberg de la signification : le contenu manifeste (partie émergée) est directement lisible ; le contenu latent (partie immergée, bien plus vaste) doit être reconstruit par le lecteur.}
\end{center}
\end{popfigure}

\courssub{Les indices du contenu latent}

Comment le lecteur peut-il découvrir un sens qui n'est pas écrit ? En repérant des \cle{indices} que l'auteur a semés dans son texte. Quatre grandes familles d'indices sont à connaître.

\textbf{1. L'ironie.} Dire le contraire de ce que l'on pense, en laissant entendre ce contraire. Ses marqueurs : l'\textbf{exagération} (hyperbole), l'\textbf{éloge excessif} qui sonne faux, la \textbf{contradiction} entre les paroles et les faits décrits, le ton solennel appliqué à un objet insignifiant. Lorsque le commandant d'\emph{Une vie de boy} ou du \emph{Vieux Nègre et la médaille} célèbre pompeusement la « générosité » de la France coloniale devant un homme ruiné par elle, l'écart entre les mots et la réalité crée l'ironie.

\textbf{2. Les connotations.} Un mot ne possède pas seulement un sens propre (sa \emph{dénotation}) : il porte des valeurs, des émotions, des jugements. « Récompense », « fidélité », « civilisation » ont une connotation positive ; placés dans la bouche d'un personnage dont les actes contredisent ces valeurs, ils se retournent contre lui.

\textbf{3. Les présupposés.} Ce qu'un énoncé admet sans le dire. Quand le commandant déclare que Meka est « un ami de la France », il présuppose que la France est, elle, digne d'amitié et de gratitude — ce que tout le roman s'emploie à contester.

\textbf{4. Les non-dits.} Ce que le texte tait alors qu'il devrait le dire : les souffrances passées sous silence dans le discours officiel, les morts oubliés au moment des honneurs. Le silence lui-même devient signifiant.

\begin{methode}[Décoder le contenu latent en quatre étapes]
\begin{enumerate}
\item \textbf{Lire au premier degré} : relever le contenu manifeste (faits, paroles, actions).
\item \textbf{Repérer les anomalies} : mots à connotation forte, exagérations, contradictions entre les propos et la situation, décalages de ton, silences étranges.
\item \textbf{Confronter au contexte} : qui parle ? à qui ? à quelle époque ? que sait-on des faits réels évoqués ?
\item \textbf{Formuler le sens latent} : énoncer clairement, au second degré, ce que l'auteur veut réellement faire comprendre, en le justifiant par les indices relevés.
\end{enumerate}
\end{methode}

\courssub{Le rôle du contexte dans le décodage}

Le contenu latent ne se devine pas : il se \textbf{déduit}, et la déduction exige des connaissances contextuelles. Trois éléments de contexte sont décisifs.

\begin{itemize}
\item \textbf{L'auteur et son époque.} Ferdinand Oyono écrit \emph{Le Vieux Nègre et la médaille} en 1956, au Cameroun alors sous administration française, à la veille des indépendances africaines. Savoir cela permet de comprendre que le roman n'est pas une simple anecdote villageoise : c'est une \textbf{critique de la colonisation}.
\item \textbf{Le destinataire.} Le discours du commandant s'adresse officiellement à la foule réunie pour la cérémonie ; mais Oyono, lui, s'adresse au lecteur, à qui il révèle le ridicule et l'hypocrisie de ce discours. Il y a donc \emph{deux niveaux de communication} superposés.
\item \textbf{La situation d'énonciation.} Une cérémonie officielle impose un langage solennel et codé. C'est précisément ce solennel, appliqué à une médaille sans valeur remise à un homme spolié, qui produit l'effet comique et critique.
\end{itemize}

\begin{savaistu}[L'ironie, arme de dénonciation chez Oyono]
Ferdinand Oyono (1929-2010), écrivain camerounais et homme d'État, utilise l'ironie pour dénoncer l'absurdité du système colonial \textbf{sans jamais la nommer directement}. Dans \emph{Le Vieux Nègre et la médaille}, le mot « colonisation » et le mot « injustice » ne sont presque jamais prononcés : c'est le contraste entre les beaux discours officiels et la réalité vécue par Meka qui fait éclater la vérité. Cette stratégie, héritée de la satire classique (Voltaire, La Fontaine), permet à l'écrivain de toucher un large public tout en échappant partiellement à la censure coloniale.
\end{savaistu}

\courssub{Application : le discours du commandant}

\begin{exoresolu}[Décodage d'un passage ironique]
\textbf{Énoncé.} Lors de la cérémonie du 14 juillet, le commandant déclare en substance : « Meka est un ami de la France. La France n'oublie jamais ceux qui l'ont servie. Cette médaille est le signe de sa reconnaissance. » Or le lecteur sait que Meka a donné ses deux fils, morts pour la France, et sa terre, cédée aux missionnaires. Relevez le contenu manifeste de ce discours, puis décodez son contenu latent.
\tcblower
\textbf{Solution détaillée.}
\emph{Étape 1 — Contenu manifeste.} Le commandant prononce un éloge : Meka est présenté comme un ami de la France ; la France est présentée comme reconnaissante ; la médaille est présentée comme une récompense méritée.
\par\emph{Étape 2 — Indices du latent.} (a) \textbf{Contradiction} entre les paroles et les faits : la France « n'oublie jamais », mais elle a pris les fils et la terre de Meka sans rien lui rendre. (b) \textbf{Connotation} : « ami », « reconnaissance », « servie » forment un champ lexical de la gratitude qui sonne faux dans cette situation. (c) \textbf{Présupposé} : le discours admet que la médaille est une juste compensation, ce qui est dérisoire face aux sacrifices consentis. (d) \textbf{Non-dit} : le commandant ne mentionne ni les deux fils morts, ni la terre perdue.
\par\emph{Étape 3 — Contexte.} Discours officiel colonial, prononcé par un représentant de l'administration, devant une foule, le jour de la fête nationale française : le commandant célèbre en réalité la France, pas Meka.
\par\emph{Étape 4 — Sens latent formulé.} Oyono dénonce l'hypocrisie coloniale : la médaille est une compensation symbolique et dérisoire pour des sacrifices immenses ; la « reconnaissance » de la France est un masque qui cache l'exploitation et l'ingratitude. La cérémonie sert la propagande coloniale, non le vieil homme.
\end{exoresolu}

\courssub{Contenu latent et contraction de texte}

\begin{attention}[Le piège majeur du résumé]
Au résumé, le candidat qui ne retient que le contenu manifeste \textbf{trahit l'intention réelle de l'auteur}. Résumer le passage précédent par « le commandant récompense Meka, qui est un ami de la France » revient à écrire le contraire de ce qu'Oyono veut dire ! Il faut donc restituer le sens latent : « l'auteur montre, par l'ironie du discours officiel, l'hypocrisie de la prétendue reconnaissance coloniale ».
\par Mais attention au piège inverse : restituer le sens latent ne signifie pas \textbf{ajouter ses propres opinions}. Le résumeur ne doit ni condamner lui-même le commandant, ni inventer un sens que les indices ne justifient pas. Règle d'or : tout élément latent retenu doit pouvoir être \textbf{justifié par un indice précis du texte}.
\end{attention}

\begin{aretenir}[L'essentiel]
\begin{itemize}
\item Le \cle{contenu manifeste} est le sens explicite, littéral, directement relevable.
\item Le \cle{contenu latent} est le sens implicite, reconstruit à partir d'indices : ironie, connotations, présupposés, non-dits.
\item Le décodage du latent exige de tenir compte du \textbf{contexte} : auteur, époque, destinataire, situation d'énonciation.
\item Au résumé, il faut restituer le sens latent \emph{de l'auteur}, sans y mêler ses opinions personnelles.
\end{itemize}
\end{aretenir}

\begin{exercice}[Entraînement]
Dans le roman, les villageois admirent la médaille de Meka et l'envient, tandis que le lecteur comprend qu'elle ne lui apportera que des ennuis. 1) Quel est le contenu manifeste de l'attitude des villageois ? 2) Quel contenu latent l'écart entre leur admiration et la réalité fait-il apparaître ? 3) Quels indices du texte permettent ce décodage ?
\end{exercice}

\begin{corrige}[Exercice 1]
1) Contenu manifeste : les villageois considèrent la médaille comme un honneur prestigieux et envient Meka. 2) Contenu latent : Oyono montre que les valeurs coloniales ont été intériorisées par les colonisés eux-mêmes, qui admirent un symbole vide ; la médaille, loin d'être un bienfait, devient une source d'humiliations (l'arrestation, la prison). 3) Indices : le contraste entre l'enthousiasme des villageois et les malheurs qui frappent Meka juste après la cérémonie ; le décalage entre la valeur supposée de la médaille et son inutilité réelle ; l'ironie du narrateur qui laisse les faits démentir l'admiration collective.
\end{corrige}

\coursec{Lecture méthodique du Vieux Nègre et la médaille (n°4 et n°5)}

Dans la section précédente, nous avons appris à distinguer les \cle{contenus manifestes} (ce que le texte dit explicitement) des \cle{contenus latents} (ce qu'il suggère sans le dire). Cette compétence va maintenant être mise en pratique sur une œuvre complète : \emph{Le Vieux Nègre et la médaille} de Ferdinand Oyono. La \cle{lecture méthodique} est en effet l'étape indispensable qui précède tout résumé : on ne peut contracter un texte que si on l'a d'abord compris en profondeur, dans ce qu'il dit et dans ce qu'il cache.

\begin{savaistu}[Une œuvre majeure du programme]
\emph{Le Vieux Nègre et la médaille}, publié en 1956, est le troisième roman de l'écrivain camerounais Ferdinand Oyono. Il raconte l'histoire de Meka, un vieil homme qui a tout donné à la colonisation (ses terres, ses fils morts à la guerre, sa foi chrétienne) et qui reçoit en récompense une médaille. Loin d'être un honneur, cette médaille déclenche sa prise de conscience : il comprend qu'il a été dupé. Ce roman \cle{anticolonial} majeur de la littérature camerounaise et africaine est inscrit au programme du Probatoire et du Baccalauréat technique.
\end{savaistu}

\courssub{La démarche générale de la lecture méthodique}

Lire méthodiquement, c'est lire \emph{plusieurs fois}, avec un objectif différent à chaque lecture. On ne lit pas un texte argumentatif ou polémique comme on lit un fait divers : il faut d'abord saisir le sens général, puis observer comment l'auteur construit son propos, puis enfin dégager les idées essentielles.

\begin{methode}[Lire trois fois]
\begin{enumerate}
\item \textbf{Lecture globale} : lire le passage d'un trait, sans s'arrêter. Objectif : répondre aux questions \emph{qui ? quoi ? où ? quand ?} et formuler le sens général en une phrase.
\item \textbf{Lecture analytique} : relire lentement, crayon en main. Objectif : repérer les procédés (ironie, caricature, contraste, registres), souligner les mots importants, distinguer ce qui est dit (manifeste) de ce qui est suggéré (latent).
\item \textbf{Lecture synthétique} : relire une dernière fois pour dégager les \cle{idées essentielles}. Objectif : formuler chaque idée en une phrase, avec ses propres mots, puis les hiérarchiser.
\end{enumerate}
\end{methode}

\begin{popfigure}
\begin{center}
\begin{tikzpicture}[
  etape/.style={rectangle, rounded corners, draw=popblue, fill=popblueL, thick, minimum width=2.8cm, minimum height=1.1cm, align=center, font=\small\bfseries},
  descr/.style={below=2mm of etape, font=\scriptsize, text width=2.9cm, align=center, text=popdark},
  fleche/.style={-{Stealth[length=3mm]}, very thick, poporange}
]
\node[etape] (e1) at (0,0) {1. Situer};
\node[etape] (e2) at (4.2,0) {2. Lire};
\node[etape] (e3) at (8.4,0) {3. Analyser};
\node[etape] (e4) at (12.6,0) {4. Dégager le sens};
\node[below=2mm of e1, font=\scriptsize, text width=2.9cm, align=center, text=popdark] {identifier l'auteur, l'œuvre, le contexte du passage};
\node[below=2mm of e2, font=\scriptsize, text width=2.9cm, align=center, text=popdark] {lecture globale puis lecture attentive, à voix haute si possible};
\node[below=2mm of e3, font=\scriptsize, text width=2.9cm, align=center, text=popdark] {repérer thèse, arguments, procédés, contenus latents};
\node[below=2mm of e4, font=\scriptsize, text width=2.9cm, align=center, text=popdark] {formuler les idées essentielles et le sens global};
\draw[fleche] (e1) -- (e2);
\draw[fleche] (e2) -- (e3);
\draw[fleche] (e3) -- (e4);
\end{tikzpicture}
\end{center}
\legende{Frise des quatre étapes de la lecture méthodique : situer le passage dans l'œuvre, lire, analyser, puis dégager le sens. Ces étapes préparent directement la contraction de texte.}
\end{popfigure}

\courssub{Lecture méthodique n°4 : la remise de la médaille et le discours officiel}

\textbf{Situation du passage.} Le passage étudié se situe au cœur du roman : le 14 juillet, jour de la fête nationale française, le commandant de la circonscription remet solennellement une médaille à Meka, en présence des autorités coloniales et des villageois. C'est un passage à la fois \cle{argumentatif} et \cle{polémique} : derrière le discours officiel d'éloge se cache une violente critique de la colonisation.

\textbf{Lecture globale.} À la première lecture, on retient le sens manifeste : un vieil Africain est honoré par l'administration coloniale pour ses services rendus à la France. Le discours du commandant multiplie les éloges : Meka est présenté comme un ami fidèle, un modèle, un homme bon et généreux.

\textbf{Lecture analytique : repérage de la thèse implicite.} Mais Oyono ne dit jamais directement « la colonisation est hypocrite ». Sa \cle{thèse} est \emph{implicite} : elle se devine à travers les procédés d'écriture. Le lecteur attentif la reconstruit : \emph{la colonisation se pare de générosité alors qu'elle exploite et méprise ceux qu'elle prétend honorer}. C'est la dénonciation de l'\cle{hypocrisie coloniale}.

\begin{definition}[Thèse implicite]
La \cle{thèse implicite} est l'idée principale que l'auteur défend sans jamais l'énoncer clairement. Le lecteur doit la déduire des procédés (ironie, contraste, caricature) et de l'écart entre ce qui est dit et ce qui est montré. Elle relève du \cle{contenu latent} du texte.
\end{definition}

\textbf{Lecture analytique : les procédés.} Trois procédés dominent dans ce passage.

\begin{itemize}
\item \textbf{L'ironie} : l'auteur fait dire aux personnages le contraire de ce qu'il pense. Le commandant célèbre la « grandeur » de Meka, mais cet éloge exagéré sonne faux : il cache la volonté de ridiculiser le vieil homme et de faire de lui un instrument de propagande.
\item \textbf{La caricature} : les traits des personnages sont exagérés jusqu'au grotesque. Le commandant, suant dans son uniforme, débitant des phrases creuses ; les notables blancs s'arrachant Meka pour l'exhiber : chacun est réduit à un masque ridicule.
\item \textbf{Le contraste entre les discours et les réalités} : on célèbre la « fraternité » coloniale le jour même où Meka, ivre et perdu, sera abandonné sous la pluie, puis arrêté et humilié par les forces de l'ordre de ceux qui venaient de le décorer. L'écart entre les belles paroles et les actes est le ressort principal de la dénonciation.
\end{itemize}

\begin{exemplebox}[Contenu manifeste et contenu latent dans le discours du commandant]
Dans le discours officiel, le commandant déclare en substance : « Meka est un grand homme, un ami de la France, un exemple pour tous. »
\begin{itemize}
\item \textbf{Contenu manifeste} : Meka est un homme honorable que la France récompense.
\item \textbf{Contenu latent} : l'éloge est tellement exagéré et général qu'il devient vide de sens. Le commandant ne connaît même pas vraiment Meka. L'éloge cache la volonté de le ridiculiser et de l'utiliser comme preuve de la « bienveillance » coloniale. La médaille est un leurre : elle ne coûte rien à l'administration et ne rend ni les terres, ni les fils, ni la dignité de Meka.
\end{itemize}
\end{exemplebox}

\begin{attention}[Ne pas confondre le propos du personnage et celui de l'auteur]
Dans un texte ironique, ce que dit le personnage (le commandant) n'est pas ce que pense l'auteur (Oyono). Erreur fréquente à l'examen : écrire « Oyono fait l'éloge du colonisateur ». C'est l'inverse : Oyono met l'éloge dans la bouche du commandant pour mieux le détruire. Toujours se demander : \emph{que veut l'auteur en faisant parler ainsi son personnage ?}
\end{attention}

\courssub{Lecture méthodique n°5 : la prise de conscience de Meka}

\textbf{Situation du passage.} Après la cérémonie, Meka, fêté et enivré, s'égare en pleine nuit sous l'orage. Il est arrêté par les gendarmes, battu et enfermé. C'est là, dans l'humiliation, que s'opère sa \cle{prise de conscience} : le passage devient un texte de \emph{réflexion}, où le personnage repense toute sa vie.

\textbf{Lecture globale.} Le sens manifeste : Meka comprend que la médaille n'est rien, que les Blancs qu'il a servis toute sa vie le méprisent, et qu'il a perdu son identité, sa terre et sa famille pour rien.

\textbf{Lecture analytique.} Le passage repose sur :
\begin{itemize}
\item le \textbf{monologue intérieur} et le discours indirect libre, qui nous font entrer dans la pensée de Meka ;
\item le \textbf{contraste} entre la lumière et les honneurs de la veille et la nuit, la boue et la prison du lendemain ;
\item la \textbf{symbolique de la médaille} : objet brillant mais inutile, image de toutes les promesses coloniales ;
\item le \textbf{registre pathétique}, qui suscite la pitié du lecteur pour le vieil homme brisé.
\end{itemize}

\textbf{Lecture synthétique : dégagement des idées essentielles.} À cette étape, on formule chaque idée du passage en une phrase, avec ses propres mots.

\begin{exoresolu}[Dégager les idées essentielles du passage de la prise de conscience]
\textbf{Énoncé.} À partir du passage où Meka, enfermé après la fête, repense à sa vie, dégagez quatre idées essentielles, formulées chacune en une phrase.
\tcblower
\textbf{Solution détaillée.} On applique la troisième lecture (synthétique) : on repère les grands moments de la réflexion de Meka et on les reformule.
\begin{enumerate}
\item Meka prend conscience que la médaille reçue n'a aucune valeur réelle.
\item Il comprend que les colons qu'il a servis toute sa vie le méprisent et l'ont trompé.
\item Il réalise qu'il a sacrifié sa terre, ses fils et sa foi ancestrale pour rien.
\item Il éprouve douleur, honte et révolte devant le gâchis de son existence.
\end{enumerate}
Chaque idée est \emph{reformulée} : aucune n'est recopiée mot à mot du roman, mais toutes sont fidèles à sa pensée.
\end{exoresolu}

\begin{attention}[Idée essentielle n'est pas phrase copiée]
Une \cle{idée essentielle} n'est pas une phrase recopiée du texte : c'est une idée \emph{reformulée} avec ses propres mots. Copier une phrase du roman, même importante, ne prouve pas qu'on l'a comprise. La reformulation est la preuve de la compréhension. Au résumé, toute phrase copiée est sanctionnée.
\end{attention}

\courssub{De la lecture méthodique au résumé : le tableau des idées essentielles}

La lecture méthodique n'est pas une fin en soi : elle \emph{prépare la contraction de texte}. Le résumé consiste à conserver les idées essentielles, à les hiérarchiser, puis à les reformuler dans un texte bref et personnel. Le travail se visualise ainsi :

\begin{popfigure}
\begin{center}
\begin{tabularx}{\linewidth}{|>{\raggedright\arraybackslash}X|>{\raggedright\arraybackslash}X|>{\raggedright\arraybackslash}X|}
\hline
\rowcolor{popblueL} \textbf{Passage (lecture méthodique)} & \textbf{Idées essentielles dégagées} & \textbf{Reformulation pour le résumé} \\
\hline
Discours du commandant à la cérémonie & 1. L'administration honore Meka en paroles. \par 2. Cet éloge est hypocrite et intéressé. & Lors de la fête, l'administration couvre Meka d'éloges qui ne cachent qu'hypocrisie et calcul. \\
\hline
Nuit d'errance, arrestation et prison & 3. La réalité coloniale rattrape brutalement Meka. & Aussitôt après les honneurs, Meka est abandonné, puis arrêté et humilié par ceux-là mêmes qui l'ont décoré. \\
\hline
Méditation de Meka en prison & 4. Prise de conscience : sa vie entière a été un leurre. & Dans la prison, Meka comprend que la médaille est vide et qu'il a sacrifié sa vie pour rien. \\
\hline
\end{tabularx}
\end{center}
\legende{Tableau de progression du résumé : du passage analysé aux idées essentielles, puis à leur reformulation. Ce tableau, construit collectivement en classe, constitue le plan du futur résumé.}
\end{popfigure}

\textbf{Mise en commun.} En classe, chaque groupe propose ses idées essentielles ; le tableau collectif est construit au tableau en trois temps :
\begin{enumerate}
\item on \emph{regroupe} les idées identiques formulées différemment ;
\item on \emph{hiérarchise} : idées principales (la thèse, la prise de conscience) d'abord, idées secondaires (exemples, détails) ensuite — ces dernières pourront être supprimées au résumé ;
\item on \emph{reformule} chaque idée retenue en une phrase claire, personnelle et correctement liée aux autres.
\end{enumerate}

\begin{aretenir}[Ce qu'il faut retenir]
\begin{itemize}
\item La lecture méthodique comporte trois lectures : globale (sens général), analytique (procédés, contenus latents), synthétique (idées essentielles).
\item Dans \emph{Le Vieux Nègre et la médaille}, la thèse implicite d'Oyono est la dénonciation de l'hypocrisie coloniale, portée par l'ironie, la caricature et le contraste entre discours et réalités.
\item La prise de conscience de Meka transforme le récit en texte de réflexion : la médaille, symbole des promesses coloniales, se révèle vide.
\item Une idée essentielle est une idée reformulée en une phrase, jamais une phrase copiée.
\item La lecture méthodique prépare directement le résumé : dégager, hiérarchiser, reformuler.
\end{itemize}
\end{aretenir}

\begin{exercice}[À vous de jouer]
Relisez le passage de la remise de la médaille. 1) Relevez deux marques d'ironie dans le discours du commandant et expliquez leur contenu latent. 2) Dégagez trois idées essentielles du passage, formulées en une phrase chacune. 3) Classez ces idées de la plus importante à la moins importante et justifiez votre classement.
\end{exercice}

\begin{corrige}[Exercice : éléments de réponse]
1) Marques d'ironie possibles : les superlatifs exagérés (« grand homme », « ami fidèle de la France ») alors que le commandant ne connaît pas Meka ; l'insistance sur la « gratitude » de la France alors qu'elle a pris ses terres et ses fils. Contenu latent : la France coloniale se donne bonne conscience à peu de frais. 2) Idées essentielles : l'administration honore Meka publiquement ; cet honneur est une mise en scène hypocrite ; la médaille ne compense en rien les sacrifices de Meka. 3) L'idée principale est l'hypocrisie de la cérémonie (c'est la thèse implicite) ; les deux autres sont ses illustrations.
\end{corrige}

\coursec{Les techniques de réduction du texte}

La lecture méthodique du \emph{Vieux Nègre et la médaille} nous a permis, dans les séances précédentes, de repérer la thèse de Ferdinand Oyono, ses arguments et la progression de sa démonstration contre l'illusion coloniale. Nous savons désormais \emph{comprendre} un texte argumentatif. Il reste à franchir une étape décisive pour l'examen : savoir le \emph{réduire} sans le trahir. C'est l'objet de la \cle{contraction de texte}, exercice au programme du Probatoire et du Baccalauréat technique, qui repose sur trois techniques précises : la \cle{suppression}, la \cle{généralisation} et la \cle{condensation}.

\courssub{Le principe de la contraction de texte}

\courssub{Qu'est-ce que contracter un texte ?}

Imaginez que vous deviez raconter à un ami pressé un match de football de quatre-vingt-dix minutes en trois minutes. Vous ne pouvez pas tout dire : vous gardez le score, les buteurs, les moments décisifs, et vous laissez de côté les passes anodines, les remplacements sans conséquence, les descriptions des tribunes. Votre récit est plus court, mais il reste \emph{fidèle} : votre ami sait exactement ce qui s'est passé. La contraction de texte fonctionne de la même manière, avec des règles strictes.

\begin{definition}[La contraction de texte]
La \cle{contraction de texte} (ou résumé) est un exercice qui consiste à réduire un texte à une fraction donnée de sa longueur (le plus souvent le \textbf{quart}), en conservant intégralement son sens, sa thèse et sa progression logique, \textbf{sans aucun commentaire ni opinion personnelle}. Le résumé est rédigé avec les mots du candidat, dans un français correct, et il ne contient ni jugement (« l'auteur a raison »), ni appréciation (« ce beau texte »), ni ajout étranger au texte.
\end{definition}

\begin{attention}[Ce que la contraction n'est pas]
La contraction n'est ni un commentaire, ni une critique, ni une paraphrase libre. Trois fautes éliminatoires à l'examen :
\begin{itemize}
\item donner son avis (« je pense qu'Oyono exagère ») ;
\item recopier de longues phrases du texte sans les reformuler ;
\item dépasser ou rester fortement en dessous de la longueur imposée.
\end{itemize}
\end{attention}

\courssub{Le calcul préalable : combien de mots ?}

Avant d'écrire la moindre phrase, il faut \emph{compter}. Le sujet indique toujours le coefficient de réduction : « Résumez ce texte au quart de sa longueur », « en 150 mots environ », etc.

\begin{methode}[Préparer la contraction]
\begin{enumerate}
\item \textbf{Compter} les mots du texte (ou estimer : nombre de mots d'une ligne moyenne $\times$ nombre de lignes).
\item \textbf{Diviser} par le coefficient demandé : pour le quart, on divise par 4.
\item \textbf{Prévoir la marge} : une tolérance de $\pm 10\,\%$ est admise.
\item \textbf{Lire} le texte deux fois et repérer la thèse, les arguments, les articulations logiques.
\item \textbf{Réduire paragraphe par paragraphe}, en appliquant les trois techniques (suppression, généralisation, condensation).
\item \textbf{Compter} les mots du résumé rédigé et ajuster.
\end{enumerate}
\end{methode}

\begin{exemplebox}[Un calcul type d'examen]
Le texte proposé compte \textbf{600 mots}. Le résumé est demandé au quart :
\[ 600 \div 4 = 150 \text{ mots} \]
Avec la tolérance de $\pm 10\,\%$ : $150 \times 0{,}9 = 135$ et $150 \times 1{,}1 = 165$. Le résumé doit donc contenir \textbf{entre 135 et 165 mots}. Un résumé de 90 mots ou de 220 mots serait sanctionné, même s'il est bien écrit.
\end{exemplebox}

\begin{popfigure}
\begin{center}
\begin{tikzpicture}[
  etape/.style={draw=popdark, thick, rounded corners, fill=popblueL, text width=3.4cm, align=center, minimum height=1.1cm, font=\small},
  fleche/.style={-{Stealth[length=3mm]}, very thick, poporange},
  pct/.style={font=\small\bfseries, text=popdark}
]
\node[etape, fill=popgoldL, align=center] (t0) at (0,0){\textbf{Texte de départ}\\ 600 mots};
\node[pct] at (0,-1.05) {100 \%};
\node[etape, align=center] (t1) at (4.2,0){\textbf{Suppression}\\ répétitions, exemples, détails};
\node[etape, align=center] (t2) at (8.4,0){\textbf{Généralisation}\\ énumérations $\to$ termes génériques};
\node[etape, align=center] (t3) at (12.6,0){\textbf{Condensation}\\ fusion des phrases};
\node[etape, fill=popgreenL, align=center] (t4) at (12.6,-2.6){\textbf{Résumé final}\\ 150 mots};
\node[pct] at (10.9,-2.6) {25 \%};
\draw[fleche] (t0) -- (t1);
\draw[fleche] (t1) -- (t2);
\draw[fleche] (t2) -- (t3);
\draw[fleche] (t3) -- (t4);
\draw[popline, dashed, thick] (0,-0.6) -- (0,-2.2);
\node[font=\small\itshape, text=popmuted, rotate=90] at (-0.35,-1.4) {réduction progressive};
\end{tikzpicture}
\end{center}
\legende{Schéma en entonnoir de la contraction : le texte de départ (100 \%) traverse successivement les trois techniques de réduction pour aboutir au résumé (25 \%).}
\end{popfigure}

\courssub{Technique 1 : la suppression}

\courssub{Le principe : couper sans blesser}

La \cle{suppression} est la technique la plus intuitive : elle consiste à éliminer purement et simplement les éléments qui ne sont pas indispensables à la compréhension du raisonnement. Mais attention : tout ne se supprime pas. Il faut distinguer l'\emph{ossature} du texte (ce qui porte le sens) de sa \emph{parure} (ce qui l'embellit ou l'illustre).

\begin{propriete}[Ce qu'on supprime et ce qu'on conserve]
\textbf{On supprime obligatoirement :}
\begin{itemize}
\item les \textbf{répétitions} et les reformulations d'une même idée ;
\item les \textbf{exemples} et les anecdotes qui illustrent un argument déjà énoncé ;
\item les \textbf{citations} d'auteurs ;
\item les \textbf{comparaisons}, métaphores et figures de style ornementales ;
\item les \textbf{détails descriptifs} (couleurs, lieux, gestes) et les \textbf{précisions chiffrées secondaires} ;
\item les adjectifs et adverbes d'intensité non essentiels (« très », « vraiment », « magnifique »).
\end{itemize}
\textbf{On conserve obligatoirement :}
\begin{itemize}
\item la \textbf{thèse} de l'auteur ;
\item les \textbf{arguments} qui la soutiennent ;
\item la \textbf{progression logique} (les étapes du raisonnement et les connecteurs qui les marquent).
\end{itemize}
\end{propriete}

\begin{attention}[Le piège classique : l'exemple et l'argument]
Supprimer un exemple ne signifie \textbf{pas} supprimer l'argument qu'il illustre. Si le texte dit : « L'éducation transforme les hommes. Ainsi, Meka, qui n'a jamais fréquenté l'école, ne comprend pas la valeur de la médaille qu'on lui remet », on supprime l'exemple de Meka, mais on \emph{conserve} l'argument : « l'éducation transforme les hommes ». Le résumé sans argument serait amputé ; le résumé avec l'exemple serait trop long.
\end{attention}

\begin{exemplebox}[Suppression en action]
\textbf{Phrase originale (33 mots) :} « Le vieux Meka, ému jusqu'aux larmes, le cœur battant très fort comme un tam-tam de village, s'avança lentement, très lentement, vers l'estrade magnifique où l'attendait le commandant. »

\textbf{Éléments supprimés :} « ému jusqu'aux larmes » (détail), « le cœur battant très fort comme un tam-tam de village » (comparaison ornementale), « lentement, très lentement » (répétition), « magnifique » (adjectif d'appréciation).

\textbf{Forme réduite (11 mots) :} « Meka s'avança vers l'estrade où l'attendait le commandant. »
\end{exemplebox}

\courssub{Technique 2 : la généralisation}

\courssub{Le principe : du particulier au générique}

La \cle{généralisation} consiste à remplacer une énumération de termes particuliers par un seul \cle{terme générique} qui les englobe tous. C'est l'économie maximale : plusieurs mots sont remplacés par un ou deux, sans perte de sens essentiel.

\begin{exemplebox}[Généralisations typiques]
\begin{center}
\begin{tabularx}{\linewidth}{|X|X|}
\hline
\rowcolor{popblueL} \textbf{Énumération du texte} & \textbf{Terme générique} \\
\hline
manguiers, palmiers, bananiers, goyaviers & arbres fruitiers \\
\hline
pagnes, chemises, pantalons, robes & vêtements \\
\hline
le maire, le sous-préfet, le commandant, le curé & les autorités locales \\
\hline
rire, pleurer, s'indigner, s'enthousiasmer & réagir avec émotion \\
\hline
le lundi, le mardi, le mercredi, le jeudi & les jours de la semaine \\
\hline
\end{tabularx}
\end{center}
\end{exemplebox}

\begin{attention}[Condition d'application de la généralisation]
Le terme générique ne doit pas \emph{élargir} le sens au-delà de ce que dit le texte. Si le texte énumère « manguiers, palmiers, bananiers », écrire « la végétation » serait trop vague (cela inclurait les herbes, les lianes...), et écrire « les arbres » serait déjà acceptable, mais « les arbres fruitiers » est le plus fidèle. La généralisation doit rester \textbf{exacte} : ni trop étroite, ni trop large.
\end{attention}

\courssub{Technique 3 : la condensation}

\courssub{Le principe : fondre plusieurs phrases en une seule}

La \cle{condensation} est la technique la plus exigeante : elle consiste à \emph{fusionner} plusieurs phrases ou propositions en une seule phrase qui en exprime l'essentiel. On supprime les mots de liaison redondants, on remplace les subordonnées par des compléments, on ramène deux verbes à un seul.

\begin{exemplebox}[Condensation en action]
\textbf{Texte original (3 phrases, 30 mots) :} « Meka reçut une médaille. Cette médaille récompensait ses terres données à la mission. Elle récompensait aussi ses deux fils morts pour la France. »

\textbf{Forme condensée (1 phrase, 15 mots) :} « Meka reçut une médaille récompensant ses terres données à la mission et ses fils morts pour la France. »

Les deux dernières phrases, qui partageaient le même verbe « récompensait », ont été fusionnées grâce à la conjonction de coordination « et ».
\end{exemplebox}

\begin{popfigure}
\begin{center}
\begin{tikzpicture}[font=\small]
\node[text width=13.5cm, align=left] at (0,0) {%
\textbf{Étape 1 (original, 26 mots) :} Le vieux nègre, \textcolor{poporange}{[qui était très vieux et très fatigué]}, comprit \textcolor{poporange}{[enfin, mais beaucoup trop tard]}, que la médaille \textcolor{poporange}{[brillante et dorée]} n'était \textcolor{poporange}{[en réalité]} qu'un leurre.};
\node[text width=13.5cm, align=left] at (0,-2.1) {%
\textbf{Étape 2 (suppression des détails) :} Le vieux nègre comprit \textcolor{poporange}{[trop tard]} que la médaille n'était qu'un leurre.};
\node[text width=13.5cm, align=left] at (0,-3.7) {%
\textbf{Étape 3 (forme réduite, 13 mots) :} \textbf{Meka comprit que la médaille n'était qu'un leurre.}};
\draw[-{Stealth}, very thick, popblue] (-6.4,-0.9) -- (-6.4,-1.5);
\draw[-{Stealth}, very thick, popblue] (-6.4,-2.9) -- (-6.4,-3.3);
\end{tikzpicture}
\end{center}
\legende{Réduction progressive d'une phrase : les éléments entre crochets et en couleur (détails, répétitions, ornements) sont supprimés étape par étape jusqu'à la forme essentielle.}
\end{popfigure}

\courssub{Entraînement guidé : réduire un paragraphe de 120 mots à 30 mots}

\begin{exoresolu}[Réduction guidée d'un paragraphe argumentatif]
\textbf{Énoncé.} Résumez le paragraphe suivant (120 mots) en 30 mots environ (le quart), en appliquant les trois techniques de réduction.

« L'école est la clé du développement de nos villages. En effet, un enfant instruit devient un adulte capable de lire un contrat, de calculer ses récoltes, de soigner sa famille en suivant les prescriptions d'un médecin. Prenons l'exemple de Ngoumou, petite localité de la région du Centre : depuis l'ouverture d'un collège en 1998, le nombre de commerçants sachant tenir une comptabilité a triplé, passant de 12 à 36 boutiques bien gérées. On le voit donc clairement, mes chers amis, l'instruction transforme la vie quotidienne, elle change tout, absolument tout, comme le soleil qui fait jaillir le maïs de la terre rouge. »

\tcblower
\textbf{Solution détaillée, étape par étape.}

\textbf{Étape 1 — Repérage de l'ossature.} Thèse : « l'école est la clé du développement des villages ». Argument : « l'instruit sait lire, calculer, se soigner ». Conclusion : « l'instruction transforme la vie quotidienne ».

\textbf{Étape 2 — Suppression.} On élimine : l'exemple chiffré de Ngoumou (« 1998 », « 12 à 36 boutiques »), l'apostrophe « mes chers amis », la répétition « elle change tout, absolument tout », la comparaison « comme le soleil qui fait jaillir le maïs de la terre rouge ».

\textbf{Étape 3 — Généralisation.} « lire un contrat, calculer ses récoltes, soigner sa famille » devient « gérer sa vie quotidienne » ou « s'occuper utilement de ses affaires et de sa santé ».

\textbf{Étape 4 — Condensation.} On fusionne thèse, argument et conclusion en deux phrases liées par un connecteur logique.

\textbf{Résumé proposé (30 mots) :} « L'école est indispensable au développement des villages, car l'instruction rend l'adulte capable de gérer ses affaires et sa santé : elle transforme toute la vie quotidienne. »

\textbf{Vérification :} $120 \div 4 = 30$ mots ; notre résumé en compte 30 ; la thèse, l'argument et la progression (thèse $\to$ justification $\to$ conclusion) sont conservés ; aucune opinion personnelle n'a été ajoutée.
\end{exoresolu}

\begin{exercice}[À vous de contracter]
Résumez en 25 mots environ le passage suivant (100 mots) : « Le sport est nécessaire à la santé des jeunes. D'abord, il fortifie le cœur, les muscles, les os et les poumons. Ensuite, il apprend la discipline, le respect des règles, le goût de l'effort. Ainsi, à Douala, le jeune Etame, qui s'entraînait chaque matin à cinq heures avant les cours, est devenu champion régional tout en réussissant brillamment son baccalauréat technique avec une moyenne de 15,4. Le sport, véritablement, forme l'homme tout entier, comme le feu de la forge façonne le fer. »
\end{exercice}

\begin{corrige}[Exercice : contraction du passage sur le sport]
\textbf{Corrigé type (25 mots) :} « Le sport est indispensable aux jeunes : il renforce leur organisme et leur enseigne la discipline et l'effort. Il forme ainsi l'homme tout entier. »

\textbf{Commentaire.} On a supprimé l'exemple d'Etame et ses précisions chiffrées (« cinq heures », « 15,4 »), la comparaison finale (« comme le feu de la forge ») ; on a généralisé « le cœur, les muscles, les os et les poumons » en « leur organisme » ; on a conservé la thèse (« le sport est nécessaire »), les deux arguments (santé, discipline) et la progression marquée par les deux points et le connecteur « ainsi ».
\end{corrige}

\begin{aretenir}[L'essentiel de la section]
\begin{itemize}
\item La contraction réduit un texte à une fraction imposée (souvent le quart), \textbf{sans trahir le sens ni ajouter d'opinion}.
\item On calcule d'abord la longueur cible : nombre de mots $\div$ coefficient, avec une tolérance de $\pm 10\,\%$.
\item Trois techniques complémentaires : la \cle{suppression} (répétitions, exemples, citations, figures ornementales), la \cle{généralisation} (énumération $\to$ terme générique exact), la \cle{condensation} (fusion de phrases).
\item On conserve toujours : la \textbf{thèse}, les \textbf{arguments}, la \textbf{progression logique}.
\item Supprimer un exemple ne supprime pas l'argument qu'il illustre.
\end{itemize}
\end{aretenir}

Ces techniques de réduction supposent une maîtrise des outils de liaison de la phrase : c'est en manipulant les conjonctions de coordination et de subordination que l'on pourra condenser et reformuler avec correction. C'est l'objet de la section suivante.

\coursec{Les outils de liaison : conjonctions de coordination et de subordination}

Dans la section précédente, nous avons appris à \cle{réduire} un texte : supprimer les redondances, les exemples superflus, les répétitions. Mais réduire ne suffit pas. Pour obtenir un résumé fluide et dense, il faut encore \emph{resserrer} les phrases les unes aux autres, c'est-à-dire les \cle{fusionner}. Or cette fusion obéit à des règles grammaticales précises : celles de la \cle{coordination} et de la \cle{subordination}. C'est l'objet de cette section, qui constitue l'outil technique majeur de la contraction de texte.

\courssub{Intuition : deux façons d'attacher les phrases}

Observez ces deux formulations issues de notre lecture du \emph{Vieux Nègre et la Médaille} :

\begin{enumerate}
\item \emph{Meka est vieux. Meka reçoit une médaille.}
\item \emph{Meka, qui est vieux, reçoit une médaille.}
\end{enumerate}

Dans la première version, les deux phrases sont posées côte à côte, comme deux wagons non attelés. Dans la seconde, la première idée a été \emph{intégrée} dans la seconde grâce au mot \emph{qui}. Le sens est le même, mais la seconde version est plus courte et plus élégante. Toute la grammaire de cette section repose sur cette intuition : on peut \textbf{juxtaposer} des éléments de même rang (coordination) ou \textbf{enchâsser} un élément dans un autre (subordination).

\courssub{La coordination : juxtaposer des éléments de même fonction}

\begin{definition}[La coordination]
La \cle{coordination} est un procédé grammatical qui unit deux éléments de \textbf{même nature} et de \textbf{même fonction} (deux mots, deux groupes de mots ou deux propositions) au moyen d'une \cle{conjonction de coordination}. Les deux éléments restent indépendants : chacun pourrait, grammaticalement, se tenir sans l'autre.
\end{definition}

\begin{aretenir}[Les sept conjonctions de coordination]
Il n'existe que \textbf{sept} conjonctions de coordination en français. Il faut les connaître par cœur, dans l'ordre traditionnel :
\begin{center}
\textbf{mais -- ou -- et -- donc -- or -- ni -- car}
\end{center}
Moyen mnémotechnique : \emph{« Mais où est donc Ornicar ? »}
\end{aretenir}

\begin{exemplebox}[Coordination dans \emph{Le Vieux Nègre et la Médaille}]
\emph{« Meka est fier, \textbf{car} on lui a décerné une médaille. »}

Ici, \emph{car} relie deux propositions indépendantes et introduit une \textbf{cause}. Chacune des deux propositions pourrait exister seule : \emph{Meka est fier.} / \emph{On lui a décerné une médaille.} La coordination exprime simplement le lien logique entre elles.
\end{exemplebox}

\begin{center}
\begin{tabularx}{\linewidth}{|l|X|X|}
\hline
\rowcolor{popblue!40} \textbf{Conjonction} & \textbf{Valeur logique} & \textbf{Exemple} \\
\hline
et & addition & Il lut la lettre \emph{et} sourit. \\
\hline
ou & alternative & Viens \emph{ou} reste. \\
\hline
ni & négation additive & Il ne mange \emph{ni} ne boit. \\
\hline
mais & opposition & Il est pauvre \emph{mais} digne. \\
\hline
donc & conséquence & Il pleut, \emph{donc} je reste. \\
\hline
or & transition, argument & Il se tait, \emph{or} il devrait parler. \\
\hline
car & cause & Il rit, \emph{car} il est heureux. \\
\hline
\end{tabularx}
\end{center}

\courssub{La subordination : intégrer une proposition dans une autre}

\begin{definition}[La proposition subordonnée]
Une \cle{proposition subordonnée} est une proposition qui \textbf{dépend} d'une \cle{proposition principale} et qui \textbf{ne peut pas se tenir seule}. Elle est introduite :
\begin{itemize}
\item soit par une \textbf{conjonction de subordination} (\emph{que, parce que, bien que, si, lorsque, puisque...}) $\rightarrow$ proposition subordonnée conjonctive (complétive ou circonstancielle) ;
\item soit par un \textbf{pronom relatif} (\emph{qui, que, où, dont...}) $\rightarrow$ proposition subordonnée relative (fonction d'épithète).
\end{itemize}
Elle occupe une fonction dans la phrase (complément d'objet, épithète, complément circonstanciel...).
\end{definition}

Comparez :
\begin{itemize}
\item \emph{Meka est fier.} $\rightarrow$ phrase complète, autonome.
\item \emph{Parce qu'on lui a décerné une médaille.} $\rightarrow$ impossible de laisser cette proposition seule : elle attend sa principale : \emph{Meka est fier \textbf{parce qu'}on lui a décerné une médaille.}
\end{itemize}

La subordination crée donc une \textbf{hiérarchie} : la proposition principale porte le message essentiel, la subordonnée apporte une précision (cause, temps, condition, description...). C'est exactement le mouvement du résumé : distinguer l'essentiel de l'accessoire, puis organiser l'accessoire autour de l'essentiel.

\begin{popfigure}
\begin{tikzpicture}[
  grow=down,
  level distance=1.35cm,
  level 1/.style={sibling distance=5.6cm},
  level 2/.style={sibling distance=4.4cm},
  every node/.style={draw, rounded corners=2pt, align=center, font=\small, inner sep=3pt},
  edge from parent/.style={draw=popmuted, -{Stealth[length=2mm]}}
]
\node[fill=popblue!35, align=center]{\textbf{Phrase complexe}\\ \emph{Meka, bien qu'il soit vieux, se rendit}\\\emph{ à la cérémonie parce qu'il était fier.}}
  child { node[fill=popgreen!35, align=center]{\textbf{Proposition principale}\\ \emph{Meka se rendit}\\\emph{ à la cérémonie}} }
  child { node[fill=poporange!35] {\textbf{Subordonnées}}
    child { node[fill=poppurple!30, align=center]{concession (circonst.)\\ \emph{bien qu'il soit vieux}\\ (subjonctif)} }
    child { node[fill=poppink!30, align=center]{cause (circonst.)\\ \emph{parce qu'il était fier}\\ (indicatif)} }
  };
\end{tikzpicture}
\legende{Arbre de la phrase complexe : la proposition principale (vert) porte le message essentiel ; les subordonnées (violet, rose) s'y rattachent comme des branches et expriment des circonstances (concession, cause).}
\end{popfigure}

\courssub{Les valeurs des conjonctions de subordination et des relatifs}

Chaque outil de subordination exprime un \textbf{rapport logique} précis. Bien le choisir, c'est préserver le sens du texte de départ : une erreur de conjonction ou de pronom relatif dans un résumé est une erreur de sens.

\begin{center}
\begin{tabularx}{\linewidth}{|l|X|X|}
\hline
\rowcolor{popblue!40} \textbf{Valeur logique} & \textbf{Outils principaux (conjonctions / relatifs)} & \textbf{Exemple} \\
\hline
Cause & parce que, puisque, comme, vu que & \emph{Comme} il pleuvait, le match fut reporté. \\
\hline
Conséquence & si bien que, de sorte que, au point que & Il plut, \emph{si bien que} le match fut reporté. \\
\hline
Concession & bien que, quoique, même si, alors que & \emph{Bien qu'}il soit vieux, il marcha longtemps. \\
\hline
Condition & si, à condition que, pourvu que & \emph{Si} tu résumes bien, tu réussiras. \\
\hline
Temps & quand, lorsque, dès que, pendant que, avant que & \emph{Lorsque} la lettre arriva, Meka sourit. \\
\hline
But & pour que, afin que, de peur que & Il se leva tôt \emph{afin qu'}on le remarque. \\
\hline
Comparaison & comme, ainsi que, plus... que & Il marcha \emph{comme} un jeune homme. \\
\hline
\multicolumn{3}{|c|}{\textbf{Outils de subordination relative (description, identification)}} \\
\hline
Épithète / Identification & qui (sujet), que (COD), dont (objet prép.), où (lieu/temps) & \emph{Meka, \textbf{qui} est vieux, reçoit une médaille.} \\
\hline
\end{tabularx}
\end{center}

\begin{attention}[Pièges de mode et de nature -- Vigilance Bac]
\begin{itemize}
\item \emph{\textbf{Bien que}} et \emph{\textbf{quoique}} sont suivis du \textbf{subjonctif}, jamais de l'indicatif : on écrit \emph{bien qu'il \textbf{soit} vieux}, et non \emph{bien qu'il est vieux}. De même \emph{avant que} et \emph{pour que} (but) exigent le subjonctif.
\item \emph{\textbf{Car}} est une conjonction de \textbf{coordination} : elle relie deux propositions indépendantes (pas de virgule obligatoire devant, mais possible). \emph{\textbf{Parce que}}, en revanche, introduit une \textbf{subordonnée} (virgule souvent omise si la subordonnée est en fin de phrase). Ne confondez jamais les deux dans l'analyse grammaticale.
\item \emph{Si} exprimant la condition est suivi de l'\textbf{indicatif} (présent ou imparfait selon le temps de la principale), jamais du conditionnel : \emph{si tu viens}, \emph{si tu venais} ; non \emph{si tu viendrais}.
\item \emph{\textbf{Comme}} peut marquer la \textbf{cause} (subordination, indicatif : \emph{Comme il est tard, je pars}) ou la \textbf{comparaison} (subordination ou coordination selon structure). Ne pas confondre.
\end{itemize}
\end{attention}

\courssub{Le rôle des outils de liaison dans le résumé}

Pourquoi cette grammaire est-elle au cœur de la contraction de texte ? Parce que la subordination (conjonctive ou relative) permet de \textbf{fusionner} des phrases et donc de \textbf{gagner des mots} sans perdre d'idées. Comptons sur un exemple simple :

\begin{itemize}
\item Version juxtaposée (coordination ou phrases séparées) : \emph{Il pleuvait. Le match fut reporté.} $\rightarrow$ 7 mots, deux phrases.
\item Version subordonnée (cause) : \emph{Comme il pleuvait, le match fut reporté.} $\rightarrow$ 7 mots, une seule phrase, le lien logique devenu explicite.
\item Version resserrée (nominalisation + subordination relative) : \emph{La pluie fit reporter le match.} $\rightarrow$ 6 mots.
\end{itemize}

Le gain paraît modeste sur une phrase ; sur un texte de vingt lignes à ramener au quart, la fusion systématique des phrases par subordination fait gagner des lignes entières tout en améliorant le style et la cohérence.

\begin{methode}[Transformer des phrases juxtaposées en phrase complexe]
\begin{enumerate}
\item \textbf{Identifier le rapport logique} entre les deux phrases : cause, conséquence, opposition, temps, condition, description ?
\item \textbf{Choisir l'idée principale} : c'est celle qui deviendra la proposition principale (souvent le résultat, le fait central, l'action principale).
\item \textbf{Choisir l'outil} adapté : conjonction de subordination (circonstancielle) \textbf{ou} pronom relatif (si l'idée secondaire décrit un nom de la principale).
\item \textbf{Construire la subordonnée} autour de l'idée secondaire, en veillant au mode (subjonctif après \emph{bien que}, \emph{avant que}, \emph{pour que}, \emph{quoique}).
\item \textbf{Vérifier} que la phrase obtenue est correcte, plus courte, et fidèle au sens initial (pas de contresens logique).
\end{enumerate}
\end{methode}

\begin{exoresolu}[Fusion de trois phrases en une phrase complexe]
\textbf{Énoncé.} Fusionnez les trois phrases suivantes en une seule phrase complexe correcte, sans changer le sens : \emph{Meka était très âgé. Il se rendit à la cérémonie. On lui remettait une médaille.}
\tcblower
\textbf{Solution détaillée.}
\begin{enumerate}
\item \textbf{Analyse des rapports logiques :}
\begin{itemize}
\item Entre « Meka était très âgé » et « Il se rendit à la cérémonie » : \textbf{opposition / concession} (malgré son âge, il y alla).
\item Entre « On lui remettait une médaille » et « Il se rendit à la cérémonie » : \textbf{cause} (il y alla \emph{parce que} on le décorait).
\end{itemize}
\item \textbf{Choix de la principale :} \emph{Il se rendit à la cérémonie} (fait central, action principale).
\item \textbf{Choix des outils :}
\begin{itemize}
\item Concession $\rightarrow$ \emph{bien que} + \textbf{subjonctif} (imparfait du subjonctif car la principale est au passé simple).
\item Cause $\rightarrow$ \emph{parce que} + indicatif (imparfait).
\end{itemize}
\item \textbf{Phrase obtenue :} \emph{\textbf{Bien qu'il fût très âgé, Meka se rendit à la cérémonie parce qu'on lui remettait une médaille.}}
\end{enumerate}
\textbf{Vérification :} une seule phrase au lieu de trois ; le subjonctif \emph{fût} est bien employé après \emph{bien que} ; le sens est intégralement conservé ; la phrase est fluide.
\end{exoresolu}

\begin{exercice}[À vous de fusionner]
Fusionnez chaque groupe de phrases en une seule phrase complexe, en choisissant l'outil de liaison approprié (conjonction de subordination ou pronom relatif) :
\begin{enumerate}
\item \emph{La pluie tombait sans cesse. Le village resta isolé.}
\item \emph{Meka avait perdu sa médaille. Il garda sa dignité.}
\item \emph{Tu auras identifié les idées essentielles. Tu pourras rédiger ton résumé.}
\item \emph{Le chef de village parla. Sa voix était tremblante.}
\end{enumerate}
\end{exercice}

\begin{corrige}[Exercice : fusion de phrases]
\begin{enumerate}
\item \textbf{Rapport de cause} (la pluie cause l'isolement) : \emph{Parce que la pluie tombait sans cesse, le village resta isolé.} \\
\textit{Variante conséquence :} \emph{La pluie tombait sans cesse, si bien que le village resta isolé.}
\item \textbf{Rapport de concession} (malgré la perte, dignité gardée) : \emph{Bien qu'il eût perdu sa médaille, Meka garda sa dignité.} (Subjonctif obligatoire après \emph{bien que} ; plus-que-parfait du subjonctif pour l'antériorité sur le passé simple \emph{garda}).
\item \textbf{Rapport de temps} (antériorité) ou \textbf{condition} : \emph{Lorsque tu auras identifié les idées essentielles, tu pourras rédiger ton résumé.} (Futur antérieur après \emph{lorsque} pour action future antérieure). \\
\textit{Variante condition :} \emph{Si tu identifies les idées essentielles, tu pourras rédiger ton résumé.}
\item \textbf{Rapport de description} (attribut du sujet « chef ») $\rightarrow$ \textbf{Subordination relative} : \emph{Le chef de village, \textbf{dont} la voix était tremblante, parla.} Ou : \emph{Le chef de village, \textbf{la voix} tremblante, parla.} (Participial, encore plus concis).
\end{enumerate}
\end{corrige}

\begin{savaistu}[Le Vieux Nègre et la Médaille : la langue de l'humiliation et de la dignité]
Dans l'œuvre de Ferdinand Oyono, le style alterne phrases courtes, coordonnées (effet de hachure, de solitude, de fait brut : \emph{« Meka est vieux. Il attend. »}) et phrases complexes subordonnées (quand la pensée, le souvenir ou la dignité reprennent le dessus). Maîtriser la subordination, c'est donc aussi comprendre comment l'auteur construit le personnage : la \textbf{coordination} isole, la \textbf{subordination} relie et donne de la profondeur. Votre résumé doit refléter cette tension.
\end{savaistu}

\begin{aretenir}[L'essentiel de la section]
\begin{itemize}
\item La \cle{coordination} juxtapose des éléments de même fonction grâce à sept conjonctions : \emph{mais, ou, et, donc, or, ni, car}. Elle place les idées sur un pied d'égalité.
\item La \cle{subordination} intègre une proposition dépendante dans une proposition principale. Elle s'opère par \textbf{conjonctions de subordination} (circonstancielles : cause, temps, concession...) ou \textbf{pronoms relatifs} (relatives : description, identification). Chaque outil exprime une valeur logique précise.
\item Dans le résumé, la subordination est l'outil roi : elle fusionne les phrases, resserre le texte et rend explicites les liens logiques (cause, concession, temps, description).
\item \textbf{Vigilance Bac :} \emph{car} coordonne (deux principales), \emph{parce que} subordonne (principale + subordonnée) ; \emph{bien que}, \emph{quoique}, \emph{avant que}, \emph{pour que} exigent le \textbf{subjonctif} ; \emph{si} (condition) + indicatif.
\end{itemize}
\end{aretenir}

Ces outils de liaison maîtrisés, nous pourrons passer à l'étape suivante : la \emph{reformulation} proprement dite des idées essentielles, c'est-à-dire leur réécriture dans un style personnel et condensé.

\coursec{La reformulation des idées essentielles}

Après avoir maîtrisé les \cle{outils de liaison} — conjonctions de coordination et de subordination — qui assurent la cohésion syntaxique à l'intérieur de la phrase, nous abordons maintenant l'étape décisive du résumé : la \cle{reformulation}. Posséder des connecteurs ne suffit pas ; il faut savoir \emph{dire autrement} ce que le texte source exprime, sans le trahir, sans l'alourdir, et en respectant la contrainte de longueur imposée par l'épreuve. C'est le passage de la \emph{réduction} (supprimer le superflu) à la \emph{réécriture} (produire un texte autonome et fidèle).

\courssub{Définition et enjeux de la reformulation}

\begin{definition}[La reformulation]
La \textbf{reformulation} est l'opération qui consiste à \textbf{restituer le sens d'un énoncé source en utilisant ses propres mots et ses propres structures syntaxiques}, tout en conservant intégralement l'information, la portée argumentative et la logique interne du texte original. Elle s'oppose au \emph{recopiage} (copier-coller de segments) et à la \emph{paraphrase libre} (qui ajoute, commente ou déforme).
\end{definition}

Au Probatoire comme au Baccalauréat technique, la reformulation est le \textbf{critère discriminant} entre une copie moyenne et une copie d'excellence. Le correcteur vérifie trois exigences indissociables :
\begin{enumerate}
    \item \textbf{Fidélité sémantique :} aucune idée force n'est omise, aucune nuance n'est altérée, aucun jugement personnel n'est introduit.
    \item \textbf{Autonomie langagière :} le vocabulaire et la syntaxe sont renouvelés (synonymes, conversion actif/passif, nominalisation/verbalisation, regroupement d'idées).
    \item \textbf{Économie verbale :} la reformulation est plus concise que l'original ; elle condense l'information sans la vider de sa substance.
\end{enumerate}

\courssub{Les techniques de reformulation}

La reformulation efficace mobilise un « outillage » linguistique précis. Voici les quatre leviers principaux, illustrés sur un extrait type.

\begin{exemplebox}[Extrait de référence (inspiré de \emph{Le Vieux Nègre et la Médaille})]
\textit{« Meka, accablé par le poids des ans et la honte cuisante de son inutilité, comprit avec une lucidité douloureuse que la médaille brillait sur sa poitrine comme un miroir aux alouettes, un leurre grossier tendu par l'administration coloniale pour masquer l'exploitation systématique de son peuple. »}
\end{exemplebox}

\begin{center}
\begin{tabularx}{\linewidth}{|l|X|X|}
\hline
\rowcolor{popblueL}
\textbf{Technique} & \textbf{Application correcte} & \textbf{Erreur fréquente (recopiage / déformation)} \\
\hline
\textbf{1. Substitution lexicale (synonymes, hyperonymes)} & Meka, \textbf{vieux et humilié}, \textbf{prit conscience} que la médaille \textbf{n'était qu'une illusion}, un \textbf{piège} colonial pour \textbf{cacher l'oppression}. & Meka, \textbf{accablé par le poids des ans et la honte}, \textbf{comprit} que la médaille \textbf{brillait comme un miroir aux alouettes}... (segments de 6 à 10 mots recopiés) \\
\hline
\textbf{2. Conversion diathétique (actif $\leftrightarrow$ passif / impersonnel)} & \textbf{La lucidité frappa Meka :} la médaille \textbf{apparaissait comme} un leurre colonial. & Meka comprit que la médaille \textbf{était tendue par l'administration}... (conservation de la structure agent-patient sans gain de concision) \\
\hline
\textbf{3. Nominalisation / Verbalisation (transcatégoriel)} & \textbf{L'humiliation de Meka} s'accompagna de \textbf{la prise de conscience du caractère illusoire} de sa récompense. & Meka \textbf{fut humilié et comprit}... (suite de verbes faibles, pas de condensation nominale) \\
\hline
\textbf{4. Regroupement / Fusion d'idées (condensation)} & \textbf{Humilié, Meka démasqua la médaille :} un leurre colonial masquant l'exploitation de son peuple. & Meka \textbf{comprit que la médaille était un leurre.} \textbf{Elle masquait l'exploitation.} (deux phrases au lieu d'une, perte de fluidité) \\
\hline
\end{tabularx}
\end{center}

\begin{popfigure}
\begin{tikzpicture}[
    node distance=1.2cm and 1.5cm,
    box/.style={draw=popblue, thick, rounded corners, fill=popblue!5, minimum width=3.5cm, minimum height=1.2cm, align=center, font=\small},
    arrow/.style={-Stealth, thick, popdark},
    labelbox/.style={font=\footnotesize\itshape, text=popmuted}
]
\node[box, align=center] (orig){Phrase originale\\ (segment repéré)};
\node[box, right=of orig, align=center] (tech){Choix de la technique\\ (synonymie, voix,\\ nominalisation, regroupement)};
\node[box, right=of tech, align=center] (reform){Ébauche reformulée\\ (mots propres,\\ structure nouvelle)};
\node[box, below=of tech, align=center] (check){Vérification :\\ Fidélité ? Concision ?\\ Fluidité ?};
\node[box, right=of reform, align=center] (final){Phrase finale\\ intégrée au résumé};
\draw[arrow] (orig) -- node[above, labelbox] {1. Isoler} (tech);
\draw[arrow] (tech) -- node[above, labelbox] {2. Transformer} (reform);
\draw[arrow] (reform.south) |- node[above, labelbox, pos=0.75] {3. Contrôler} (check.east);
\draw[arrow] (check.north) -- node[right, labelbox] {Si défaut : retour} (tech.south);
\draw[arrow] (reform) -- node[above, labelbox] {4. Insérer} (final);
\end{tikzpicture}
\legende{Schéma du processus mental de reformulation d'un segment repéré. La boucle de vérification (étape 3) est cruciale : on ne valide la reformulation que si elle est fidèle, plus courte et bien liée aux phrases voisines ; sinon, on retourne au choix de la technique.}
\end{popfigure}

\courssub{La fidélité au texte source : règles d'or}

\begin{attention}[Sanction du recopiage au Bac]
\textbf{Au Probatoire et au Baccalauréat technique, la copie de segments du texte original (plus de 4 à 5 mots consécutifs) est considérée comme un manquement grave à la consigne de « reformuler ».} Elle entraîne une pénalité forte sur la note de \emph{reformulation} et peut faire basculer la note globale en dessous de la moyenne. \textbf{Interdiction absolue :} recopier une définition, une citation, une image forte, une énumération. \textbf{Obligation :} tout passer au crible de ses propres mots.
\end{attention}

\begin{savaistu}[Le compte-mots : exigence formelle]
L'indication du \textbf{nombre exact de mots} à la fin du résumé est une \textbf{exigence formelle de l'épreuve}. Un écart supérieur à $\pm 10\%$ par rapport à la longueur prescrite (souvent 120 mots $\pm 10\%$, soit 108 à 132 mots) coûte des points sur la note de \emph{concision} et de \emph{respect de la consigne}. Comptez \emph{après} la rédaction définitive (mots-outils inclus, nombres en chiffres comptent pour un mot). Notez : \textbf{Nombre de mots : 118}.
\end{savaistu}

La fidélité impose aussi de \textbf{ne rien ajouter} : ni exemple personnel, ni connaissance extérieure, ni commentaire (« c'est triste », « cela montre que... »), ni développement implicite. Le résumé n'est pas une dissertation. Il est le \emph{miroir réduit} du texte.

\courssub{La cohérence du résumé : articulation des idées}

Une suite de phrases reformulées isolément ne fait pas un résumé. Le \cle{texte résumé} doit être lisible comme un texte à part entière. Cela suppose :
\begin{itemize}
    \item \textbf{Des connecteurs logiques adaptés :} \emph{donc, par conséquent, en effet, cependant, de plus, enfin} (repris de la section précédente sur les conjonctions).
    \item \textbf{Une progression thématique :} on suit l'ordre des idées forces du texte (souvent : situation initiale $\rightarrow$ élément déclencheur $\rightarrow$ prise de conscience $\rightarrow$ conséquence/ouverture).
    \item \textbf{L'absence de répétitions lexicales inutiles :} utiliser des pronoms, des reprises nominales variées (\emph{le protagoniste, ce dernier, l'ancien combattant, Meka}).
    \item \textbf{Le temps des verbes :} généralement le \textbf{présent de narration} (ou passé simple si le texte est au passé et que le résumé adopte une valeur narrative), mais \emph{un seul temps principal} maintenu tout au long.
\end{itemize}

\courssub{Méthode complète en 5 étapes : de la lecture à la copie finale}

\begin{methode}[La méthode du résumé en 5 étapes (version Bac)]
\begin{enumerate}
    \item \textbf{Lecture analytique (2 lectures) :} 1$^{\text{re}}$ lecture : compréhension globale (sujet, thèse, ton). 2$^{\text{e}}$ lecture : repérage des \textbf{idées forces} (une par paragraphe ou par séquence logique) $\rightarrow$ surligner / numéroter dans la marge.
    \item \textbf{Réduction (brouillon) :} Rayer le superflu (exemples, détails, répétitions, adjectifs qualificatifs non essentiels, connecteurs de transition internes). Ne garder que le \textbf{squelette argumentatif} : sujets, verbes, compléments indispensables, mots-clés.
    \item \textbf{Reformulation segment par segment (brouillon) :} Pour chaque idée force retenue, \textbf{réécrire en français propre} (synonymes, voix passive/active, nominalisation, regroupement). \textbf{Ne pas regarder le texte original} pendant cette phase (le cacher si nécessaire). Vérifier : fidélité ? concision ? pas de recopiage $> 4$ mots ?
    \item \textbf{Rédaction suivie (brouillon) :} Enchaîner les phrases reformulées avec des \textbf{connecteurs logiques} adaptés. Unifier le temps verbal. Vérifier la fluidité (lecture à voix basse).
    \item \textbf{Vérification finale \& Copie propre :} \textbf{Compter les mots} (corps du texte seulement). Ajuster (couper ou développer légèrement) pour entrer dans la fourchette $\pm 10\%$. Reporter le total : \textbf{Nombre de mots : XXX}. Recopier soigneusement sur la copie d'examen (lisibilité, pas de ratures).
\end{enumerate}
\end{methode}

\begin{popfigure}
\begin{tikzpicture}[
    node distance=1.0cm and 0.8cm,
    stepbox/.style={draw=poporange, thick, rounded corners, fill=poporange!10, minimum width=2.8cm, minimum height=1.8cm, align=center, font=\small\bfseries},
    arrow/.style={-Stealth, very thick, poporange}
]
\node[stepbox, align=center] (e1){\textbf{1. LIRE}\\\emph{Analytique}\\\textcolor{popmuted}{$\times$ 2}};
\node[stepbox, right=of e1, align=center] (e2){\textbf{2. RÉDUIRE}\\\emph{Squelette}\\\textcolor{popmuted}{rayer le superflu}};
\node[stepbox, right=of e2, align=center] (e3){\textbf{3. REFORMULER}\\\emph{Mots propres}\\\textcolor{popmuted}{synonymes, voix, nom.}};
\node[stepbox, right=of e3, align=center] (e4){\textbf{4. RÉDIGER}\\\emph{Enchaîner}\\\textcolor{popmuted}{connecteurs, temps unique}};
\node[stepbox, right=of e4, align=center] (e5){\textbf{5. VÉRIFIER}\\\emph{Compter \& Copier}\\\textcolor{popmuted}{$\pm 10\%$, propreté}};
\draw[arrow] (e1) -- (e2);
\draw[arrow] (e2) -- (e3);
\draw[arrow] (e3) -- (e4);
\draw[arrow] (e4) -- (e5);
\draw[arrow, popgreen, bend left=30] (e5.north) to node[above, font=\tiny, text=popgreen] {Si hors fourchette} (e4.north);
\draw[arrow, poppink, bend right=30] (e4.south) to node[below, font=\tiny, text=poppink] {Si infidèle / recopié} (e3.south);
\end{tikzpicture}
\legende{Diagramme de la méthode complète en 5 étapes. Les boucles de rétroaction (vertes et roses) rappellent que le processus est itératif : on revient en arrière si le contrôle final révèle un défaut (longueur, fidélité, recopiage).}
\end{popfigure}

\courssub{Exemple d'application sur \emph{Le Vieux Nègre et la Médaille} (Lecture n°5)}

Appliquons la méthode sur le passage climactique où Meka, après la cérémonie, reste seul avec sa médaille (fin du roman, lecture méthodique n°5).

\begin{exoresolu}[Reformulation d'un passage clé]
\textbf{Énoncé :} Reformulez le passage suivant en \textbf{35 mots maximum} (extrait adapté, fin du roman) :
\textit{« Meka restait là, immobile, la médaille au cou. Il ne ressentait plus ni froid ni faim. Une seule pensée l'habitait : il avait été dupé. La médaille n'était pas un honneur, c'était une chaîne. Une chaîne dorée qui l'attachait à jamais à ses bourreaux. Il avait servi l'ennemi, croyant servir son pays. »}
\tcblower
\textbf{Solution détaillée :}

\begin{tabularx}{\linewidth}{|l|X|}
\hline
\rowcolor{popblueL}
\textbf{Étape} & \textbf{Travail} \\
\hline
\textbf{1. Idées forces repérées} & 1. Meka immobile, médaille au cou. 2. Insensibilité physique (froid/faim). 3. Prise de conscience : duperie. 4. Médaille = chaîne dorée (symbole). 5. Attachement aux bourreaux. 6. Erreur : servi l'ennemi au lieu du pays. \\
\hline
\textbf{2. Réduction (mots-clés)} & Meka immobile, médaille, insensible, conscience duperie, médaille = chaîne dorée, lié aux bourreaux, servi ennemi $\neq$ pays. \\
\hline
\textbf{3. Reformulation (brouillon)} & \textit{Ébauche 1 :} Meka, figé, ne sent plus rien. Il réalise qu'on l'a trompé : sa médaille est une chaîne dorée l'enchaînant à ses oppresseurs. Il a servi l'ennemi, pas son pays. (28 mots) \\
\hline
\textbf{4. Polissage \& Liaison} & \textit{Version finale :} \textbf{Figé, insensible au froid comme à la faim, Meka réalise soudain sa duperie : la médaille, chaîne dorée, le lie à jamais à ses bourreaux ; il a servi l'ennemi, croyant servir son pays.} \\
\hline
\textbf{5. Compte-mots} & \textbf{Nombre de mots : 32} (dans la fourchette 35 $\pm 10\%$, soit 32 à 38 mots). \\
\hline
\end{tabularx}

\textbf{Analyse des choix de reformulation :}
\begin{itemize}
    \item \emph{« restait là, immobile »} $\rightarrow$ \textbf{Figé} (participe unique, plus concis).
    \item \emph{« ne ressentait plus ni froid ni faim »} $\rightarrow$ \textbf{insensible au froid comme à la faim} (structure comparative, économie de « ne... plus »).
    \item \emph{« Une seule pensée l'habitait : il avait été dupé »} $\rightarrow$ \textbf{réalise soudain sa duperie} (nominalisation + verbe unique).
    \item \emph{« La médaille n'était pas un honneur, c'était une chaîne. Une chaîne dorée... »} $\rightarrow$ \textbf{la médaille, chaîne dorée,} (apposition, suppression de la phrase nominale redondante).
    \item \emph{« qui l'attachait à jamais à ses bourreaux »} $\rightarrow$ \textbf{le lie à jamais à ses bourreaux} (verbe au présent de narration, pronom relatif absorbé).
    \item \emph{« Il avait servi l'ennemi, croyant servir son pays »} $\rightarrow$ \textbf{il a servi l'ennemi, croyant servir son pays} (conservé car déjà dense ; « croyant », participe présent, garde la nuance d'erreur subjective).
\end{itemize}
\end{exoresolu}

\courssub{Grille d'auto-évaluation (à utiliser sur brouillon avant la copie propre)}

Avant de recopier votre résumé, cochez honnêtement chaque critère. Un « Non » = retour à l'étape correspondante.

\begin{center}
\begin{tabularx}{\linewidth}{|c|l|X|c|}
\hline
\rowcolor{popblueL}
\textbf{N°} & \textbf{Critère} & \textbf{Question de contrôle} & \textbf{OK ?} \\
\hline
1 & \textbf{Fidélité} & Toutes les idées forces du texte sont-elles présentes ? Aucune idée ajoutée ? Aucun contresens ? & $\square$ \\
\hline
2 & \textbf{Reformulation} & Aucun segment de $> 4$ mots recopié ? Vocabulaire personnel ? Structures syntaxiques variées (passif, nominalisation) ? & $\square$ \\
\hline
3 & \textbf{Liaison / Cohérence} & Les idées s'enchaînent-elles logiquement ? Connecteurs appropriés ? Temps verbal unique ? Références (pronoms) claires ? & $\square$ \\
\hline
4 & \textbf{Langue} & Orthographe, grammaire, syntaxe correctes ? Ponctuation rigoureuse ? Style neutre (pas de « je », pas de familiarité) ? & $\square$ \\
\hline
5 & \textbf{Longueur} & Nombre de mots compté \textbf{exactement} ? Dans la fourchette $\pm 10\%$ ? Total indiqué en fin de texte ? & $\square$ \\
\hline
\end{tabularx}
\end{center}

\begin{attention}[Pièges de dernière minute]
\begin{itemize}
    \item \textbf{Oublier le compte-mots} ou l'écrire \emph{avant} la rédaction définitive (le chiffre doit correspondre à la copie finale).
    \item \textbf{Changer de temps} en cours de route (ex. : présent puis passé simple). Choisissez le \textbf{présent de narration} (standard au Bac) et tenez-le.
    \item \textbf{Laisser des abréviations} (etc., c.-à-d., M.) ou des symboles : tout s'écrit en toutes lettres.
    \item \textbf{Ne pas relire à voix basse} : l'oreille détecte les lourdeurs, les répétitions, les ruptures de rythme que l'œil manque.
\end{itemize}
\end{attention}

\courssub{Exercice d'entraînement}

\begin{exercice}[Reformulation guidée — Sujet type Bac]
\textbf{Texte support (extrait fictif, style administratif / rapport) :}
\textit{« Il convient de souligner que la mise en œuvre du nouveau programme d'électrification rurale a connu des retards significatifs. Ces retards sont imputables, pour une large part, à la difficulté d'acheminer le matériel lourd vers les sites d'implantation, du fait de l'état dégradé des pistes d'accès. Par ailleurs, les populations locales n'ont pas été suffisamment associées aux phases de concertation préalable, ce qui a engendré des résistances passives lors de l'installation des poteaux. Néanmoins, les premières lignes mises sous tension ont permis de constater un impact immédiat sur l'activité économique des villages concernés : multiplication des petits ateliers de soudure, allongement des heures d'ouverture des commerces, amélioration de la conservation des produits périssables. »}

\textbf{Consigne :} Résumez ce texte en \textbf{60 mots $\pm 10\%$ (soit 54 à 66 mots)}. Vous appliquerez la méthode en 5 étapes. Indiquez le nombre de mots à la fin.
\end{exercice}

\begin{corrige}[Exercice — Reformulation guidée]
\textbf{Étapes 1 \& 2 (Repérage \& Réduction) :}
Idées forces : 1. Retards du programme d'électrification rurale. 2. Cause 1 : acheminement du matériel difficile (pistes dégradées). 3. Cause 2 : populations non concertées $\rightarrow$ résistances. 4. Résultat positif (lignes sous tension) : relance de l'économie locale (ateliers, commerces, conservation).

\textbf{Étape 3 (Reformulation segment par segment) :}
\begin{itemize}
    \item \textit{« mise en œuvre... a connu des retards »} $\rightarrow$ \textbf{Le programme d'électrification rurale accuse un retard notable.}
    \item \textit{« imputables... difficulté d'acheminer... état dégradé des pistes »} $\rightarrow$ \textbf{Il provient de l'acheminement malaisé du matériel par des pistes dégradées.}
    \item \textit{« populations pas associées... résistances passives »} $\rightarrow$ \textbf{L'absence de concertation locale a suscité des résistances à l'installation.}
    \item \textit{« Néanmoins... impact immédiat... multiplication... allongement... amélioration »} $\rightarrow$ \textbf{Pourtant, la mise sous tension des premières lignes relance déjà l'économie villageoise : ateliers, commerces prolongés, meilleure conservation.}
\end{itemize}

\textbf{Étape 4 (Rédaction suivie \& Liaison) :}
\textbf{Le programme d'électrification rurale accuse un retard notable, dû à l'acheminement malaisé du matériel par des pistes dégradées et à l'absence de concertation locale, source de résistances. Pourtant, la mise sous tension des premières lignes relance déjà l'économie villageoise : multiplication des ateliers, allongement des heures de commerce, meilleure conservation des denrées.}

\textbf{Étape 5 (Vérification \& Compte-mots) :}
Comptage : Le(1) programme(2) d'électrification(3) rurale(4) accuse(5) un(6) retard(7) notable(8), dû(9) à(10) l'acheminement(11) malaisé(12) du(13) matériel(14) par(15) des(16) pistes(17) dégradées(18) et(19) à(20) l'absence(21) de(22) concertation(23) locale(24), source(25) de(26) résistances(27). Pourtant(28), la(29) mise(30) sous(31) tension(32) des(33) premières(34) lignes(35) relance(36) déjà(37) l'économie(38) villageoise(39) : multiplication(40) des(41) ateliers(42), allongement(43) des(44) heures(45) de(46) commerce(47), meilleure(48) conservation(49) des(50) denrées(51).
\textbf{Nombre de mots : 51.}

\textbf{Attention :} 51 mots $<$ 54 (borne basse). Il faut développer légèrement pour atteindre la fourchette.

\textbf{Ajustement (ajout de précisions utiles, sans idée nouvelle) :}
\textbf{Le programme d'électrification rurale accuse un retard notable, imputable à l'acheminement difficile du matériel lourd par des pistes dégradées et au défaut de concertation des populations, qui a provoqué des résistances. Néanmoins, la mise sous tension des premières lignes dynamise déjà l'économie locale : ateliers de soudure, commerces ouverts plus tard, conservation améliorée des produits périssables.}

\textbf{Nouveau comptage :} Le(1) programme(2) d'électrification(3) rurale(4) accuse(5) un(6) retard(7) notable(8), imputable(9) à(10) l'acheminement(11) difficile(12) du(13) matériel(14) lourd(15) par(16) des(17) pistes(18) dégradées(19) et(20) au(21) défaut(22) de(23) concertation(24) des(25) populations(26), qui(27) a(28) provoqué(29) des(30) résistances(31). Néanmoins(32), la(33) mise(34) sous(35) tension(36) des(37) premières(38) lignes(39) dynamise(40) déjà(41) l'économie(42) locale(43) : ateliers(44) de(45) soudure(46), commerces(47) ouverts(48) plus(49) tard(50), conservation(51) améliorée(52) des(53) produits(54) périssables(55).
\textbf{Nombre de mots : 55. Dans la fourchette [54 ; 66]. Valide.}
\end{corrige}

\begin{aretenir}[Synthèse : Les 3 piliers d'une reformulation réussie]
\begin{enumerate}
    \item \textbf{Le « Dire autrement » (Autonomie) :} Synonymes, conversion de voix, nominalisation, regroupement. \textbf{Zéro recopiage} ($> 4$ mots).
    \item \textbf{Le « Dire juste » (Fidélité) :} Toutes les idées forces, rien que les idées forces, dans l'ordre logique, sans commentaire personnel.
    \item \textbf{Le « Dire bien » (Cohérence \& Concision) :} Connecteurs logiques, temps unique, phrases complètes, longueur respectée $\pm 10\%$, compte-mots exact indiqué.
\end{enumerate}
\end{aretenir}

\coursec{Compte rendu et remédiation}

Après avoir appris à repérer les idées essentielles d'un texte argumentatif et à les reformuler avec ses propres mots, il reste une dernière étape, souvent négligée mais décisive : \cle{évaluer} ce que l'on a produit, \cle{corriger} ses erreurs et \cle{progresser}. C'est tout l'objet du compte rendu et de la remédiation. Un résumé n'est jamais définitivement « fini » à la première version : c'est en comparant sa copie à un corrigé, en identifiant précisément ses écarts et en s'entraînant sur des exercices ciblés que l'élève transforme une compétence fragile en compétence solide pour le jour du Probatoire ou du Baccalauréat technique.

\courssub{Correction collective d'un résumé type}

\begin{definition}[Le compte rendu]
Le \cle{compte rendu} est le moment où la classe confronte sa production à une référence : le corrigé du professeur ou un résumé collectif construit au tableau. Il ne s'agit pas de « donner la bonne réponse » à recopier, mais de rendre visibles les \cle{réussites} (ce qui est conforme aux attentes) et les \cle{écarts} (ce qui s'en éloigne), afin que chacun comprenne \emph{pourquoi} sa production est réussie ou perfectible.
\end{definition}

La correction collective se déroule en trois temps. D'abord, le professeur projette ou dicte un \cle{résumé type} (le corrigé). Ensuite, la classe relit ce corrigé critère par critère : les idées essentielles sont-elles toutes présentes ? Sont-elles reformulées ? Les liaisons sont-elles assurées ? Enfin, chaque élève annote sa propre copie en marge, sans la recopier : il souligne en vert ce qui est réussi, en rouge ce qui manque ou ce qui est fautif.

\begin{methode}[Comparer son résumé au corrigé : les 5 critères de la grille]
Pour auto-évaluer votre résumé, vous vérifiez successivement cinq critères :
\begin{enumerate}
\item \textbf{Fidélité} : la thèse de l'auteur et toutes les idées essentielles sont présentes, sans déformation ni ajout d'opinion personnelle.
\item \textbf{Reformulation} : aucune phrase n'est copiée du texte ; le vocabulaire et les tournures sont les vôtres.
\item \textbf{Liaison} : les idées sont enchaînées par des connecteurs logiques (d'abord, ensuite, en effet, cependant, donc...) et des conjonctions correctement employées.
\item \textbf{Langue} : orthographe, grammaire et accords sont corrects ; les verbes d'énonciation sont au présent ou au passé selon la consigne.
\item \textbf{Longueur} : le nombre de mots respecte la fourchette imposée (par exemple le quart du texte, à 10 \% près).
\end{enumerate}
Chaque critère non respecté est un écart à traiter par une remédiation ciblée.
\end{methode}

\begin{exemplebox}[Une correction collective en pratique]
Sur un texte de 400 mots de Ferdinand Oyono à résumer en 100 mots, le corrigé collectif compte 4 idées essentielles. La classe constate : la majorité a bien conservé la thèse (réussite), mais la moitié des copies a recopié la phrase d'ouverture du texte (écart de reformulation) et un tiers dépasse 130 mots (écart de longueur). Le professeur note ces deux écarts au tableau : ce seront les priorités de remédiation.
\end{exemplebox}

\courssub{Typologie des erreurs fréquentes}

L'observation de nombreuses copies d'élèves de Terminale technique permet de classer les erreurs en quatre grandes familles. Les reconnaître, c'est déjà commencer à les éviter.

\begin{enumerate}
\item \textbf{La copie du texte.} L'élève recopie des phrases entières, parfois en les « coupant ». Or un résumé qui copie n'est pas un résumé : il ne prouve ni la compréhension ni la capacité de reformulation, et il est sanctionné à l'examen.
\item \textbf{L'ajout d'opinions.} L'élève commente (« je trouve que l'auteur a raison », « cette histoire est triste »). Le résumé est un exercice de \cle{fidélité} : il rend compte de la pensée de l'auteur, jamais de celle du lecteur.
\item \textbf{La perte de la thèse.} En supprimant trop, l'élève garde des exemples ou des détails et oublie l'idée directrice. Le résumé devient une suite de faits sans fil conducteur : c'est l'erreur la plus grave, car elle trahit le texte.
\item \textbf{Le non-respect de la longueur.} Résumé trop long (l'élève n'a pas assez réduit) ou trop court (il a sacrifié des idées essentielles). La longueur imposée fait partie de la consigne : la dépasser de plus de 10 \% est une faute de contrat de communication.
\end{enumerate}

\begin{attention}[La remédiation n'est pas une punition]
Beaucoup d'élèves vivent la correction comme un jugement : « j'ai eu une mauvaise note, je suis nul en résumé ». C'est une erreur de perspective. La \cle{remédiation} est une \textbf{seconde chance d'apprentissage ciblée} : elle part de l'erreur précise que \emph{vous} avez commise et \emph{vous} propose un exercice précis pour la dépasser. Un élève qui comprend ses erreurs et s'entraîne dessus progresse plus vite qu'un élève qui réussit par hasard. Au Baccalauréat technique, ce sont les erreurs travaillées en classe qui ne se reproduisent plus dans la copie.
\end{attention}

\courssub{Tableau de remédiation : de l'erreur à l'exercice correctif}

Le principe de la remédiation efficace est simple : à chaque erreur correspond une \cle{cause}, et à chaque cause un \cle{exercice correctif} adapté. On ne « refait » pas globalement le résumé : on répare la pièce défectueuse.

\begin{popfigure}
\begin{center}
\rowcolors{2}{popblueL}{white}
\begin{tabularx}{\linewidth}{|l|X|X|}
\hline
\rowcolor{popblue}\textcolor{white}{\textbf{Erreur observée}} & \textcolor{white}{\textbf{Cause probable}} & \textcolor{white}{\textbf{Exercice correctif proposé}} \\
\hline
Copie de phrases du texte & Peur de mal dire ; reformulation non maîtrisée & Reformuler 5 phrases isolées (changement de vocabulaire et de structure) avant tout résumé complet \\
\hline
Ajout d'opinions personnelles & Confusion entre résumé et commentaire & Souligner dans un résumé fautif toutes les marques de subjectivité (je, à mon avis) et les supprimer \\
\hline
Perte de la thèse & Hiérarchisation défaillante : exemples pris pour des idées & Classer 8 énoncés d'un texte en « idée essentielle / exemple / détail » \\
\hline
Résumé trop long & Suppression insuffisante & Exercice de suppression : rayer dans un paragraphe tout ce qui peut disparaître sans perte de sens \\
\hline
Résumé trop court & Généralisation excessive & Exercice de fusion de phrases : réunir deux phrases proches sans perdre aucune information \\
\hline
Enchaînements absents ou fautifs & Conjonctions de coordination et de subordination mal maîtrisées & Compléter un résumé à trous avec mais, donc, car, parce que, bien que, puis justifier chaque choix \\
\hline
\end{tabularx}
\end{center}
\legende{Tableau de remédiation : chaque erreur renvoie à une cause identifiée et à un exercice correctif précis. L'élève ne s'entraîne que sur ses propres écarts.}
\end{popfigure}

\courssub{Activités de remédiation différenciées}

Tous les élèves n'ont pas les mêmes erreurs : il serait absurde de faire faire le même exercice à toute la classe. La remédiation est donc \cle{différenciée} : chacun travaille sur son point faible, identifié grâce à la grille des cinq critères. Trois ateliers types couvrent l'essentiel des besoins.

\textbf{Atelier 1 : la suppression.} On donne un paragraphe chargé de répétitions, d'exemples et de détails. Consigne : rayer tout ce qui peut disparaître sans que le sens essentiel soit perdu, puis compter les mots restants. Cet exercice entraîne directement la technique de réduction par élimination.

\textbf{Atelier 2 : la généralisation.} On donne une énumération (« les manguiers, les palmiers, les bananiers et les goyaviers bordaient la cour ») et l'on demande de la remplacer par un terme générique (« des arbres fruitiers bordaient la cour »). L'élève apprend à remonter du particulier au général, geste fondamental de la contraction.

\textbf{Atelier 3 : la fusion de phrases.} On donne deux ou trois phrases courtes exprimant des idées proches ; il faut les fondre en une seule phrase correcte, en utilisant une conjonction de coordination ou de subordination adaptée. Cet atelier travaille à la fois la réduction et la liaison.

\begin{exemplebox}[Un parcours de remédiation individualisé]
Un élève qui copie le texte ne doit pas refaire immédiatement un résumé complet : il recommencerait la même erreur. Il s'entraîne \textbf{d'abord} sur la reformulation de phrases isolées : on lui donne « Le vieux Meka demeura immobile, écrasé par la honte de cette médaille dérisoire », et il produit par exemple « Meka resta figé, humilié par cette récompense sans valeur ». Après cinq phrases réussies, il passe à un paragraphe, puis seulement ensuite à un résumé complet. La remédiation progresse du simple au complexe.
\end{exemplebox}

\begin{exoresolu}[Atelier de fusion de phrases]
Énoncé : fusionnez en une seule phrase les trois énoncés suivants, sans perdre aucune information essentielle : « Meka a reçu une médaille. Cette médaille récompense sa longue patience. Meka se sent pourtant humilié. »
\tcblower
Solution détaillée : on repère d'abord les relations logiques entre les phrases : la deuxième phrase \emph{explique} la première (cause), la troisième \emph{s'oppose} aux deux premières (opposition). On choisit donc une conjonction de subordination de cause et un connecteur d'opposition. Proposition : « Meka a reçu une médaille qui récompensait sa longue patience, mais il se sentait humilié. » Variante avec subordination : « Bien que la médaille reçue par Meka récompensât sa longue patience, il se sentait humilié. » On vérifie : aucune information perdue (médaille, patience, humiliation), une seule phrase, liaisons correctes.
\end{exoresolu}

\courssub{Réécriture individuelle du résumé}

Après la correction collective et les ateliers, chaque élève \cle{réécrit} intégralement son résumé, chez lui ou en classe, en s'appuyant sur la grille d'auto-évaluation remplie. Cette réécriture n'est pas une simple copie du corrigé : l'élève repart de \emph{sa} première version, corrige ses écarts un à un et produit une version améliorée, qu'il accompagne d'une courte justification (« j'ai supprimé l'exemple du village car c'était un détail ; j'ai remplacé la phrase copiée par une reformulation ; j'ai ajouté \emph{cependant} pour marquer l'opposition »). Justifier sa démarche est un savoir-faire exigé par le programme : c'est ce qui transforme l'exercice en apprentissage durable.

\begin{methode}[Réussir sa réécriture en 4 étapes]
\begin{enumerate}
\item Relire votre grille d'auto-évaluation et lister vos écarts par ordre de gravité (perte de la thèse d'abord, longueur ensuite).
\item Corriger chaque écart avec la technique apprise en atelier (suppression, généralisation, fusion, reformulation).
\item Relire à voix haute pour vérifier la fluidité et les liaisons.
\item Compter les mots et vérifier le respect de la fourchette imposée avant de rendre la copie.
\end{enumerate}
\end{methode}

\courssub{Bilan des acquis : quiz rapide}

La séquence se clôt par un quiz de validation, à faire en dix minutes, sans document. Il porte sur les trois piliers de l'unité : la structure argumentative, les outils de liaison et les techniques de réduction.

\begin{exercice}[Quiz de fin d'unité]
\begin{enumerate}
\item Quels sont les trois éléments de la structure d'un texte argumentatif ?
\item Une idée essentielle est-elle le plus souvent exprimée par un exemple ou par un énoncé général ?
\item Classez les mots suivants : \emph{mais, car, donc, parce que, bien que, et, cependant} — conjonctions de coordination ou de subordination ?
\item Vrai ou faux : dans un résumé, on peut donner son avis si l'on est d'accord avec l'auteur.
\item Citez les trois techniques de réduction étudiées.
\item Un texte fait 600 mots ; la consigne impose le quart. Quelle fourchette de mots est acceptable (à 10 \% près) ?
\end{enumerate}
\end{exercice}

\begin{corrige}[Quiz de fin d'unité]
\begin{enumerate}
\item Une thèse (le point de vue défendu), des arguments (qui la justifient) et des exemples (qui illustrent les arguments).
\item Par un énoncé général ; l'exemple est un détail que le résumé élimine.
\item Coordination : \emph{mais, donc, et, car}. Subordination : \emph{parce que, bien que} ; \emph{cependant} est un connecteur (adverbe), ni l'une ni l'autre.
\item Faux : le résumé doit être strictement fidèle à la pensée de l'auteur, sans aucune opinion personnelle.
\item La suppression (des détails et répétitions), la généralisation (du particulier au général) et la fusion (de phrases proches).
\item Le quart de 600 est 150 mots ; à 10 \% près, la fourchette acceptable va de 135 à 165 mots.
\end{enumerate}
\end{corrige}

\begin{aretenir}[L'essentiel de la section]
Le compte rendu et la remédiation transforment l'erreur en progrès. On compare son résumé au corrigé selon cinq critères : \cle{fidélité}, \cle{reformulation}, \cle{liaison}, \cle{langue}, \cle{longueur}. Les erreurs fréquentes sont la copie du texte, l'ajout d'opinions, la perte de la thèse et le non-respect de la longueur. À chaque erreur correspond un exercice correctif ciblé : suppression, généralisation, fusion de phrases ou reformulation. La réécriture individuelle, justifiée à partir de la grille d'auto-évaluation, consolide les acquis. La remédiation n'est pas une punition : c'est une seconde chance d'apprendre, et c'est elle qui prépare efficacement à l'épreuve de contraction de texte du Baccalauréat technique.
\end{aretenir}

\newpage

\coursec{Activité d'intégration}

\coursec{Leçon 15 : Formuler les idées essentielles d'un texte à résumer}

\courssub{Objectifs de l'activité}

À l'issue de cette activité d'intégration, l'élève doit être capable de :

\begin{itemize}
\item repérer la \cle{structure argumentative} d'un texte (thèse, arguments, exemples, connecteurs logiques) ;
\item distinguer les \cle{contenus manifestes} (ce qui est dit explicitement) des \cle{contenus latents} (ce qui est suggéré, sous-entendu) ;
\item appliquer les \cle{techniques de réduction} (suppression des exemples, des répétitions, des figures de style ; regroupement des idées) ;
\item \cle{reformuler} les idées essentielles avec ses propres mots, sans copier le texte ;
\item employer correctement les \cle{conjonctions de coordination et de subordination} pour lier les idées ;
\item produire un \cle{résumé} de 150 mots ($\pm$ 10 \%) conforme aux exigences du Probatoire / Baccalauréat technique, puis s'auto-évaluer à l'aide d'une grille.
\end{itemize}

\courssub{Situation}

La région du Nord-Ouest compte de nombreux villages producteurs de café, de maïs et d'ignames. Pourtant, pendant la saison des pluies, les pistes deviennent impraticables : les produits pourrissent au village, les malades peinent à rejoindre les hôpitaux et les élèves manquent l'école. Un journal local publie l'éditorial suivant (environ 600 mots), intitulé \emph{Des pistes pour vivre} :

\begin{exemplebox}[Texte à résumer : « Des pistes pour vivre » (extrait représentatif)]
« Qui a déjà vu un camion chargé de café renversé dans la boue, à dix kilomètres de la première route goudronnée ? Moi, oui. Et j'ai vu les yeux du planteur qui regardait sa récolte d'une année entière se mélanger à la terre rouge. Voilà le vrai visage de notre sous-développement : ce n'est pas la paresse, ce n'est pas le manque de volonté, c'est l'absence de pistes rurales.

Nos paysans ne demandent ni l'aumône ni les discours. Ils demandent une route. Car une piste, c'est d'abord la liberté : la liberté de vendre sa récolte au juste prix au lieu de la brader au premier acheteur venu ; la liberté de conduire une femme en couches à l'hôpital avant qu'il ne soit trop tard ; la liberté pour nos enfants d'aller à l'école même quand le ciel se fâche.

Les chiffres parlent d'eux-mêmes. Dans notre département, près de 60 \% des villages sont enclavés au moins quatre mois par an. Le coût du transport d'un sac de maïs double, parfois triple, pendant la saison des pluies. Résultat : le producteur vend à perte, le consommateur achète cher, et l'État perd des recettes fiscales. Tout le monde est perdant.

Certains objectent que les pistes coûtent cher. C'est vrai. Mais que coûte l'inaction ? Chaque année, des tonnes de denrées pourrissent faute de débouchés. Chaque année, des jeunes quittent le village, découragés, grossissant les rangs des chômeurs en ville. Une piste bien entretenue coûte moins cher que l'exode rural qu'elle aurait pu empêcher.

D'ailleurs, l'expérience de nos voisins le prouve : là où les pistes ont été réhabilitées, les revenus agricoles ont augmenté, les centres de santé ont vu leur fréquentation progresser et les marchés hebdomadaires ont repris vie. La route n'est pas un luxe : elle est la condition première de tout développement local.

Nous demandons donc aux autorités, aux élites et aux communautés elles-mêmes de faire des pistes rurales une priorité absolue. Que chacun y mette du sien : l'État par le financement, les communes par l'entretien, les villageois par la participation. Le développement ne descendra pas du ciel : il passera par nos pistes, ou il ne passera pas. »
\end{exemplebox}

À titre de comparaison, on relit le passage du \emph{Vieux Nègre et la médaille} (Ferdinand Oyono) où le commandant, lors de la cérémonie de remise de la médaille, déclare en substance que la France et l'Afrique marchent désormais « main dans la main », que Meka est un « bon Français », un ami fidèle. Ce discours solennel, plein de belles promesses, contraste avec la réalité que Meka découvrira bientôt : mépris, arrestation, humiliation. Là encore, il faut lire au-delà des mots.

\textbf{Données de travail :} texte source $\approx$ 600 mots ; résumé attendu : 150 mots ($\pm$ 10 \%, soit entre 135 et 165 mots) ; au moins \textbf{trois conjonctions de subordination} exigées ; durée : 2 séances (travail de groupe puis production individuelle).

\courssub{Consignes}

\textbf{Travail en groupe (séance 1) :}

\begin{enumerate}
\item Lisez le texte en entier, puis relevez la \cle{thèse} défendue par l'éditorialiste. Formulez-la en une phrase avec vos propres mots.
\item Construisez le tableau de la structure argumentative : pour chaque argument, indiquez le paragraphe, l'idée essentielle et l'exemple ou le chiffre qui l'illustre.
\item Relevez deux passages où le contenu est \cle{latent} (sous-entendu, ironie, question rhétorique) et dites ce que l'auteur suggère réellement.
\item Comparez avec le discours du commandant dans \emph{Le Vieux Nègre et la médaille} : quel contenu latent se cache derrière ses belles paroles ? Quelle leçon de lecture en tirez-vous pour tout texte argumentatif ?
\item Appliquez les techniques de réduction : barrez mentalement les exemples, les répétitions, les questions rhétoriques ; regroupez les idées voisines. Vérifiez que vous conservez environ un quart du volume initial.
\end{enumerate}

\textbf{Travail individuel (séance 2) :}

\begin{enumerate}
\setcounter{enumi}{5}
\item Rédigez un résumé de 150 mots ($\pm$ 10 \%) qui rend compte de la thèse et des arguments essentiels, sans exemple chiffré détaillé, sans phrase recopiée, en employant au moins trois conjonctions de subordination différentes (par exemple : \emph{parce que, puisque, bien que, si, lorsque, afin que}).
\item Auto-évaluez votre production à l'aide de la grille ci-dessous, puis remettez-la pour la mise en commun et la remédiation.
\end{enumerate}

\begin{center}
\begin{tabularx}{\linewidth}{|l|c|X|}
\hline
\rowcolor{popblueL} \textbf{Critère} & \textbf{Points} & \textbf{Indicateur} \\
\hline
Fidélité à la thèse & 4 & La thèse de l'auteur est restituée sans déformation ni opinion personnelle. \\
\hline
Arguments essentiels & 4 & Les 3 ou 4 arguments majeurs sont présents, les exemples supprimés. \\
\hline
Reformulation & 4 & Aucune phrase copiée ; vocabulaire personnel correct. \\
\hline
Liaison des idées & 4 & Au moins 3 conjonctions de subordination bien employées ; enchaînement logique. \\
\hline
Respect du volume et langue & 4 & 135 à 165 mots ; orthographe et syntaxe correctes. \\
\hline
\textbf{Total} & \textbf{20} & \\
\hline
\end{tabularx}
\end{center}

\courssub{Ressources à mobiliser}

\begin{aretenir}[Boîte à outils du résumé]
\begin{itemize}
\item \textbf{Structure argumentative :} thèse (point de vue défendu) + arguments (raisons) + exemples (illustrations à supprimer dans le résumé).
\item \textbf{Contenu manifeste / latent :} le manifeste est dit ; le latent est suggéré (ironie, questions rhétoriques, écart entre paroles et réalité, comme chez Oyono).
\item \textbf{Réduction :} supprimer exemples, répétitions, figures ; regrouper les idées proches ; viser environ le quart du texte.
\item \textbf{Reformulation :} changer les mots et les tournures, jamais les idées ; ne pas donner son avis.
\item \textbf{Conjonctions de coordination :} \emph{mais, ou, et, donc, or, ni, car}. \textbf{Conjonctions de subordination :} \emph{que, parce que, puisque, bien que, si, lorsque, quoique, afin que}.
\end{itemize}
\end{aretenir}

\courssub{Démarche et corrigé commenté}

\begin{exoresolu}[Résolution guidée de l'activité]
Énoncé : produire le résumé de 150 mots de l'éditorial « Des pistes pour vivre » en suivant les consignes 1 à 7.
\tcblower
\textbf{Étape 1 — Dégager la thèse.} Le titre, la fin du premier paragraphe et la chute convergent : \emph{la construction et l'entretien des pistes rurales sont la condition première du développement local dans la région du Nord-Ouest}. C'est le contenu manifeste central.

\textbf{Étape 2 — Tableau de la structure argumentative.}

\begin{center}
\begin{tabularx}{\linewidth}{|c|X|X|}
\hline
\rowcolor{popblueL} \textbf{Paragraphe} & \textbf{Idée essentielle (à garder)} & \textbf{Exemple / chiffre (à supprimer)} \\
\hline
1 & L'enclavement, et non la paresse, explique le sous-développement. & Le camion de café renversé dans la boue. \\
\hline
2 & Une piste, c'est la liberté de vendre, de se soigner, d'étudier. & La femme en couches, les écoliers. \\
\hline
3 & L'enclavement ruine producteurs, consommateurs et État. & 60 \% de villages enclavés ; prix du transport qui triple. \\
\hline
4 & L'inaction coûte plus cher que l'investissement. & L'exode rural des jeunes. \\
\hline
5 & L'expérience des régions voisines prouve l'efficacité des pistes. & Marchés et centres de santé relancés. \\
\hline
6 & Appel à la mobilisation de tous : État, communes, villageois. & La formule finale frappante. \\
\hline
\end{tabularx}
\end{center}

\textbf{Étape 3 — Contenus latents.} (a) La question initiale « Qui a déjà vu... ? » suggère que les décideurs, eux, n'ont jamais rien vu : critique implicite des responsables éloignés du terrain. (b) « Nos paysans ne demandent ni l'aumône ni les discours » sous-entend que les autorités se contentent de promesses verbales. (c) Chez Oyono, le discours du commandant (« bon Français », « main dans la main ») cache un contenu latent de domination et d'hypocrisie : les belles paroles masquent le mépris réel. Leçon : dans tout texte argumentatif, il faut lire ce qui est dit \emph{et} ce qui est tu.

\textbf{Étape 4 — Réduction et regroupement.} On supprime les anecdotes, les chiffres précis et les formules rhétoriques ; on regroupe : arguments économiques (paragraphes 3 et 4), arguments sociaux (paragraphe 2), preuve par l'expérience (paragraphe 5), appel final (paragraphe 6).

\textbf{Étape 5 — Modèle de résumé (150 mots).}

\emph{Selon l'éditorialiste, le sous-développement de la région du Nord-Ouest ne vient pas de la paresse des paysans, mais de l'enclavement des villages, parce que les pistes rurales font défaut. Il affirme qu'une piste offre la liberté de vendre sa récolte au juste prix, de se soigner à temps et d'envoyer ses enfants à l'école. Sur le plan économique, l'enclavement pénalise tout le monde : le producteur vend à perte, le consommateur paie cher et l'État perd des recettes. Bien que leur construction coûte cher, les pistes reviennent moins cher que l'inaction, puisque celle-ci provoque le gaspillage des récoltes et l'exode des jeunes. L'expérience des régions voisines montre d'ailleurs que, lorsque les pistes sont réhabilitées, les revenus et les services sociaux progressent. L'auteur appelle donc l'État, les communes et les villageois à s'unir afin que les pistes rurales deviennent une priorité, car le développement local passera par elles ou ne sera pas.}

Vérification : 150 mots ; conjonctions de subordination utilisées : \emph{parce que, bien que, puisque, lorsque, afin que} (cinq, soit plus que le minimum exigé) ; aucune phrase copiée ; thèse et quatre arguments restitués ; aucun avis personnel.

\textbf{Étape 6 — Auto-évaluation.} Chaque élève se note sur 20 avec la grille, puis la classe confronte les productions en mise en commun ; les points faibles (copie de phrases, oubli de la thèse, subordonnées mal construites) font l'objet d'une remédiation ciblée.
\end{exoresolu}

\begin{attention}[Pièges fréquents au Probatoire / Baccalauréat technique]
\begin{itemize}
\item \textbf{Copier des phrases du texte} : c'est la faute la plus sanctionnée. On reformule toujours.
\item \textbf{Donner son avis} (« je pense que l'auteur a raison... ») : le résumé restitue la pensée de l'auteur, pas la vôtre.
\item \textbf{Garder les exemples et les chiffres} : ils disparaissent au profit de l'idée qu'ils illustrent.
\item \textbf{Dépasser ou sous-estimer le volume} : comptez vos mots ; hors de la fourchette 135--165, des points sont perdus.
\item \textbf{Confondre coordination et subordination} : \emph{car} coordonne, \emph{parce que} subordonne ; vérifiez la nature de chaque connecteur.
\end{itemize}
\end{attention}

\courssub{Bilan de l'activité}

Cette activité a permis de mobiliser en une seule tâche l'ensemble des savoir-faire de la leçon. Les élèves ont d'abord \textbf{lu pour comprendre} : repérage de la thèse et des arguments, distinction entre contenus manifestes et contenus latents, mise en relation avec le discours trompeur du commandant chez Oyono, qui a montré qu'un texte argumentatif se lit aussi entre les lignes. Ils ont ensuite \textbf{lu pour réduire} : suppression des exemples et des figures, regroupement des idées voisines, respect du quart du volume. Ils ont enfin \textbf{écrit pour reformuler} : production d'un résumé personnel, lié par des conjonctions de subordination, fidèle à la pensée de l'auteur. L'auto-évaluation par grille et la remédiation collective ont rendu chaque élève acteur de ses progrès, conformément aux exigences de l'épreuve de contraction de texte du Probatoire / Baccalauréat technique : fidélité, brièveté, clarté et correction de la langue.

\newpage

\coursec{Méthodes et exercices résolus}

\coursec{III : Formuler les idées essentielles d'un texte à résumer}

\courssub{Le texte argumentatif : structure et enjeux}

\begin{methode}[Analyser la structure d'un texte argumentatif]
    \textbf{Objectif :} Identifier la thèse, les arguments et les exemples pour comprendre la logique du texte avant toute contraction.
    \begin{enumerate}
        \item \textbf{Repérer la thèse :} C'est l'idée principale défendue par l'auteur (souvent en introduction ou conclusion). Se demander : « Que veut-il prouver ? ».
        \item \textbf{Identifier les arguments :} Ce sont les raisons logiques qui soutiennent la thèse. Les repérer grâce aux connecteurs logiques (\cle{donc}, \cle{car}, \cle{puisque}, \cle{en effet}, \cle{par conséquent}).
        \item \textbf{Distinguer les exemples/illustrations :} Faits, anecdotes, statistiques, citations qui concrétisent les arguments. Ils sont souvent introduits par : \cle{par exemple}, \cle{ainsi}, \cle{notamment}, \cle{comme le montre}.
        \item \textbf{Schématiser la structure :} Thèse $\rightarrow$ Argument 1 (+ Exemple) $\rightarrow$ Argument 2 (+ Exemple) $\rightarrow$ ... $\rightarrow$ Conclusion.
        \item \textbf{Évaluer l'enjeu :} Quel est le but ? Convaincre (raison) ? Persuader (émotion) ? Dénoncer ? Proposer ?
    \end{enumerate}
    \textbf{Reflexe Bac :} Sur ton brouillon, note la thèse en une phrase et liste les arguments avec des mots-clés (pas des phrases entières).
\end{methode}

\begin{exoresolu}[Analyse structurelle d'un extrait argumentatif]
    \textbf{Énoncé :} Lisez l'extrait suivant et identifiez la thèse, les arguments principaux et leur nature (logique, autorité, exemple), puis schématisez la structure.
    \begin{quote}
        \textit{« L'école ne doit pas seulement instruire, elle doit éduquer. En effet, l'instruction sans éducation produit des esprits brillants mais sans conscience morale. L'histoire nous le prouve : les plus grands tyrans étaient souvent des hommes très cultivés. Par conséquent, former le caractère est aussi urgent que transmettre le savoir. C'est pourquoi l'éducation civique doit redevenir une priorité absolue dans nos programmes. »}
    \end{quote}
    \tcblower
    \textbf{Solution détaillée :}
    \begin{enumerate}
        \item \textbf{Thèse :} « L'école ne doit pas seulement instruire, elle doit éduquer » (phrase 1). L'auteur défend l'idée que l'éducation morale est indissociable de l'instruction intellectuelle.
        \item \textbf{Arguments :}
            \begin{itemize}
                \item \textbf{Argument 1 (Logique / Causalité) :} « L'instruction sans éducation produit des esprits brillants mais sans conscience morale. » (Introduit par \cle{En effet}). C'est une explication théorique.
                \item \textbf{Argument 2 (Exemple historique / Preuve par l'exemple) :} « L'histoire nous le prouve : les plus grands tyrans étaient souvent des hommes très cultivés. » Illustration concrète pour valider l'argument 1.
            \end{itemize}
        \item \textbf{Conclusion / Conséquence pratique :} « Par conséquent, former le caractère est aussi urgent que transmettre le savoir. C'est pourquoi l'éducation civique doit redevenir une priorité absolue. » (Proposition de solution).
        \item \textbf{Schéma structurel :}
        \begin{center}
            \begin{tabular}{|l|p{10cm}|}\hline
            \textbf{Élément} & \textbf{Contenu} \\ \hline
            Thèse & École = Instruction + Éducation (morale) \\ \hline
            Arg 1 & Instruction seule = danger (esprit sans conscience) \\ \hline
            Exemple & Tyrans cultivés (preuve historique) \\ \hline
            Conclusion & Urgence éducation civique / Priorité programmes \\ \hline
            \end{tabular}
        \end{center}
    \end{enumerate}
    \textbf{Conclusion :} Le texte suit un plan déductif (Thèse $\rightarrow$ Arguments $\rightarrow$ Conséquence). La contraction devra garder : Thèse + Arg 1 synthétisé + Conclusion.
\end{exoresolu}

\courssub{Contenus manifestes et contenus latents}

\begin{methode}[Distinguer le dit et le non-dit (manifestes / latents)]
    \textbf{Objectif :} Ne pas se contenter de résumer ce qui est écrit explicitement, mais restituer la portée réelle du texte (ironie, sous-entendus, valeurs défendues).
    \begin{enumerate}
        \item \textbf{Contenu Manifeste (Explicite) :} Ce que le texte \cle{dit} littéralement. Informations factuelles, arguments énoncés, thèse affichée. Question : « Que lit-on noir sur blanc ? ».
        \item \textbf{Contenu Latent (Implicite) :} Ce que le texte \cle{suggère}, \cle{sous-entend} ou \cle{présuppose}. Intentions réelles, critiques voilées, valeurs culturelles, idéologie. Questions : « Pourquoi le dit-il ainsi ? », « Que cache cette formulation ? », « Quel est le public visé ? ».
        \item \textbf{Outils de détection du latent :}
            \begin{itemize}
                \item Le ton (ironique, sarcastique, lyrique, neutre).
                \item Les modalisateurs (adverbes : \cle{heureusement}, \cle{malheureusement}, \cle{soi-disant} ; verbes : \cle{sembler}, \cle{prétendre}).
                \item Les figures de style (antiphrase, euphémisme, litote).
                \item Les présupposés (ce qui est tenu pour acquis sans être dit).
            \end{itemize}
        \item \textbf{Application à la contraction :} Le résumé/contraction doit rendre compte du \textbf{sens global} (manifeste + latent), pas seulement de la surface textuelle.
    \end{enumerate}
    \textbf{Attention :} Ne pas \emph{inventer} du latent. S'appuyer sur des indices textuels précis.
\end{methode}

\begin{exoresolu}[Analyse des contenus manifestes et latents]
    \textbf{Énoncé :} Analysez les contenus manifestes et latents de cet extrait (inspiré de \emph{Le Vieux Nègre et la Médaille}).
    \begin{quote}
        \textit{« Le Gouverneur remit la médaille au vieux Meka. Ce dernier, tremblant, la regarda longuement. "Une belle chose", murmura-t-il. La foule applaudit. Le Vieux Nègre rentra chez lui, la médaille sur la poitrine, mais le ventre vide. »}
    \end{quote}
    \tcblower
    \textbf{Solution détaillée :}
    \begin{enumerate}
        \item \textbf{Contenus Manifestes (Factuels / Expliques) :}
            \begin{itemize}
                \item Action : Remise de médaille par le Gouverneur au vieux Meka.
                \item Réaction physique : Meka tremble, regarde la médaille, dit « Une belle chose ».
                \item Réaction collective : La foule applaudit.
                \item Situation finale : Meka rentre chez lui, médaille sur la poitrine, ventre vide.
            \end{itemize}
        \item \textbf{Contenus Latents (Implicites / Sous-entendus) :}
            \begin{itemize}
                \item \textbf{L'ironie / Le décalage :} « Une belle chose » (litote/euphémisme) masque le mépris ou l'incompréhension de la valeur symbolique imposée. Le « tremblant » suggère la peur ou la vieillesse, pas l'émotion patriotique attendue.
                \item \textbf{La critique du colonialisme :} La médaille (symbole de la reconnaissance coloniale) est opposée au « ventre vide » (réalité physiologique, misère). Le latent crie : \cle{La décoration ne nourrit pas son homme}.
                \item \textbf{L'aliénation / La comédie :} La foule applaudit (conformisme, spectacle) mais ne résout pas la faim. Le « ventre vide » est le vrai sujet, tu par le protocole.
                \item \textbf{La dignité bafouée :} « Rentra chez lui » : retour à la réalité misérable, la médaille devient un fardeau ironique (« sur la poitrine » vs « ventre vide »).
            \end{itemize}
        \item \textbf{Synthèse pour contraction :} Il ne suffit pas d'écrire « Meka reçoit une médaille et a faim ». Il faut formuler : \cle{La remise solennelle d'une médaille coloniale à un vieux paysan affamé révèle l'absurdité et la cruauté d'un système qui honore symboliquement ceux qu'il laisse mourir de faim.}
    \end{enumerate}
    \textbf{Conclusion :} Le latent (la critique sociale) est plus important que le manifeste (le protocole) pour comprendre l'enjeu du texte.
\end{exoresolu}

\courssub{Lecture méthodique du Vieux Nègre et la Médaille (n°4 et n°5)}

\begin{methode}[Lecture méthodique d'une œuvre au programme : Fiche de synthèse]
    \textbf{Objectif :} Construire une fiche de lecture exploitable pour le résumé, le commentaire ou la dissertation.
    \begin{enumerate}
        \item \textbf{Identité de l'œuvre :} Titre complet, Auteur (Ferdinand Oyono), Genre (Roman / Satire coloniale), Date (1956), Contexte (Cameroun colonial, fin années 50).
        \item \textbf{Résumé de l'intrigue (noyau narratif) :} Meka, vieux paysan, attend la médaille pour avoir donné ses terres/son fils. Cérémonie humiliante. Retour à la misère. Mort symbolique/réelle.
        \item \textbf{Personnages clés :} Meka (victime lucide puis résignée), Kelara (femme lucide, voix de la raison), Le Gouverneur / Commandant / Prêtre (figures de l'autorité coloniale hypocrite), Les « évolués » (traîtres ou perdus).
        \item \textbf{Thèmes majeurs :} \cle{Critique du colonialisme} (aliénation, spoliation, hypocrisie), \cle{Le leurre des honneurs} (médaille vs pain), \cle{La condition paysanne} (pauvreté, dignité), \cle{Le langage} (français approximatif des administrateurs vs langue vernaculaire riche), \cle{La mort} (fin absurde).
        \item \textbf{Procédés stylistiques (pour contraction) :} Ironie, antiphrase, zoom interne (pensées de Meka), oralité transcrite, répétitions (rituel), contraste (solennité / misère).
        \item \textbf{Lectures n°4 \& n°5 (Programme) :} Se concentrer sur \textbf{la structure argumentative de la satire} (comment le récit \emph{argumente} contre le colonialisme) et \textbf{la reformulation des idées essentielles} (passer du récit à l'idée abstraite : ex: « La scène de la remise » $\rightarrow$ « La mascarade protocolaire »).
    \end{enumerate}
    \textbf{Reflexe Bac :} Avoir une « fiche personnages/thèmes » prête à l'emploi pour ne pas relire le livre pendant l'épreuve.
\end{methode}

\begin{exoresolu}[De la scène narrative à l'idée essentielle (Lecture n°5)]
    \textbf{Énoncé :} À partir du passage de la cérémonie (chapitre 6 env.), formulez l'idée essentielle en une phrase argumentative (thèse + argument) en éliminant les détails anecdotiques.
    \begin{quote}
        \textit{« Le Commandant prit la parole. Il parla de la France, de la civilisation, de l'honneur. Meka ne comprenait pas grand-chose. Il voyait des bouches bouger. Soudain, on lui épingla la médaille. Il eut chaud au cou. Le Gouverneur lui serra la main : "Brave homme, Meka." Meka ne sut quoi répondre. Il pensa à sa femme, à sa terre, à son fils mort. »}
    \end{quote}
    \tcblower
    \textbf{Solution détaillée :}
    \begin{enumerate}
        \item \textbf{Étape 1 : Segmenter (Manifeste).}
            \begin{itemize}
                \item Discours officiel : France, civilisation, honneur (rhétorique vide).
                \item Incompréhension de Meka : « bouches bouger », « ne comprenait pas ».
                \item Acte protocolaire : Épinglage, serrement de main, formule toute faite (« Brave homme »).
                \item Réaction interne : Chaleur physique (gêne), silence, pensées réelles (femme, terre, fils mort).
            \end{itemize}
        \item \textbf{Étape 2 : Identifier le Latent (Argumentatif).}
            \begin{itemize}
                \item Opposition : Langue de bois coloniale / Réalité vécue du paysan.
                \item La médaille = symbole vide (ne répare pas la terre perdue, le fils mort).
                \item Meka = figure de l'exploité qui perçoit le vide du symbole.
            \end{itemize}
        \item \textbf{Étape 3 : Reformulation (Idée essentielle).}
            \textit{Ébauche :} Meka ne comprend pas le discours et pense à ses malheurs pendant qu'on lui donne une médaille.
            \textit{Correction (niveau Bac) :} \textbf{La cérémonie officielle, par son langage abstrait et ses gestes codifiés, met en lumière le fossé infranchissable entre la rhétorique coloniale de la "reconnaissance" et la réalité concrète de la spoliation et du deuil vécus par le paysan.}
        \item \textbf{Vérification :} Thèse (fossé rhétorique/réalité) + Argument (langage abstrait vs pensées concrètes) + Vocabulaire précis (rhétorique, spoliation, fossé).
    \end{enumerate}
    \textbf{Conclusion :} La contraction transforme le récit (chronologie) en analyse (logique argumentative).
\end{exoresolu}

\courssub{Les techniques de réduction du texte}

\begin{methode}[Techniques de réduction (Contraction) : Supprimer, Regrouper, Généraliser]
    \textbf{Objectif :} Réduire un texte à 25\% (ou 1/4) de sa longueur initiale sans perdre les idées forces.
    \begin{enumerate}
        \item \textbf{Suppression (Élagage) :}
            \begin{itemize}
                \item Supprimer les \textbf{exemples}, \textbf{illustrations}, \textbf{détails descriptifs}, \textbf{répétitions}, \textbf{incises} (ex: \textit{comme on dit}, \textit{selon moi}).
                \item Supprimer les \textbf{connecteurs de remplissage} (ex: \textit{enfin}, \textit{bref}, \textit{du reste}) si la logique reste claire.
            \end{itemize}
        \item \textbf{Regroupement / Fusion (Condensation syntaxique) :}
            \begin{itemize}
                \item Transformer \textbf{paragraphe} $\rightarrow$ \textbf{phrase} ; \textbf{phrase} $\rightarrow$ \textbf{groupe nominal} ; \textbf{proposition} $\rightarrow$ \textbf{participes / infinitifs / adjectifs}.
                \item Ex: « Meka est vieux. Il est fatigué. Il attend. » $\rightarrow$ \cle{Vieux et fatigué, Meka attend.}
            \end{itemize}
        \item \textbf{Généralisation / Abstraction (Montée en généralité) :}
            \begin{itemize}
                \item Remplacer une \textbf{énumération} par un \textbf{terme hyperonyme} (général).
                \item Ex: « Il a perdu sa terre, son fils, sa santé. » $\rightarrow$ \cle{Il a tout perdu.} / \cle{Il est dépossédé de l'essentiel.}
                \item Remplacer une \textbf{scène narrative} par un \textbf{concept} : « Le gouverneur lui serre la main, lui donne une médaille, fait un discours » $\rightarrow$ \cle{La mascarade protocolaire.}
            \end{itemize}
        \item \textbf{Reformulation (Réécriture) :} Changer le vocabulaire (synonymes, antonymes, changement de catégorie grammaticale) pour gagner en concision et éviter le plagiat (copier-coller).
        \item \textbf{Contrôle final :} Compter les mots (cible $\pm$ 10\%). Vérifier : Thèse présente ? Arguments principaux ? Logique respectée ? Français correct ?
    \end{enumerate}
    \textbf{Formule magique :} \cle{Supprimer l'accessoire $\rightarrow$ Regrouper le dispersé $\rightarrow$ Généraliser le concret $\rightarrow$ Reformuler le lourd.}
\end{methode}

\begin{exoresolu}[Application des techniques de réduction]
    \textbf{Énoncé :} Contractez le texte suivant à \textbf{environ 40 mots} (texte source $\approx$ 160 mots). Indiquez les opérations effectuées.
    \begin{quote}
        \textit{« Meka s'était levé très tôt ce matin-là. Il avait mis son plus beau pagne, celui qu'il ne portait que les grands jours. Il avait marché pendant deux heures sous le soleil brûlant pour arriver à la résidence du Commandant. Ses pieds étaient tout gonflés, ses jambes tremblaient de fatigue. Quand il arriva, il y avait déjà beaucoup de monde. Tout le monde poussait, on se bousculait. Le soleil tapait fort sur la place. Meka avait soif, il avait faim. Mais il ne voulait pas partir. Il voulait voir la médaille. Enfin, le Commandant arriva. Il fit un long discours que personne n'écoutait vraiment. Puis il appela Meka. Meka avança, la tête basse. On lui mit la médaille au cou. C'était fini. Il fit demi-tour et rentra chez lui, plus fatigué encore, le ventre toujours vide. »}
    \end{quote}
    \tcblower
    \textbf{Solution détaillée :}
    \begin{enumerate}
        \item \textbf{Analyse source (163 mots) :} Cible $\approx$ 40 mots ($\pm$ 4).
        \item \textbf{Opérations de réduction :}
            \begin{itemize}
                \item \textbf{Supprimés :} Détails vestimentaires (pagne), durée marche (2h), sensations physiques précises (pieds gonflés, jambes tremblent, soif, faim), description foule (poussées, soleil), attitude passive (tête basse), retour détaillé.
                \item \textbf{Regroupés :} Préparation + Marche + Attente $\rightarrow$ \textit{Après un pénible trajet}. Discours + Remise $\rightarrow$ \textit{La cérémonie protocolaire}.
                \item \textbf{Généralisés :} « Pagne, marche, fatigue, faim, soif » $\rightarrow$ \textit{Dans la misère}. « Médaille au cou, rentre vide » $\rightarrow$ \textit{Reçoit une distinction vaine}.
            \end{itemize}
        \item \textbf{Proposition de contraction (39 mots) :}
            \begin{quote}
                \textit{« Après un pénible trajet dans la misère, Meka endure une cérémonie protocolaire où un long discours précède la remise d'une médaille. Il repart, exténué, toujours affamé, conscient du dérisoire de cet honneur vide. »}
            \end{quote}
        \item \textbf{Vérification :} Thèse (honneur vide/misère) conservée. Arguments (pénible trajet, discours long, médaille, faim) conservés. Syntaxe fluide (participes, groupes nominaux). Compte : 39 mots. \textbf{Validé.}
    \end{enumerate}
\end{exoresolu}

\courssub{Les outils de liaison : conjonctions de coordination et de subordination}

\begin{methode}[Maîtriser les connecteurs logiques pour la cohérence du résumé]
    \textbf{Objectif :} Utiliser les conjonctions pour articuler logiquement les idées essentielles dans la contraction (cause, conséquence, opposition, concession, but).
    \begin{enumerate}
        \item \textbf{Coordination (Liaison entre propositions indépendantes / groupes de même fonction) :}
            \begin{center}
                \begin{tabular}{|c|l|l|}\hline
                \textbf{Relation} & \textbf{Conjonctions} & \textbf{Exemple contraction} \\ \hline
                Addition / Énumération & \cle{et, ni, ou, puis, aussi} & \textit{Meka est vieux \textbf{et} fatigué.} \\ \hline
                Opposition / Contraste & \cle{mais, or, car (rare)} & \textit{Il reçoit une médaille \textbf{mais} a faim.} \\ \hline
                Choix / Alternative & \cle{ou, ou bien, soit...soit} & \textit{Instruction \textbf{ou} éducation ?} \\ \hline
                \end{tabular}
            \end{center}
        \item \textbf{Subordination (Liaison proposition principale $\rightarrow$ proposition subordonnée) :}
            \begin{center}
                \begin{tabular}{|c|l|l|}\hline
                \textbf{Relation logique} & \textbf{Conjonctions / Locutions} & \textbf{Exemple contraction} \\ \hline
                Cause & \cle{parce que, comme, puisque, vu que} & \textit{\textbf{Comme} il a faim, la médaille est vaine.} \\ \hline
                Conséquence & \cle{si bien que, de sorte que, au point que} & \textit{Il est vieux, \textbf{si bien que} il tremble.} \\ \hline
                But & \cle{afin que, pour que, de peur que} & \textit{Il marche \textbf{afin de} recevoir la médaille.} \\ \hline
                Concession & \cle{bien que, quoique, encore que + subj.} & \textit{\textbf{Bien qu'}il soit décoré, il reste pauvre.} \\ \hline
                Condition & \cle{si, pourvu que, à condition que} & \textit{\textbf{Si} l'école n'éduque pas, elle échoue.} \\ \hline
                Temps / Antériorité & \cle{lorsque, quand, dès que, après que + ind.} & \textit{\textbf{Après que} la cérémonie finit, il part.} \\ \hline
                \end{tabular}
            \end{center}
        \item \textbf{Stratégie contraction :} Privilégier la \textbf{subordination} (plus compacte) à la coordination ou aux phrases juxtaposées.
            \begin{itemize}
                \item \textit{Lourd :} « Meka a faim. Donc la médaille ne sert à rien. » (2 phrases, coordination/adverbe).
                \item \textit{Compact :} \cle{La médaille ne sert à rien \textbf{puisque} Meka a faim.} (1 phrase, subordination causale).
            \end{itemize}
        \item \textbf{Attention au mode :} Concession (\cle{bien que}, \cle{quoique}) $\rightarrow$ \textbf{Subjonctif}. Cause/Conséquence/Temps (indicateurs) $\rightarrow$ \textbf{Indicatif} (sauf \cle{avant que}, \cle{pour que} $\rightarrow$ Subj).
    \end{enumerate}
\end{methode}

\begin{exoresolu}[Réécriture connectée pour la concision]
    \textbf{Énoncé :} Reliez les idées suivantes en \textbf{une seule phrase complexe} (max 30 mots) en utilisant au moins une subordination de cause, une de concession et une de conséquence.
    \textit{Idées :} 1. Meka est vieux. 2. Il marche longtemps. 3. Il reçoit une médaille. 4. La médaille ne le nourrit pas. 5. Il reste digne.
    \tcblower
    \textbf{Solution détaillée :}
    \begin{enumerate}
        \item \textbf{Choix des connecteurs :}
            \begin{itemize}
                \item Cause : \cle{Parce que} / \cle{Comme} / \cle{Puisque} (pour marcher longtemps $\leftarrow$ vieux ? Non, marcher longtemps $\leftarrow$ veut médaille. Ou médaille vaine $\leftarrow$ ne nourrit pas).
                \item Concession : \cle{Bien que} / \cle{Quoique} (Vieux $\rightarrow$ marche ; Médaille $\rightarrow$ pas nourri ; Digne $\rightarrow$ misère).
                \item Conséquence : \cle{Si bien que} / \cle{De sorte que}.
            \end{itemize}
        \item \textbf{Construction pas à pas :}
            \begin{align*}
                &\text{Base : Meka reçoit une médaille.} \\
                &\text{Concession (âge/effort) : } \textbf{Bien que vieux,} \text{ il marche longtemps pour la recevoir.} \\
                &\text{Cause (inutilité) : } \textbf{Mais comme} \text{ cette médaille ne le nourrit pas,} \\
                &\text{Conséquence (dignité) : } \textbf{il reste digne, si bien que} \text{ l'honneur officiel devient dérisoire.}
            \end{align*}
        \item \textbf{Phrase finale (26 mots) :}
            \begin{quote}
                \textit{« \textbf{Bien que vieux}, Meka marche longtemps pour recevoir une médaille \textbf{qui, ne le nourrissant pas, devient dérisoire} ; \textbf{pourtant}, il conserve sa dignité. »}
            \end{quote}
            \textit{(Note : utilisation du gérondif "ne le nourrissant pas" = cause + compacité ; "pourtant" = coordination adverse pour la dignité).}
        \item \textbf{Variante plus subordonnée (24 mots) :}
            \textit{« \textbf{Quoique vieux}, Meka marche pour une médaille \textbf{vaine puisqu'elle ne le nourrit pas}, \textbf{mais il garde sa dignité.} »}
        \item \textbf{Vérification :} 1 phrase. Relations logiques claires (Concession, Cause, Opposition/Concession). Vocabulaire précis (vaine, dérisoire). Subjonctif après \cle{quoique} (\textit{quoique vieux} $\leftarrow$ ellipsis de \textit{soit}).
    \end{enumerate}
    \textbf{Conclusion :} La subordination permet de densifier l'information : 5 idées $\rightarrow$ 1 phrase.
\end{exoresolu}

\courssub{La reformulation des idées essentielles}

\begin{methode}[Réformuler sans trahir : Paraphrase et Transposition]
    \textbf{Objectif :} Exprimer les idées du texte source avec ses propres mots (vocabulaire, syntaxe) pour prouver la compréhension et éviter le plagiat.
    \begin{enumerate}
        \item \textbf{Changement de catégorie grammaticale (Transposition) :}
            \begin{itemize}
                \item Verbe $\rightarrow$ Nom (Nominalisation) : \textit{Il exploite} $\rightarrow$ \cle{Son exploitation}.
                \item Adjectif $\rightarrow$ Nom : \textit{Il est cruel} $\rightarrow$ \cle{Sa cruauté}.
                \item Proposition $\rightarrow$ Groupe Prépositionnel / Participial : \textit{Parce qu'il est pauvre} $\rightarrow$ \cle{De par sa pauvreté} / \cle{Pauvre}.
            \end{itemize}
        \item \textbf{Substitution lexicale (Synonymie / Hyperonymie / Périphrase) :}
            \begin{itemize}
                \item \textit{Donner une médaille} $\rightarrow$ \cle{Décorer}, \cle{Récompenser}, \cle{Distinction honorifique}.
                \item \textit{Ventre vide} $\rightarrow$ \cle{Faim}, \cle{Dénutrition}, \cle{Misère}, \cle{Inanition}.
                \item \textit{Vieux Nègre} $\rightarrow$ \cle{Le paysan âgé}, \cle{Le protagoniste}, \cle{La victime coloniale}.
            \end{itemize}
        \item \textbf{Changement de voix / perspective :}
            \begin{itemize}
                \item Actif $\rightarrow$ Passif (si agent important) : \textit{Le Gouverneur remet la médaille} $\rightarrow$ \cle{La médaille est remise par le Gouverneur}.
                \item Discours Direct $\rightarrow$ Indirect : \textit{« Brave homme »} $\rightarrow$ \cle{Il le qualifie de brave homme}.
            \end{itemize}
        \item \textbf{Règle d'or :} \cle{Sens identique + Forme nouvelle}. Ne pas ajouter d'idée personnelle. Garder la neutralité (sauf si le texte est subjectif, garder la subjectivité de l'auteur).
        \item \textbf{Entraînement :} Prendre une phrase du cours, la reformuler 3 façons différentes sur brouillon.
    \end{enumerate}
\end{methode}

\begin{exoresolu}[Exercices de reformulation ciblée]
    \textbf{Énoncé :} Reformulez les segments suivants en appliquant la technique indiquée entre parenthèses. La phrase doit être plus concise.
    \begin{enumerate}
        \item \textit{« Meka a marché pendant deux heures sous le soleil. »} (Nominalisation + Suppression détail)
        \item \textit{« Le Commandant a fait un discours très long que personne n'écoutait. »} (Passif + Adjectif)
        \item \textit{« Il a reçu une médaille mais il n'avait rien à manger. »} (Concession subordonnée + Vocabulaire abstrait)
        \item \textit{« Kelara a dit à son mari que cette médaille ne servait à rien. »} (Discours indirect + Verbe de parole précis)
    \end{enumerate}
    \tcblower
    \textbf{Solution détaillée :}
    \begin{enumerate}
        \item \textbf{Nominalisation :} \cle{Une marche de deux heures sous le soleil.} (Ou plus concis : \cle{Une longue marche exténuante.}) — \textit{Gain : Verbe $\rightarrow$ Nom, suppression "pendant".}
        \item \textbf{Passif + Adjectif :} \cle{Un long discours inaudible fut prononcé par le Commandant.} (Ou : \cle{Le discours interminable du Commandant fut ignoré.}) — \textit{Gain : "que personne n'écoutait" $\rightarrow$ "inaudible/ignoré".}
        \item \textbf{Concession + Abstrait :} \cle{Bien que décoré, il meurt de faim.} (Ou : \cle{La distinction honorifique ne saurait combler sa faim.}) — \textit{Gain: "reçu une médaille" $\rightarrow$ "décoré/distinction"; "rien à manger" $\rightarrow$ "faim". Structure concessive compacte.}
        \item \textbf{Discours indirect + Verbe précis :} \cle{Kelara dénonce l'inutilité de la médaille auprès de son mari.} (Ou : \cle{Kelara persuade son mari de la vanité de cette distinction.}) — \textit{Gain: "a dit que... ne servait à rien" $\rightarrow$ "dénonce l'inutilité / persuade de la vanité". Verbe de parole à valeur argumentative.}
    \end{enumerate}
    \textbf{Conclusion :} La reformulation est le moteur de la contraction : elle permet de dire « plus avec moins ».
\end{exoresolu}

\courssub{Compte rendu et remédiation}

\begin{methode}[Auto-évaluation et correction type Bac (Grille de relecture)]
    \textbf{Objectif :} Vérifier sa production (contraction/résumé) avant de rendre la copie en appliquant les critères officiels.
    \begin{enumerate}
        \item \textbf{Critère 1 : Respect de la consigne (Longueur).} Compter les mots. \textbf{Tolérance : $\pm$ 10\% de la cible.} (Ex: Cible 100 mots $\rightarrow$ entre 90 et 110). \cle{Pénalité :} Hors fourchette = note plafonnée.
        \item \textbf{Critère 2 : Fidélité aux idées essentielles (Contenu).} \checkmark Thèse présente ? \checkmark Arguments principaux (2 ou 3) présents ? \checkmark Exemples/détails \textbf{exclus} ? \checkmark Pas d'invention, pas d'interprétation personnelle ?
        \item \textbf{Critère 3 : Cohérence et Cohésion (Forme).} \checkmark Connecteurs logiques variés et justes (cause, conséq, oppos, concess) ? \checkmark Pas de phrases hachées (style télégraphique) ? \checkmark Références anaphoriques claires (qui fait quoi ?) ?
        \item \textbf{Critère 4 : Qualité de la langue (Orthographe/Grammaire/Syntaxe).} \checkmark Accords (sujet/verbe, GN) ? \checkmark Temps (cohérence récit/analyse) ? \checkmark Ponctuation ? \checkmark Vocabulaire précis (pas de répétitions maladroites) ?
        \item \textbf{Critère 5 : Présentation.} \checkmark Paragraphe unique (sauf consigne contraire) ? \checkmark Nombre de mots indiqué en fin de texte ? \checkmark Lisible ?
        \item \textbf{Remédiation type :}
            \begin{itemize}
                \item \textit{Trop long ?} $\rightarrow$ Couper les adjectifs, fusionner phrases, généraliser exemples.
                \item \textit{Trop court ?} $\rightarrow$ Rétablir un argument manquant, expliciter un lien logique.
                \item \textit{Hors sujet ?} $\rightarrow$ Re-lire le texte source, identifier la thèse.
                \item \textit{Français approximatif ?} $\rightarrow$ Simplifier la syntaxe (Sujet-Verbe-Complément), vérifier accords de base.
            \end{itemize}
    \end{enumerate}
    \textbf{Reflexe final :} \cle{Relire à voix basse (mentalement) pour "entendre" les fautes de syntaxe.}
\end{methode}

\begin{exoresolu}[Auto-correction d'une contraction d'élève (Simulation Bac)]
    \textbf{Énoncé :} Voici une proposition d'élève pour la contraction du texte sur Meka (cible 40 mots). Évaluez-la selon la grille (Longueur, Contenu, Forme, Langue) et proposez une version corrigée.
    \textit{Proposition élève (52 mots) :}
    \textit{« Meka est un vieux paysan qui marche deux heures sous le soleil pour aller chercher sa médaille. Il a très faim et soif. Le Commandant fait un discours très long. Après il donne la médaille à Meka. Meka est content mais il a encore faim. Il rentre chez lui triste. C'est injuste. »}
    \tcblower
    \textbf{Solution détaillée :}
    \begin{enumerate}
        \item \textbf{Évaluation (Grille) :}
            \begin{itemize}
                \item \textbf{Longueur :} 52 mots \textbf{(Hors fourchette 36-44)}. \textit{Pénalité majeure.}
                \item \textbf{Contenu :} Détails inutiles (2h, soif, content/triste). Idée essentielle (injustice/vide symbole) mal formulée ("C'est injuste" = jugement élève). Manque : opposition protocole/réalité.
                \item \textbf{Forme :} Phrases courtes, juxtaposées, style "récit d'enfant". Pas de subordination. Répétitions (Meka x3, médaille x2).
                \item \textbf{Langue :} Correct mais familier ("chercher sa médaille", "très long", "très faim").
            \end{itemize}
            \textbf{Note estimée : 5/20} (Hors format + contenu faible + style faible).
        \item \textbf{Remédiation (Réécriture ciblée 39 mots) :}
            \begin{quote}
                \textit{« \textbf{Après une marche exténuante}, le vieux paysan Meka \textbf{subit une cérémonie protocolaire} où un discours creux précède la remise d'une médaille. \textbf{Vaine distinction :} il rentre \textbf{affamé}, conscient du mépris colonial masqué par les honneurs. »}
            \end{quote}
        \item \textbf{Analyse des corrections :}
            \begin{itemize}
                \item \textbf{Coupe :} "deux heures", "soif", "content", "triste", "C'est injuste".
                \item \textbf{Nominalisation :} "marche deux heures" $\rightarrow$ "marche exténuante" ; "Commandant fait discours" $\rightarrow$ "cérémonie protocolaire / discours creux".
                \item \textbf{Subordination/Concession :} "Il a faim / Il reçoit" $\rightarrow$ "Vaine distinction : il rentre affamé" (Ellipse causative).
                \item \textbf{Latent restitué :} "mépris colonial masqué par les honneurs" (Remplace "C'est injuste" par analyse textuelle).
            \end{itemize}
    \end{enumerate}
    \textbf{Conclusion :} La remédiation transforme un récit subjectif et long en une analyse objective, dense et argumentative.
\end{exoresolu}

\begin{attention}[Les 5 erreurs qui coûtent des points au Bac]
    \begin{enumerate}
        \item \textbf{Ne pas compter les mots (Hors fourchette) :} C'est l'erreur n°1. Une contraction à 60 mots pour une cible de 40 = \textbf{note plafonnée à 8/20} souvent, même si le contenu est bon. \cle{Comptez et indiquez le nombre de mots en fin de copie.}
        \item \textbf{Conserver les exemples et détails anecdotiques :} Raconter l'histoire (Meka met son pagne, marche 2h, a soif) au lieu d'analyser la structure (Marche exténuante $\rightarrow$ Cérémonie $\rightarrow$ Ironie). \cle{Supprimez le "quoi" (anecdotique) pour garder le "pourquoi" (argumentatif).}
        \item \textbf{Copier-coller des phrases du texte (Plagiat / Absence de reformulation) :} Le correcteur cherche \textbf{votre} langue. Recopier des segments entiers prouve que vous n'avez pas synthétisé. \cle{Reformulez systématiquement : nominalisation, synonymes, changement de voix.}
        \item \textbf{Oublier la logique argumentative (Connecteurs) :} Produire une liste de phrases juxtaposées : « Meka marche. Il arrive. On lui donne médaille. Il a faim. » $\rightarrow$ Perte de points \textbf{Cohérence}. \cle{Utilisez \textbf{parce que, bien que, si bien que, alors que} pour lier les idées.}
        \item \textbf{Injecter son opinion personnelle :} Écrire « C'est scandaleux », « Je trouve que c'est triste », « Le Gouverneur est méchant ». \cle{Le résumé/contraction est objectif : il rend compte du texte, il ne le juge pas.} Utilisez : « Le texte montre que... », « L'auteur dénonce... », « La satire révèle... ».
    \end{enumerate}
\end{attention}

\newpage

\coursec{Exercices}

\courssub{Je vérifie mes connaissances}

\begin{exercice}[QCM : le texte argumentatif et la contraction]
Pour chaque question, entoure la bonne réponse.
\begin{enumerate}
\item La thèse d'un texte argumentatif est :
\begin{enumerate}
\item un exemple concret qui illustre une idée ;
\item l'idée principale que l'auteur défend ;
\item une anecdote destinée à amuser le lecteur.
\end{enumerate}
\item Dans un résumé, les exemples doivent être :
\begin{enumerate}
\item recopiés intégralement ;
\item supprimés ou remplacés par un terme générique ;
\item développés avec des détails personnels.
\end{enumerate}
\item Un texte de 560 mots résumé au quart doit contenir environ :
\begin{enumerate}
\item 560 mots ; \item 280 mots ; \item 140 mots.
\end{enumerate}
\item La conjonction \emph{parce que} exprime :
\begin{enumerate}
\item la concession ; \item la cause ; \item le but.
\end{enumerate}
\end{enumerate}
\end{exercice}

\begin{exercice}[Vrai ou faux ? Justifie ta réponse]
Indique si chaque affirmation est vraie ou fausse, puis justifie en une ou deux phrases.
\begin{enumerate}
\item Le contenu latent d'un texte est ce que l'auteur dit explicitement, mot pour mot.
\item Dans \emph{Le Vieux Nègre et la médaille}, le discours de remise de la médaille est un éloge sincère de Meka.
\item Une reformulation réussie ne reprend aucun segment de plus de quatre mots du texte original.
\item Les conjonctions de coordination relient deux propositions dont l'une dépend de l'autre.
\end{enumerate}
\end{exercice}

\begin{exercice}[Définitions]
Définis en deux ou trois phrases chacun des termes suivants, en donnant pour chacun un exemple :
\begin{enumerate}
\item la thèse ;
\item l'argument ;
\item le contenu latent ;
\item la généralisation (dans la technique de réduction).
\end{enumerate}
\end{exercice}

\begin{exercice}[Classement : coordination ou subordination ?]
Classe les conjonctions et locutions conjonctives suivantes dans un tableau à deux colonnes : \textbf{conjonctions de coordination} / \textbf{conjonctions de subordination}. Précise ensuite la relation logique exprimée par chacune (cause, but, temps, concession, opposition, addition...).
\begin{center}
mais — parce que — donc — bien que — quand — car — afin que — ou — puisque — cependant (attention au piège !)
\end{center}
\end{exercice}

\courssub{Je m'entraîne}

\begin{exercice}[Repérage dans un éditorial]
Recopie le texte suivant sur ta feuille, puis : \textbf{(a)} souligne la thèse ; \textbf{(b)} encadre trois arguments ; \textbf{(c)} barre les exemples.

\emph{« L'éducation des jeunes filles est la clé du développement de notre pays. D'abord, une femme instruite transmet le savoir à ses enfants : à Maroua, les fillettes dont la mère a fréquenté l'école réussissent mieux au primaire. Ensuite, scolariser les filles retarde les mariages précoces, comme on l'a observé dans plusieurs villages de la région du Centre. Enfin, une jeune fille éduquée devient une travailleuse capable de faire vivre sa famille : pensons à ces institutrices de Bafoussam qui soutiennent leurs parents âgés. Refuser l'école aux filles, c'est donc refuser l'avenir à tout le Cameroun. »}
\end{exercice}

\begin{exercice}[Fusion de phrases]
Transforme chaque couple de phrases simples en une phrase complexe, en utilisant la conjonction de subordination indiquée entre parenthèses.
\begin{enumerate}
\item Meka a reçu une médaille. Il a servi fidèlement l'administration coloniale. (parce que)
\item Les villageois admirent Meka. Ils ignorent la véritable signification de cette récompense. (alors que)
\item Meka donnera sa terre aux Blancs. Il obtiendra leur reconnaissance. (afin que)
\item Le vieil homme est fier. Il sent une gêne au fond de lui-même. (bien que)
\item La cérémonie s'achèvera. Meka rentrera seul dans sa case. (quand)
\item Kelara aime son mari. Elle désapprouve sa naïveté. (quoique)
\item Les Blancs parlent d'amitié. Ils méprisent les Africains. (tandis que)
\item Meka a perdu ses terres. Il a cru aux promesses des missionnaires. (puisque)
\end{enumerate}
\end{exercice}

\begin{exercice}[Reformulation sans emprunt]
Réécris chacune des cinq phrases suivantes (inspirées du \emph{Vieux Nègre et la médaille}) en conservant le sens, mais sans reprendre aucun segment de plus de quatre mots de la phrase originale.
\begin{enumerate}
\item « Le vieux Meka attendait la médaille que les Blancs devaient lui remettre. »
\item « Il avait donné ses terres à la mission et ses fils à la guerre. »
\item « Les villageois venaient de loin pour assister à la grande cérémonie. »
\item « Le commandant prononça un long discours plein de belles promesses. »
\item « Cette médaille brillante ne pouvait pas rendre à Meka ce qu'il avait perdu. »
\end{enumerate}
\end{exercice}

\begin{exercice}[Généralisation des énumérations]
Remplace chaque énumération concrète par un terme générique adapté, puis justifie ton choix en une phrase.
\begin{enumerate}
\item « Les manguiers, les goyaviers et les papayers entouraient la case. »
\item « Il vendit ses chèvres, ses poules et ses moutons pour payer la fête. »
\item « Les houes, les machettes et les paniers étaient rangés derrière la porte. »
\item « Le commandant, le missionnaire et le docteur prirent place à la tribune. »
\item « Elle prépara le ndolé, le miondo et le poulet DG pour les invités. »
\item « Les taxis, les motos et les bus encombraient les rues de Douala. »
\item « Il lut le journal, écouta la radio et regarda la télévision. »
\item « Les cahiers, les stylos et les règles manquaient dans le cartable. »
\item « Sa sœur, son frère et ses cousins vinrent le féliciter. »
\item « La pluie, le vent et la chaleur épuisèrent les marcheurs. »
\end{enumerate}
\end{exercice}

\courssub{Je me prépare au Bac}

\begin{exercice}[Contraction de texte — type Baccalauréat technique]
Lis attentivement le texte suivant (560 mots), puis réponds aux questions.

\begin{itshape}«ans nombre de nos villages, on entend encore qu'il est inutile d'envoyer une fille à l'école : tôt ou tard, dira-t-on, elle quittera le foyer pour se marier, et l'argent investi dans ses études profitera à une autre famille. Ce raisonnement, en apparence économique, est en réalité un piège qui appauvrit toute la nation. L'éducation des jeunes filles est en effet la condition première du développement du Cameroun.

D'abord, une mère instruite élève des enfants en meilleure santé et plus savants. Les statistiques le montrent : dans les régions où la scolarisation des filles progresse, la mortalité infantile recule et la réussite scolaire s'améliore. Une femme qui sait lire comprend les ordonnances du médecin, suit le calendrier des vaccinations et aide ses enfants à faire leurs devoirs. Pensons à ces mères de l'Extrême-Nord qui, grâce aux cours d'alphabétisation, vérifient désormais les bulletins de leurs fils. Instruire une fille, c'est donc instruire toute une génération.

Ensuite, l'école protège la jeune fille elle-même. Tant qu'elle reste sur les bancs de la classe, elle échappe au mariage précoce et à ses drames : grossesses dangereuses, dépendance totale, renoncement à tout projet personnel. Chaque année de scolarité gagnée recule l'âge du mariage et ouvre un horizon de choix. La jeune fille instruite peut devenir institutrice, infirmière, agronome ou commerçante avisée ; elle connaît ses droits et sait les défendre.

Enfin, le pays tout entier y gagne. Une nation qui prive la moitié de ses enfants du savoir se prive de la moitié de ses talents. Nos entreprises, nos hôpitaux, nos administrations ont besoin de toutes les intelligences. Le contre-exemple des pays qui ont massivement scolarisé leurs filles est éclatant : leur économie s'est transformée en une génération.

Certes, les frais de scolarité pèsent sur les familles pauvres, et les traditions ont leur valeur. Mais aucune tradition digne de ce nom ne peut exiger l'ignorance, et aucun calcul à courte vue ne doit sacrifier l'avenir. Offrons donc à chaque fille du Cameroun, de Waza à Kribi, la même chance qu'à son frère : celle d'apprendre, de choisir et de construire. »\end{itshape}

\begin{enumerate}
\item Dégage la thèse de l'auteur en une phrase.
\item Releve les trois arguments principaux et, pour chacun, l'exemple qui l'illustre.
\item Résume ce texte au quart de sa longueur (environ 140 mots, marge de plus ou moins 10 \%). Tu supprimeras les exemples, tu généraliseras les énumérations et tu reformuleras les idées essentielles sans reprendre de segment de plus de quatre mots.
\item Indique le nombre exact de mots de ton résumé à la fin de ta copie.
\end{enumerate}
\end{exercice}

\begin{exercice}[Contenu latent et ironie — \emph{Le Vieux Nègre et la médaille}]
Dans le roman de Ferdinand Oyono, le commandant remet solennellement une médaille à Meka devant tout le village, en le félicitant d'avoir donné ses terres à la mission et ses deux fils à la guerre, morts « pour la France ».
\begin{enumerate}
\item Rappelle ce qu'est le contenu manifeste et le contenu latent d'un discours.
\item Releve dans ta mémoire du roman (ou dans l'extrait étudié en classe) deux procédés ironiques utilisés par Oyono dans la scène de remise de la médaille.
\item Rédige un paragraphe d'une dizaine de lignes expliquant ce que l'auteur dénonce réellement derrière l'éloge apparent : tu montreras comment l'ironie transforme une cérémonie d'honneur en condamnation du système colonial.
\item Justifie ta démarche : pourquoi Oyono choisit-il l'ironie plutôt que la critique directe ?
\end{enumerate}
\end{exercice}

\begin{exercice}[Épreuve complète : réception et production]
\textbf{Texte support.} \emph{« Meka avait attendu ce jour toute sa vie. La médaille brillait sur sa poitrine, les Blancs applaudissaient, le village l'admirait. Pourtant, quand la nuit tomba et que les invités furent partis, le vieil homme resta seul devant sa case, la médaille à la main, et il pensa à ses deux fils qui ne reviendraient jamais, à ses terres qui ne lui appartenaient plus. »}

\begin{enumerate}
\item \textbf{Compréhension.} Quel contraste l'extrait met-il en valeur ? Releve deux champs lexicaux opposés.
\item \textbf{Langue.} Dans l'extrait, repère : (a) une conjonction de coordination et précise sa valeur ; (b) deux conjonctions de subordination et précise la relation logique de chacune (temps, cause, concession...).
\item \textbf{Réécriture.} Transforme les trois premières phrases de l'extrait en un seul énoncé complexe comportant au moins une subordonnée de temps et une subordonnée de concession.
\item \textbf{Résumé.} Résume l'extrait en 25 à 30 mots, sans aucune reprise de segment de plus de quatre mots.
\item \textbf{Expression personnelle.} En dix lignes maximum, réponds à la question : une récompense peut-elle compenser une perte irréparable ? Appuie ta réponse sur un exemple tiré du roman ou de la vie courante au Cameroun.
\end{enumerate}
\end{exercice}

\newpage

\coursec{Corrigés des exercices}

\begin{corrige}[Exercice 1 — QCM : le texte argumentatif et la contraction]
\begin{enumerate}
\item Réponse \textbf{b} : la thèse est \cle{l'idée principale que l'auteur défend}. Un exemple (réponse a) ne fait qu'illustrer un argument ; une anecdote n'est jamais la thèse.
\item Réponse \textbf{b} : dans un résumé, les exemples sont \cle{supprimés} ou remplacés par un \cle{terme générique} (technique de la généralisation). Recopier les exemples (a) alourdit le résumé ; ajouter des détails personnels (c) est une faute grave : le résumé ne doit contenir aucune idée étrangère au texte.
\item Réponse \textbf{c} : $560 \div 4 = 140$ mots. Le quart de 560 mots est bien 140 mots (avec une marge tolérée de plus ou moins 10 \%, soit entre 126 et 154 mots).
\item Réponse \textbf{b} : \emph{parce que} introduit une subordonnée de \cle{cause}. La concession s'exprime par \emph{bien que}, \emph{quoique} ; le but par \emph{pour que}, \emph{afin que}.
\end{enumerate}
\end{corrige}

\begin{corrige}[Exercice 2 — Vrai ou faux ?]
\begin{enumerate}
\item \textbf{Faux.} Le contenu latent est ce que l'auteur \emph{suggère} sans le dire explicitement (sous-entendus, ironie, implicite). Ce qui est dit mot pour mot relève du \cle{contenu manifeste}.
\item \textbf{Faux.} Le discours de remise de la médaille est un éloge \emph{en apparence} seulement : Oyono y déploie une ironie mordante. Féliciter un homme d'avoir perdu ses terres et ses deux fils est en réalité une dénonciation de l'hypocrisie coloniale.
\item \textbf{Vrai.} C'est la règle pratique de la reformulation : au-delà de quatre mots consécutifs repris au texte, on parle de recopie (ou de citation, qui doit alors être signalée). Une reformulation réussie change le vocabulaire et la construction de la phrase tout en conservant le sens.
\item \textbf{Faux.} C'est l'inverse : les conjonctions de coordination (\emph{mais, ou, et, donc, or, ni, car}) relient deux éléments de \emph{même niveau}, autonomes. Ce sont les conjonctions de subordination qui introduisent une proposition dépendant d'une autre.
\end{enumerate}
\end{corrige}

\begin{corrige}[Exercice 3 — Définitions]
\begin{enumerate}
\item \textbf{La thèse} : c'est l'idée principale, le point de vue que l'auteur d'un texte argumentatif cherche à faire admettre au lecteur. \emph{Exemple :} « L'éducation des jeunes filles est la clé du développement du Cameroun. »
\item \textbf{L'argument} : c'est une raison avancée pour prouver ou justifier la thèse ; il est souvent illustré par un exemple. \emph{Exemple :} « Une mère instruite élève des enfants en meilleure santé » est un argument qui soutient la thèse précédente.
\item \textbf{Le contenu latent} : c'est le sens caché, non formulé, qu'un discours laisse deviner derrière le sens apparent (contenu manifeste). \emph{Exemple :} quand le commandant félicite Meka d'avoir « donné » ses fils, le contenu latent est que la colonisation lui a tout pris.
\item \textbf{La généralisation} : technique de réduction qui consiste à remplacer une énumération de termes concrets ou des exemples particuliers par un terme générique qui les englobe. \emph{Exemple :} « les manguiers, les goyaviers et les papayers » devient « les arbres fruitiers ».
\end{enumerate}
\end{corrige}

\begin{corrige}[Exercice 4 — Classement : coordination ou subordination ?]
\begin{center}
\begin{tabularx}{\linewidth}{|l|X|l|X|}
\hline
\rowcolor{popblueL} \textbf{Coordination} & \textbf{Relation logique} & \textbf{Subordination} & \textbf{Relation logique} \\
\hline
mais & opposition & parce que & cause \\
\hline
donc & conséquence & bien que & concession \\
\hline
car & cause & quand & temps \\
\hline
ou & alternative & afin que & but \\
\hline
& & puisque & cause \\
\hline
\end{tabularx}
\end{center}
\textbf{Le piège :} \emph{cependant} n'est \textbf{ni} une conjonction de coordination \textbf{ni} une conjonction de subordination : c'est un \cle{adverbe de liaison} (connecteur) exprimant l'opposition. Il ne pouvait donc figurer dans aucune des deux colonnes.
\begin{attention}[Erreurs fréquentes]
\emph{Car} exprime la cause mais reste une conjonction de \emph{coordination} : la proposition introduite par \emph{car} n'est pas subordonnée. De même, \emph{donc} (conséquence) coordonne, il ne subordonne pas.
\end{attention}
\end{corrige}

\begin{corrige}[Exercice 5 — Repérage dans un éditorial]
\textbf{(a) Thèse (à souligner) :} « L'éducation des jeunes filles est la clé du développement de notre pays » (reprise en conclusion : « Refuser l'école aux filles, c'est refuser l'avenir à tout le Cameroun »).

\textbf{(b) Arguments (à encadrer) :}
\begin{enumerate}
\item « une femme instruite transmet le savoir à ses enfants » ;
\item « scolariser les filles retarde les mariages précoces » ;
\item « une jeune fille éduquée devient une travailleuse capable de faire vivre sa famille ».
\end{enumerate}

\textbf{(c) Exemples (à barrer) :} « à Maroua, les fillettes dont la mère a fréquenté l'école réussissent mieux au primaire » ; « comme on l'a observé dans plusieurs villages de la région du Centre » ; « pensons à ces institutrices de Bafoussam qui soutiennent leurs parents âgés ».

\textbf{Remarque méthodologique :} les marqueurs \emph{d'abord, ensuite, enfin} signalent les trois arguments ; les deux-points et les formules comme \emph{pensons à} ou \emph{comme} annoncent les exemples.
\end{corrige}

\begin{corrige}[Exercice 6 — Fusion de phrases]
Corrections possibles (d'autres formulations correctes sont acceptables) :
\begin{enumerate}
\item Meka a reçu une médaille \textbf{parce qu}'il avait servi fidèlement l'administration coloniale. (subordonnée de cause)
\item Les villageois admirent Meka, \textbf{alors qu}'ils ignorent la véritable signification de cette récompense. (subordonnée d'opposition)
\item Meka donnera sa terre aux Blancs \textbf{afin qu}'il obtienne leur reconnaissance. (subordonnée de but ; attention au subjonctif : \emph{obtienne})
\item \textbf{Bien qu}'il soit fier, le vieil homme sent une gêne au fond de lui-même. (subordonnée de concession ; subjonctif : \emph{soit})
\item \textbf{Quand} la cérémonie s'achèvera, Meka rentrera seul dans sa case. (subordonnée de temps ; futur dans les deux propositions)
\item \textbf{Quoique} Kelara aime son mari, elle désapprouve sa naïveté. (subordonnée de concession ; subjonctif : \emph{aime})
\item Les Blancs parlent d'amitié, \textbf{tandis qu}'ils méprisent les Africains. (subordonnée d'opposition)
\item \textbf{Puisque} Meka a cru aux promesses des missionnaires, il a perdu ses terres. (subordonnée de cause)
\end{enumerate}
\begin{attention}[Points de vigilance]
\emph{Afin que}, \emph{bien que} et \emph{quoique} exigent le \cle{subjonctif}. Après \emph{quand} exprimant l'avenir, on emploie le futur simple (\emph{s'achèvera}), jamais le présent.
\end{attention}
\end{corrige}

\begin{corrige}[Exercice 7 — Reformulation sans emprunt]
Propositions de reformulation (toute version respectant le sens et la règle des quatre mots est valable) :
\begin{enumerate}
\item « Depuis longtemps, Meka espérait recevoir la décoration promise par les colons. »
\item « Il avait sacrifié ses terres au profit de la mission et perdu ses deux garçons à la guerre. »
\item « De tous les villages alentour, on accourait pour voir la fête organisée. »
\item « Le chef de district fit un discours interminable, chargé de promesses flatteuses. »
\item « Cette décoration éclatante ne remplacerait jamais tout ce que Meka avait sacrifié. »
\end{enumerate}
\textbf{Vérification de la règle :} dans chaque proposition, aucun groupe de plus de quatre mots consécutifs n'est repris de l'original ; le vocabulaire est renouvelé (\emph{médaille} $\rightarrow$ \emph{décoration} ; \emph{Blancs} $\rightarrow$ \emph{colons}) et la construction syntaxique est modifiée.
\end{corrige}

\begin{corrige}[Exercice 8 — Généralisation des énumérations]
\begin{enumerate}
\item « Les \textbf{arbres fruitiers} entouraient la case. » — Le terme générique englobe les trois espèces citées.
\item « Il vendit son \textbf{bétail} (ou : ses \textbf{animaux domestiques}) pour payer la fête. » — Chèvres, poules et moutons sont des animaux de basse-cour et d'élevage.
\item « Les \textbf{outils agricoles} étaient rangés derrière la porte. » — Houes, machettes et paniers servent aux travaux des champs.
\item « Les \textbf{autorités coloniales} (ou : les \textbf{notables}) prirent place à la tribune. » — Le terme rassemble les représentants de l'administration et des missions.
\item « Elle prépara des \textbf{plats traditionnels} pour les invités. » — Ndolé, miondo et poulet DG sont des mets camerounais.
\item « Les \textbf{moyens de transport} encombraient les rues de Douala. » — Taxis, motos et bus sont des véhicules de transport urbain.
\item « Il s'informa par tous les \textbf{médias}. » — Presse écrite, radio et télévision sont des médias.
\item « Les \textbf{fournitures scolaires} manquaient dans le cartable. » — Cahiers, stylos et règles sont du matériel d'écolier.
\item « Toute sa \textbf{famille} (ou : ses \textbf{proches}) vint le féliciter. » — Sœur, frère et cousins sont des membres de la famille.
\item « Les \textbf{intempéries} (ou : les \textbf{conditions climatiques}) épuisèrent les marcheurs. » — Pluie, vent et chaleur relèvent du climat.
\end{enumerate}
\end{corrige}

\begin{corrige}[Exercice 9 — Contraction de texte — type Baccalauréat technique]
\textbf{1. Thèse.} L'auteur soutient que l'éducation des jeunes filles est la condition première du développement du Cameroun : refuser l'école aux filles, c'est appauvrir la nation entière.

\textbf{2. Arguments et exemples.}
\begin{center}
\begin{tabularx}{\linewidth}{|X|X|}
\hline
\rowcolor{popblueL} \textbf{Argument} & \textbf{Exemple} \\
\hline
Une mère instruite élève des enfants plus sains et plus savants. & Les mères de l'Extrême-Nord alphabétisées qui vérifient les bulletins scolaires. \\
\hline
L'école protège la jeune fille du mariage précoce et de ses drames. & Les grossesses dangereuses et la dépendance totale des filles mariées tôt. \\
\hline
Le pays entier gagne à scolariser ses filles. & Les pays qui ont massivement scolarisé leurs filles ont transformé leur économie. \\
\hline
\end{tabularx}
\end{center}

\textbf{3. Résumé au quart (140 mots $\pm$ 10 \%).}

\emph{Prétendre qu'instruire une fille est un investissement perdu, puisqu'elle rejoindra un autre foyer, est un raisonnement ruineux pour la nation entière. Scolariser les filles est en effet indispensable au progrès du pays. Une mère éduquée assure à ses enfants une meilleure santé et de meilleurs résultats scolaires : former une fille, c'est former une génération. De plus, l'école préserve la jeune fille des mariages précoces et de leurs conséquences dramatiques, en lui offrant un avenir professionnel et la connaissance de ses droits. Enfin, priver la moitié de la population du savoir prive le pays de la moitié de ses talents, dont tous les secteurs ont besoin. Certes, la scolarité coûte cher et les traditions comptent, mais aucune d'elles ne saurait imposer l'ignorance. Il faut donc offrir à chaque fille les mêmes chances qu'à son frère.} (140 mots)

\textbf{4.} Nombre de mots : \textbf{140}.

\begin{attention}[Barème indicatif et points de vigilance]
Barème indicatif sur 20 : thèse (3 pts) ; arguments et exemples (6 pts) ; résumé (9 pts : respect du nombre de mots 2 pts, suppression des exemples 2 pts, généralisation 2 pts, reformulation sans reprise 2 pts, correction de la langue 1 pt) ; décompte exact (2 pts). Vigilance : aucune idée personnelle ajoutée ; aucun segment de plus de quatre mots recopié ; le résumé doit être rédigé en phrases liées, pas en simple liste de mots.
\end{attention}
\end{corrige}

\begin{corrige}[Exercice 10 — Contenu latent et ironie — \emph{Le Vieux Nègre et la médaille}]
\textbf{1.} Le \cle{contenu manifeste} est le sens explicite, littéral d'un discours : ce qui est dit. Le \cle{contenu latent} est le sens implicite, ce que le discours suggère ou cache derrière les mots.

\textbf{2.} Deux procédés ironiques :
\begin{itemize}
\item l'\textbf{éloge déplacé} : le commandant félicite Meka d'avoir « donné » ses terres et ses deux fils, transformant des pertes tragiques en mérites ;
\item le \textbf{contraste entre la solennité de la cérémonie et le néant de la récompense} : discours pompeux, fanfare et médaille brillante face à un vieil homme ruiné ; on peut ajouter l'antiphrase « morts pour la France », alors que ces fils sont morts pour une cause qui n'était pas la leur.
\end{itemize}

\textbf{3. Paragraphe de réponse.} Derrière l'éloge apparent, Oyono dénonce l'hypocrisie et la rapacité du système colonial. Le contenu manifeste du discours du commandant célèbre la fidélité et le sacrifice de Meka ; le contenu latent révèle une spoliation : les terres ont été prises au profit de la mission, les fils ont été sacrifiés dans une guerre étrangère. La médaille, simple bout de métal brillant, apparaît comme le prix dérisoire d'une vie dépossédée. L'ironie fait du discours officiel sa propre condamnation : plus le commandant s'étend sur les « bienfaits » de la colonisation, plus le lecteur mesure l'écart entre les belles paroles et la réalité misérable. La cérémonie d'honneur devient ainsi un réquisitoire.

\textbf{4. Justification de la démarche.} Oyono choisit l'ironie plutôt que la critique directe parce qu'elle est plus efficace et plus sûre : elle fait rire puis réfléchir, elle laisse le lecteur découvrir lui-même la vérité, ce qui la rend plus persuasive ; elle permettait aussi, à l'époque coloniale, de contourner la censure tout en dénonçant l'injustice avec une force qu'un pamphlet frontal n'aurait pas eue.
\end{corrige}

\begin{corrige}[Exercice 11 — Épreuve complète : réception et production]
\textbf{1. Compréhension.} L'extrait oppose l'éclat public de la cérémonie à la solitude douloureuse du soir. Champ lexical de la \textbf{gloire apparente} : \emph{médaille, brillait, applaudissaient, admirait}. Champ lexical de la \textbf{perte et de la solitude} : \emph{seul, ne reviendraient jamais, ne lui appartenaient plus}.

\textbf{2. Langue.}
\begin{itemize}
\item[(a)] Conjonction de coordination : \emph{et} (« quand la nuit tomba \textbf{et} que les invités furent partis ») : valeur d'\textbf{addition} ; on accepte aussi \emph{et} dans « la médaille à la main, \textbf{et} il pensa ».
\item[(b)] Conjonctions de subordination : \emph{quand} (« \textbf{quand} la nuit tomba ») : relation de \textbf{temps} ; \emph{que} (« et \textbf{que} les invités furent partis ») : subordonnée de \textbf{temps} reprise ; \emph{que} dans « ses deux fils \textbf{qui} ne reviendraient jamais » relève du pronom relatif (ne pas confondre). On accepte donc \emph{quand} et le \emph{que} repris après \emph{et}.
\end{itemize}

\textbf{3. Réécriture (proposition).} « Bien que la médaille brillât sur sa poitrine et que les Blancs applaudissent, Meka, qui avait attendu ce jour toute sa vie, sentit, quand le village eut fini de l'admirer, le poids de sa solitude. » — Toute phrase unique contenant une subordonnée de temps (\emph{quand}, \emph{lorsque}) et une subordonnée de concession (\emph{bien que} + subjonctif) est acceptée.

\textbf{4. Résumé (25 à 30 mots).} « Malgré l'honneur éclatant de la décoration et l'admiration de tous, Meka, redevenu solitaire, mesure l'étendue de ses pertes irréparables : enfants disparus et terres confisquées. » (27 mots)

\textbf{5. Expression personnelle (corrigé indicatif).} Une récompense ne peut compenser une perte irréparable, car ce qui est détruit à jamais — un enfant, une terre ancestrale — n'a pas de prix. La médaille de Meka en est la preuve : elle brille, mais elle ne rend ni ses fils ni ses champs ; elle n'est que le symbole dérisoire d'une gratitude hypocrite. Dans la vie courante, un dédommagement en argent versé à une famille après un accident mortel apaise peut-être les besoins matériels, mais il n'efface jamais le deuil. La vraie justice consisterait à prévenir la perte, non à la décorer après coup.

\begin{attention}[Barème indicatif et points de vigilance]
Barème sur 20 : compréhension (4 pts) ; langue (4 pts) ; réécriture (4 pts : phrase unique 1 pt, subordonnée de temps 1 pt, concession + subjonctif 1 pt, correction 1 pt) ; résumé (4 pts : nombre de mots 1 pt, reformulation 2 pts, correction 1 pt) ; expression personnelle (4 pts : prise de position claire, exemple pertinent, langue). Vigilance : ne pas confondre le \emph{que} conjonction et le \emph{qui}/\emph{que} pronom relatif ; subjonctif obligatoire après \emph{bien que} ; le résumé ne doit contenir aucune reprise de plus de quatre mots.
\end{attention}
\end{corrige}

\newpage

\coursec{Fiche bilan}

\begin{popfigure}
\begin{tikzpicture}[mindmap, grow cyclic, every node/.style={concept, text=white, font=\small},
  root concept/.append style={concept color=popblue, font=\bfseries},
  level 1/.append style={level distance=3.6cm, sibling angle=51.4},
  level 1 concept/.append style={font=\footnotesize}]
\node[root concept]{Formuler les idées essentielles d'un texte à résumer}
  child[concept color=poporange]{ node {Texte argumentatif : thèse, arguments, exemples} }
  child[concept color=popgreen]{ node {Contenus manifestes et latents} }
  child[concept color=poppurple]{ node {Lecture méthodique : \emph{Le Vieux Nègre et la médaille}} }
  child[concept color=popblue!80]{ node {Techniques de réduction : suppression, généralisation} }
  child[concept color=poppink]{ node {Outils de liaison : coordination, subordination} }
  child[concept color=popgold]{ node {Reformulation : mots propres, fidélité} }
  child[concept color=popmuted]{ node {Compte rendu et remédiation} };
\end{tikzpicture}
\legende{Carte mentale de la leçon : les sept savoirs articulés autour de la formulation des idées essentielles pour la contraction de texte.}
\end{popfigure}

\begin{bilan}
\textbf{Définitions clés}
\begin{itemize}
\item \cle{Texte argumentatif} : texte qui défend une \cle{thèse} (point de vue) à l'aide d'\cle{arguments} (raisons) illustrés par des \cle{exemples}.
\item \cle{Contenu manifeste} : ce que le texte dit explicitement. \cle{Contenu latent} : ce qu'il suggère (implicite, ironie, sous-entendus).
\item \cle{Idée essentielle} : idée sans laquelle le sens global du texte disparaît ; elle porte la thèse ou un argument majeur.
\item \cle{Contraction de texte} : réduction fidèle d'un texte (en général au quart, avec marge de 10 \%), sans commentaire ni copie de phrases entières.
\item \cle{Reformulation} : dire la même idée avec ses propres mots, sans trahir la pensée de l'auteur.
\end{itemize}

\textbf{Repères à connaître par cœur}
\begin{center}
\begin{tabularx}{\linewidth}{|l|X|}\hline
\rowcolor{popblueL} \textbf{Notion} & \textbf{À retenir} \\ \hline
Structure argumentative & Introduction (thèse) -- développement (arguments + exemples) -- conclusion. \\ \hline
Conjonctions de coordination & mais, ou, et, donc, or, ni, car : relient des éléments de même fonction. \\ \hline
Conjonctions de subordination & que, quand, comme, si, parce que, puisque, bien que, pour que : introduisent une subordonnée. \\ \hline
Techniques de réduction & suppression (détails, répétitions, exemples), généralisation (terme générique), condensation (fusion d'idées). \\ \hline
Compte de mots & Nombre de mots du résumé $\approx$ $\dfrac{\text{mots du texte}}{4}$, à $\pm$ 10 \%. \\ \hline
\end{tabularx}
\end{center}

\textbf{Méthodes express}
\begin{itemize}
\item Repérer les idées essentielles : lire le texte deux fois, souligner thèse et arguments, éliminer exemples et répétitions, dégager une idée par paragraphe.
\item Réduire : supprimer l'accessoire, généraliser les énumérations, condenser les développements, reformuler avec les connecteurs adaptés.
\item Vérifier : fidélité au sens, enchaînement logique, phrase d'ouverture présentant le texte, respect du nombre de mots.
\end{itemize}
\end{bilan}

\begin{aretenir}[Checklist avant le Bac]
\begin{itemize}
\item $\square$ Je sais définir un texte argumentatif et repérer thèse, arguments et exemples.
\item $\square$ Je distingue le contenu manifeste du contenu latent (implicite, ironie).
\item $\square$ Je connais les grandes étapes de la lecture méthodique du \emph{Vieux Nègre et la médaille}.
\item $\square$ Je maîtrise les trois techniques de réduction : suppression, généralisation, condensation.
\item $\square$ Je cite sans hésiter les sept conjonctions de coordination.
\item $\square$ Je reconnais les principales conjonctions de subordination et leur valeur logique.
\item $\square$ Je reformule une idée avec mes propres mots sans la déformer.
\item $\square$ Je calcule le nombre de mots attendu (quart du texte, à 10 \% près).
\item $\square$ Je rédige une phrase d'introduction présentant l'auteur, l'œuvre et la thèse.
\item $\square$ Je relie les idées du résumé par des connecteurs logiques variés.
\item $\square$ Je relis pour éliminer toute copie de phrase entière du texte.
\item $\square$ Je vérifie grammaire, orthographe et ponctuation avant de rendre ma copie.
\end{itemize}
\end{aretenir}

\findecours

\end{document}
